% Statistiques — Mathématiques Tle A — Cours signé OnBuch+
% Programme officiel MINESEC. Compiler avec : tectonic MA09-statistiques-a.tex
\newcommand{\DOCMATIERE}{Mathématiques}
\newcommand{\DOCNIVEAU}{Tle A}
\newcommand{\DOCMODULE}{Organisation et gestion des données}
\newcommand{\DOCLECON}{A09}
\newcommand{\DOCTITRE}{Statistiques}
% OnBuch+ — préambule « cours signés » ENRICHI (compilé avec tectonic / XeLaTeX)
% Système visuel « Pop » d'OnBuch+ : fond crème (jamais blanc), bordures encre
% épaisses, ombres « dures » décalées, titres Archivo Black en capitales.
% Version MATHÉMATIQUES : graphes (pgfplots/tikz), boîtes pédagogiques
% (définition, propriété, méthode, attention, le savais-tu, exercice résolu,
% exercice, TP, bilan), page de garde et signature OnBuch+ ancrée en bas.
%
% Macros attendues dans main.tex AVANT \input{preamble} :
%   \DOCMATIERE  \DOCNIVEAU  \DOCMODULE  \DOCLECON (numéro)  \DOCTITRE
\documentclass[11pt]{article}

\usepackage{fontspec}
\usepackage[french]{babel}
\usepackage[a4paper,margin=2cm,top=3.4cm,bottom=2.8cm,headheight=1.4cm,footskip=1.3cm]{geometry}
\usepackage{amsmath,amssymb,amsfonts}
\usepackage{mathtools} % \xmapsto, \coloneqq... souvent utilisés par les modèles
\usepackage{cancel} % simplifications barrées
% \abs{z} (module, valeur absolue) et \norm{u} (norme), utilisés par les modèles
\providecommand{\abs}{}\let\abs\relax\DeclarePairedDelimiter{\abs}{\lvert}{\rvert}
\providecommand{\norm}{}\let\norm\relax\DeclarePairedDelimiter{\norm}{\lVert}{\rVert}
\usepackage[table]{xcolor} % [table] : fournit aussi \rowcolors (alternance de lignes)
\usepackage{pagecolor}
\usepackage{tikz}
\usetikzlibrary{arrows.meta,positioning,calc,shapes.geometric,shapes.misc,shapes.arrows,decorations.pathmorphing,decorations.markings,patterns,shadows,mindmap,trees,fit,backgrounds,matrix,intersections}
\usepackage{pgfplots}
\usepackage{pgf-pie} % diagrammes circulaires \pie (statistiques)
\pgfplotsset{compat=1.18}
\usepgfplotslibrary{fillbetween,groupplots} % aires entre courbes (calcul intégral)
\usepackage{enumitem}
\usepackage{fancyhdr}
% [most] charge toutes les bibliothèques tcolorbox (listings, minted, raster,
% documentation...) et gonfle inutilement la mémoire TeX sur un document long
% (~300 boîtes) : on ne charge que ce qui est réellement utilisé.
\usepackage[skins,breakable]{tcolorbox}
\usepackage{graphicx}
\usepackage{colortbl}
\usepackage{booktabs}
\usepackage{tabularx}
\usepackage{diagbox} % case d'en-tête diagonale des tableaux à double entrée
% Colonnes extensibles centrée / gauche / droite (C, L, R), souvent utilisées par les modèles
\newcolumntype{C}{>{\centering\arraybackslash}X}
\newcolumntype{L}{>{\raggedright\arraybackslash}X}
\newcolumntype{R}{>{\raggedleft\arraybackslash}X}
\usepackage{array}
\usepackage{multicol}
\usepackage{siunitx}
% parse-numbers=false : accepte aussi \SI{2,5 \times 10^{-3}}{...}
\sisetup{locale=FR,output-decimal-marker={,},group-minimum-digits=4,parse-numbers=false}
\DeclareSIUnit\ohm{\ensuremath{\symup{\Omega}}} % U+2126 absent de la police
\usepackage{qrcode}
% \degree : commande fréquemment utilisée par les modèles pour un angle ou une
% température en dehors de \SI{}{} (non fournie par les paquets chargés).
\providecommand{\degree}{^{\circ}}
% \qed (fin de démonstration), utilisé par les modèles sans amsthm
\providecommand{\qed}{\hfill$\square$}
% \barre : double barre des valeurs interdites dans un tableau de variations (tkz-tab)
\providecommand{\barre}{\ensuremath{\|}}
% environnement proof (amsthm non chargé) utilisé par les modèles
\ifdefined\proof\else\newenvironment{proof}[1][Démonstration]{\par\noindent\textit{#1.}\ }{\qed\par}\fi
\usepackage{lastpage}
\usepackage{hyperref}
\hypersetup{hidelinks}

% ============================================================
% Palette « Pop » OnBuch+ — mêmes hex que la classe P (plus_ui.dart)
% ============================================================
\definecolor{poporange}{HTML}{F59321}
\definecolor{popdark}{HTML}{E07A0C}
\definecolor{popcream}{HTML}{FFF6E8}
\definecolor{popink}{HTML}{1C1714}
\definecolor{poppurple}{HTML}{7A5AE0}
\definecolor{popgreen}{HTML}{2E9E6A}
\definecolor{poppink}{HTML}{E0457B}
\definecolor{popblue}{HTML}{2F7DE1}
\definecolor{popgold}{HTML}{C98A12}
\definecolor{popmuted}{HTML}{8A7F76}
\definecolor{popline}{HTML}{EADFD2}
\definecolor{popgray}{HTML}{8A7F76}
\colorlet{popgrey}{popgray}
% Alias tolérants (le modèle nomme parfois les couleurs différemment)
\colorlet{popwhite}{white}
\colorlet{popblack}{popink}
\colorlet{popred}{poppink}
\colorlet{popyellow}{popgold}
% Teintes claires pour fonds de tableaux / zones
\colorlet{popblueL}{popblue!10!white}
\colorlet{poporangeL}{poporange!14!white}
\colorlet{popgreenL}{popgreen!12!white}
\colorlet{poppurpleL}{poppurple!12!white}
\colorlet{poppinkL}{poppink!12!white}
\colorlet{popgoldL}{popgold!14!white}

\pagecolor{popcream}

% ============================================================
% Typographie
% ============================================================
\defaultfontfeatures{Ligatures=TeX}
\setmainfont{PlusJakartaSans-Medium.ttf}[Path=../../../fonts/,BoldFont=PlusJakartaSans-Bold.ttf]
% Maths en sans-serif (Fira Math) avec chiffres et lettres droites de Plus
% Jakarta, pour que les formules se fondent dans le texte.
\usepackage[math-style=ISO,bold-style=ISO]{unicode-math}
\setmathfont{FiraMath-Regular.otf}[Path=../../../fonts/]
\setmathfont{PlusJakartaSans-Medium.ttf}[Path=../../../fonts/,range={up/num,up/{Latin,latin},\mathup}]
\usepackage{icomma} % virgule décimale sans espace en mode math
% Caractères mathématiques Unicode absents des polices de texte (Plus Jakarta,
% Archivo Black) : rendus par leur équivalent mathématique, en texte comme en
% formule — sinon ils s'affichent en carré vide (ex. « Arithmétique dans ℤ »).
\usepackage{newunicodechar}
\newunicodechar{ℝ}{\ensuremath{\mathbb{R}}}
\newunicodechar{ℤ}{\ensuremath{\mathbb{Z}}}
\newunicodechar{ℂ}{\ensuremath{\mathbb{C}}}
\newunicodechar{ℕ}{\ensuremath{\mathbb{N}}}
\newunicodechar{ℚ}{\ensuremath{\mathbb{Q}}}
\newunicodechar{𝔻}{\ensuremath{\mathbb{D}}}
\newunicodechar{α}{\ensuremath{\alpha}}
\newunicodechar{β}{\ensuremath{\beta}}
\newunicodechar{γ}{\ensuremath{\gamma}}
\newunicodechar{δ}{\ensuremath{\delta}}
\newunicodechar{ε}{\ensuremath{\varepsilon}}
\newunicodechar{θ}{\ensuremath{\theta}}
\newunicodechar{λ}{\ensuremath{\lambda}}
\newunicodechar{μ}{\ensuremath{\mu}}
\newunicodechar{σ}{\ensuremath{\sigma}}
\newunicodechar{τ}{\ensuremath{\tau}}
\newunicodechar{φ}{\ensuremath{\varphi}}
\newunicodechar{ψ}{\ensuremath{\psi}}
\newunicodechar{ω}{\ensuremath{\omega}}
\newunicodechar{Σ}{\ensuremath{\Sigma}}
% majuscules produites par \MakeUppercase dans les titres (α-aminés -> Α)
\newunicodechar{Α}{\ensuremath{\alpha}}
\newunicodechar{Β}{\ensuremath{\beta}}
\newunicodechar{Γ}{\ensuremath{\Gamma}}
\newunicodechar{Δ}{\ensuremath{\Delta}}
\newunicodechar{Ε}{\ensuremath{\varepsilon}}
\newunicodechar{Θ}{\ensuremath{\Theta}}
\newunicodechar{Λ}{\ensuremath{\Lambda}}
\newunicodechar{Μ}{\ensuremath{\mu}}
\newunicodechar{Τ}{\ensuremath{\tau}}
\newunicodechar{Φ}{\ensuremath{\Phi}}
\newunicodechar{Ψ}{\ensuremath{\Psi}}
\newunicodechar{Ω}{\ensuremath{\Omega}}
\newunicodechar{↦}{\ensuremath{\mapsto}}
\newunicodechar{∈}{\ensuremath{\in}}
\newunicodechar{∉}{\ensuremath{\notin}}
\newunicodechar{∧}{\ensuremath{\wedge}}
\newunicodechar{⇔}{\ensuremath{\Leftrightarrow}}
\newunicodechar{⟺}{\ensuremath{\Longleftrightarrow}}
\newunicodechar{⇒}{\ensuremath{\Rightarrow}}
\newunicodechar{⟹}{\ensuremath{\Longrightarrow}}
\newunicodechar{∘}{\ensuremath{\circ}}
\newunicodechar{∪}{\ensuremath{\cup}}
\newunicodechar{∩}{\ensuremath{\cap}}
\newunicodechar{≡}{\ensuremath{\equiv}}
\newunicodechar{⊂}{\ensuremath{\subset}}
\newunicodechar{⊥}{\ensuremath{\perp}}
\newunicodechar{∃}{\ensuremath{\exists}}
\newunicodechar{∀}{\ensuremath{\forall}}
\newunicodechar{∛}{\ensuremath{\sqrt[3]{\,}}}
\newunicodechar{∜}{\ensuremath{\sqrt[4]{\,}}}
\newunicodechar{⃗}{\ensuremath{{}^{\to}}}
\newunicodechar{ⁿ}{\ifmmode^{n}\else\textsuperscript{\ensuremath{n}}\fi}
\newunicodechar{ˣ}{\ifmmode^{x}\else\textsuperscript{\ensuremath{x}}\fi}
\newunicodechar{ᵇ}{\ifmmode^{b}\else\textsuperscript{\ensuremath{b}}\fi}
\newunicodechar{ʳ}{\ifmmode^{r}\else\textsuperscript{\ensuremath{r}}\fi}
\newunicodechar{ᵘ}{\ifmmode^{u}\else\textsuperscript{\ensuremath{u}}\fi}
\newunicodechar{ʸ}{\ifmmode^{y}\else\textsuperscript{\ensuremath{y}}\fi}
\newunicodechar{ᵃ}{\ifmmode^{a}\else\textsuperscript{\ensuremath{a}}\fi}
\newunicodechar{ᵉ}{\ifmmode^{e}\else\textsuperscript{\ensuremath{e}}\fi}
\newunicodechar{ᵏ}{\ifmmode^{k}\else\textsuperscript{\ensuremath{k}}\fi}
\newunicodechar{ᵖ}{\ifmmode^{p}\else\textsuperscript{\ensuremath{p}}\fi}
\newunicodechar{ᴺ}{\ifmmode^{N}\else\textsuperscript{\ensuremath{N}}\fi}
\newunicodechar{ᶜ}{\ifmmode^{c}\else\textsuperscript{\ensuremath{c}}\fi}
\newunicodechar{ʰ}{\ifmmode^{h}\else\textsuperscript{\ensuremath{h}}\fi}
\newunicodechar{ᵠ}{\ifmmode^{q}\else\textsuperscript{\ensuremath{q}}\fi}
\newunicodechar{ₙ}{\ifmmode_{n}\else\textsubscript{\ensuremath{n}}\fi}
\newunicodechar{ₐ}{\ifmmode_{a}\else\textsubscript{\ensuremath{a}}\fi}
\newunicodechar{ᵢ}{\ifmmode_{i}\else\textsubscript{\ensuremath{i}}\fi}
\newunicodechar{ᵣ}{\ifmmode_{r}\else\textsubscript{\ensuremath{r}}\fi}
\newunicodechar{ₖ}{\ifmmode_{k}\else\textsubscript{\ensuremath{k}}\fi}
\newunicodechar{ᵦ}{\ifmmode_{\beta}\else\textsubscript{\ensuremath{\beta}}\fi}
\newunicodechar{ₓ}{\ifmmode_{x}\else\textsubscript{\ensuremath{x}}\fi}
\newfontfamily\popdisplay{ArchivoBlack-Regular.ttf}[Path=../../../fonts/]
\newfontfamily\poplogo{SpaceGrotesk-Bold.ttf}[Path=../../../fonts/]
\newfontfamily\popsemi{PlusJakartaSans-SemiBold.ttf}[Path=../../../fonts/]
\newfontfamily\popxbold{PlusJakartaSans-ExtraBold.ttf}[Path=../../../fonts/]
\newfontfamily\popxboldls{PlusJakartaSans-ExtraBold.ttf}[Path=../../../fonts/,LetterSpace=3.0]

\setlist[enumerate]{leftmargin=1.7em,itemsep=5pt,topsep=4pt}
\setlist[itemize]{leftmargin=1.7em,itemsep=4pt,topsep=4pt,label={\color{poporange}\textbullet}}
\setlength{\parindent}{0pt}
\setlength{\parskip}{5pt}
\renewcommand{\arraystretch}{1.25}

% pgfplots : style Pop par défaut
\pgfplotsset{
  every axis/.append style={axis line style={popink,line width=1pt},tick style={popink},
    grid style={popline,line width=0.5pt},label style={font=\small},tick label style={font=\footnotesize}},
  popcurve/.style={poporange,line width=1.8pt,no markers,smooth},
}
% Même style disponible dans un \draw TikZ simple (hors pgfplots)
\tikzset{popcurve/.style={poporange,line width=1.8pt}}

% ============================================================
% Pill / squiggle
% ============================================================
\newcommand{\poppill}[3]{%
  \tikz[baseline=-0.55ex]\node[draw=popink,line width=1.2pt,rounded corners=2.6pt,%
    fill=#2,inner xsep=6.5pt,inner ysep=3pt]{%
    \popxboldls\fontsize{7}{7}\selectfont\color{#3}\MakeUppercase{#1}};%
}
\newcommand{\popsquiggle}[1]{%
  \tikz[baseline=-0.4ex]{
    \draw[#1,line width=2.6pt,line cap=round]
      plot [smooth] coordinates {(0,0.09) (0.34,0.01) (0.68,0.21) (1.02,0.01) (1.36,0.10)};
  }%
}
% Mot-clé surligné (marqueur orange sous le texte)
\newcommand{\cle}[1]{{\popxbold\color{popdark}#1}}

% ============================================================
% En-tête / pied
% ============================================================
\pagestyle{fancy}
\fancyhf{}
\renewcommand{\headrulewidth}{0pt}
\renewcommand{\footrulewidth}{0pt}
\lhead{\raisebox{-0.15ex}{\poplogo\fontsize{13}{13}\selectfont\color{popink}OnBuch\color{poporange}+}}
\rhead{\poppill{\DOCMATIERE\ \textbullet\ \DOCNIVEAU\ \textbullet\ Leçon \DOCLECON}{white}{popink}}
\lfoot{\popxbold\fontsize{7.6}{7.6}\selectfont\color{popmuted}\MakeUppercase{Cours signé OnBuch+}\ \textbullet\ {\popsemi\DOCTITRE}}
\rfoot{\tikz[baseline=-0.6ex]\node[rounded corners=8.5pt,fill=popink,minimum height=17pt,minimum width=17pt,inner xsep=5pt,inner sep=0]{\popxbold\fontsize{8}{8}\selectfont\color{popcream}\thepage\,/\,\pageref*{LastPage}};}

% ============================================================
% Titres
% ============================================================
\newcommand{\chaptitre}[2]{%
  \begin{tcolorbox}[enhanced,colback=poporange,colframe=popink,boxrule=2.6pt,arc=11pt,
      drop shadow=popink,left=15pt,right=15pt,top=12pt,bottom=11pt,before skip=3pt,after skip=18pt]
    \poppill{#2}{popink}{popcream}\par\vspace{9pt}
    {\raggedright\hyphenpenalty=10000\exhyphenpenalty=10000\popdisplay\fontsize{22}{24}\selectfont\color{popcream}\MakeUppercase{#1}\par}
  \end{tcolorbox}
}
% Section numérotée : pastille orange avec le numéro + titre Archivo
\newcounter{popsec}
\newcommand{\coursec}[1]{%
  \stepcounter{popsec}\setcounter{popsub}{0}%
  \par\vspace{16pt}\noindent
  \begin{minipage}{\linewidth}
  \tikz[baseline=-0.7ex]\node[draw=popink,line width=1.6pt,fill=poporange,rounded corners=4pt,minimum size=22pt,inner sep=0]{\popdisplay\fontsize{12}{12}\selectfont\color{popcream}\thepopsec};%
  \hspace{8pt}{\popdisplay\fontsize{15}{17}\selectfont\color{popink}\MakeUppercase{#1}}\par
  \vspace{4pt}\noindent{\color{poporange}\rule{36pt}{3.6pt}}
  \end{minipage}\par\nopagebreak\vspace{8pt}
}
\newcounter{popsub}
\newcommand{\courssub}[1]{%
  \stepcounter{popsub}%
  \par\vspace{9pt}\noindent{\popxbold\fontsize{11.5}{13}\selectfont\color{popink}{\color{poporange}\thepopsec.\thepopsub}\hspace{6pt}#1}\par\nopagebreak\vspace{4pt}
}

% ============================================================
% Boîtes pédagogiques — même recette : fond blanc, bordure épaisse colorée,
% ombre dure encre, badge plein. Titre optionnel [#1].
% ============================================================
\tcbset{popbase/.style={breakable,enhanced,colback=white,boxrule=2.3pt,arc=9pt,
  shadow={3pt}{-3pt}{0pt}{popink},left=14pt,right=14pt,top=11pt,bottom=11pt,before skip=11pt,after skip=15pt,
  fontupper=\popsemi\fontsize{10.6}{14.5}\selectfont\color{popink},
  fontlower=\popsemi\fontsize{10.6}{14.5}\selectfont\color{popink}}}
% Coche dessinée (le glyphe ✓ n'existe pas dans les polices du document)
\newcommand{\popcheck}[1]{\tikz[baseline=-0.5ex]\draw[#1,line width=1.6pt,line cap=round,line join=round](-0.13,0.0)--(-0.04,-0.09)--(0.13,0.1);}
\newcommand{\popboxhead}[3]{\poppill{#1}{#2}{white}\ifx\relax#3\relax\else\hspace{6pt}{\popxbold\fontsize{10.6}{12}\selectfont\color{popink}#3}\fi\par\vspace{7pt}}

\newtcolorbox{aretenir}[1][]{popbase,colframe=popgreen,before upper={\popboxhead{À retenir}{popgreen}{#1}}}
\newtcolorbox{exemplebox}[1][]{popbase,colframe=poporange,before upper={\popboxhead{Exemple}{poporange}{#1}}}
\newtcolorbox{definition}[1][]{popbase,colframe=popblue,before upper={\popboxhead{Définition}{popblue}{#1}}}
\newtcolorbox{propriete}[1][]{popbase,colframe=poppurple,before upper={\popboxhead{Propriété}{poppurple}{#1}}}
\newtcolorbox{methode}[1][]{popbase,colframe=popgold,before upper={\popboxhead{Méthode}{popgold}{#1}}}
\newtcolorbox{attention}[1][]{popbase,colframe=poppink,before upper={\popboxhead{Attention}{poppink}{#1}}}
\newtcolorbox{experience}[1][]{popbase,colframe=popblue,colback=popblueL,before upper={\popboxhead{Expérience}{popblue}{#1}}}
% « Le savais-tu ? » : carte violette pleine (contexte, culture, Cameroun)
\newtcolorbox{savaistu}[1][]{popbase,colback=poppurple,colframe=popink,
  fontupper=\popsemi\fontsize{10.4}{14.2}\selectfont\color{white},
  before upper={\poppill{Le savais-tu\,?}{white}{poppurple}\ifx\relax#1\relax\else\hspace{6pt}{\popxbold\color{white}#1}\fi\par\vspace{7pt}}}
% Exercice résolu : énoncé en haut, \tcblower puis la solution sur fond crème
\newcounter{popexo}\newcounter{popexr}
\newtcolorbox{exoresolu}[1][]{popbase,colframe=poporange,colbacklower=popcream,
  segmentation style={popink,line width=1.2pt,dashed},
  before upper={\stepcounter{popexr}\popboxhead{Exercice résolu \thepopexr}{poporange}{#1}},
  before lower={\poppill{Solution}{popink}{popcream}\par\vspace{6pt}}}
% Exercice (non résolu en place) — la correction est dans la partie Corrigés
\newtcolorbox{exercice}[1][]{popbase,colframe=popink,
  before upper={\stepcounter{popexo}\popboxhead{Exercice \thepopexo}{popink}{#1}}}
% Corrigé
\newtcolorbox{corrige}[1][]{popbase,colframe=popgreen,colback=popgreenL,
  before upper={\popboxhead{Corrigé}{popgreen}{#1}}}
% Objectifs / prérequis / bilan
\newtcolorbox{objectifs}{popbase,colframe=popink,colback=poporangeL,
  before upper={\popboxhead{Objectifs de la leçon}{popink}{}}}
\newtcolorbox{prerequis}{popbase,colframe=popink,colback=popblueL,
  before upper={\popboxhead{Prérequis}{popblue}{}}}
\newtcolorbox{bilan}{popbase,colframe=popink,colback=white,boxrule=2.8pt,
  before upper={\popboxhead{Fiche bilan}{poporange}{L'essentiel à connaître pour l'examen}}}
% Figure encadrée avec légende
\newtcolorbox{popfigure}[1][]{enhanced,breakable=false,colback=white,colframe=popline,boxrule=1.4pt,arc=8pt,
  left=10pt,right=10pt,top=10pt,bottom=8pt,before skip=10pt,after skip=12pt,halign=center,
  lower separated=false,halign lower=center,
  fontlower=\popsemi\fontsize{9}{11.5}\selectfont\color{popmuted},
  #1}
\newcommand{\legende}[1]{\par\vspace{6pt}{\popsemi\fontsize{9}{11.5}\selectfont\color{popmuted}#1}}

% ============================================================
% Signature OnBuch+
% ============================================================
\newcommand{\poplogoinline}[1]{{\poplogo\fontsize{#1}{#1}\selectfont\color{popink}OnBuch\color{poporange}+}}
\newcommand{\signatureonbuch}{%
  \begin{tcolorbox}[enhanced,colback=popink,colframe=popink,boxrule=2.2pt,arc=10pt,
      drop shadow=poporange,left=14pt,right=14pt,top=10pt,bottom=10pt,before skip=0pt,after skip=0pt]
    \tikz[baseline=-0.7ex]\node[circle,fill=popgreen,draw=popcream,line width=1.3pt,minimum size=22pt,inner sep=0]{\popcheck{white}};%
    \hspace{9pt}\begin{minipage}[c]{0.62\linewidth}
      {\popxbold\fontsize{12}{13}\selectfont\color{popcream}Signé\ \textbullet\ {\poplogo OnBuch\color{poporange}+}}\par\vspace{2pt}
      {\raggedright\popsemi\fontsize{8}{10.5}\selectfont\color{popline}\MakeUppercase{Équipe pédagogique OnBuch+}\\\MakeUppercase{Conforme au programme officiel MINESEC}\par}
    \end{minipage}\hfill
    {\poplogo\fontsize{22}{22}\selectfont\color{popcream}OnBuch\color{poporange}+}
  \end{tcolorbox}%
}

% ============================================================
% Page de garde
% #1 titre · #2 matière · #3 niveau · #4 module · #5 n° leçon · #6 durée ·
% #7 description / objectifs (texte ou itemize)
% ============================================================
\newcommand{\pagegarde}[7]{%
  \thispagestyle{empty}
  \newgeometry{a4paper,margin=1.7cm,top=1.6cm,bottom=1.4cm}%
  \begin{tikzpicture}[remember picture,overlay]
    \fill[poporange!18] ([xshift=-3.2cm,yshift=-3.6cm]current page.north east) circle (4.6cm);
    \fill[poppurple!14] ([xshift=2.4cm,yshift=7.6cm]current page.south west) circle (3.8cm);
  \end{tikzpicture}%
  {\poplogo\fontsize{30}{30}\selectfont\color{popink}OnBuch\color{poporange}+}\hfill
  \raisebox{0.6ex}{\poppill{Cours signé \textbullet\ Premium}{popink}{popcream}}\par
  \vspace{1pt}\popsquiggle{poppurple}\par
  \vspace{8pt}
  \begin{tcolorbox}[enhanced,colback=poporange,colframe=popink,boxrule=2.8pt,arc=13pt,
      drop shadow=popink,left=18pt,right=18pt,top=12pt,bottom=13pt,after skip=10pt]
    \poppill{#2 \textbullet\ #3}{popink}{popcream}\hspace{5pt}\poppill{#4}{white}{popink}\par\vspace{8pt}
    {\popdisplay\fontsize{14}{14}\selectfont\color{popink}LEÇON #5}\par\vspace{4pt}
    {\raggedright\hyphenpenalty=10000\exhyphenpenalty=10000\popdisplay\fontsize{25}{27}\selectfont\color{popcream}\MakeUppercase{#1}\par}
    \vspace{8pt}
    {\popxbold\fontsize{9.5}{9.5}\selectfont\color{popink}\MakeUppercase{Durée conseillée : #6}}
  \end{tcolorbox}
  \begin{tcolorbox}[enhanced,colback=white,colframe=popink,boxrule=2.2pt,arc=10pt,
      drop shadow=popink,left=16pt,right=16pt,top=10pt,bottom=10pt,after skip=8pt]
    \poppill{Dans cette leçon}{poppurple}{white}\par\vspace{5pt}
    \popsemi\fontsize{9.3}{12.6}\selectfont\color{popink}#7
  \end{tcolorbox}
  \noindent\begin{minipage}[t]{0.487\linewidth}
    \begin{tcolorbox}[enhanced,colback=popgreen,colframe=popink,boxrule=2.2pt,arc=10pt,
        drop shadow=popink,left=11pt,right=11pt,top=10pt,bottom=10pt,height=3.1cm,valign=center]
      \tcbox[colback=white,colframe=popink,boxrule=1.3pt,arc=5pt,boxsep=0pt,nobeforeafter,
        left=3pt,right=3pt,top=3pt,bottom=3pt,valign=center]{\qrcode[height=1.9cm]{https://play.google.com/store/apps/details?id=com.luvvix.onbuch&hl=fr}}%
      \hspace{8pt}\begin{minipage}[c]{0.5\linewidth}
        {\popdisplay\fontsize{11}{12.5}\selectfont\color{white}\MakeUppercase{Télécharge OnBuch}}\par\vspace{4pt}
        {\raggedright\popsemi\fontsize{8.4}{11}\selectfont\color{white}Gratuit \textbullet\ Google Play\\Tuteur IA, annales, concours, résultats\par}
      \end{minipage}
    \end{tcolorbox}
  \end{minipage}\hfill
  \begin{minipage}[t]{0.487\linewidth}
    \begin{tcolorbox}[enhanced,colback=poppurple,colframe=popink,boxrule=2.2pt,arc=10pt,
        drop shadow=popink,left=11pt,right=11pt,top=10pt,bottom=10pt,height=3.1cm,valign=center]
      \tikz[baseline=-0.7ex]\node[circle,fill=white,draw=popink,line width=1.3pt,minimum size=1.6cm,inner sep=0]{\popdisplay\fontsize{24}{24}\selectfont\color{poppurple}+};%
      \hspace{8pt}\begin{minipage}[c]{0.6\linewidth}
        {\popdisplay\fontsize{11}{12.5}\selectfont\color{white}\MakeUppercase{Passe à OnBuch+}}\par\vspace{4pt}
        {\raggedright\popsemi\fontsize{8.4}{11}\selectfont\color{white}Tous les cours signés, Tuteur IA illimité, fiches PDF — onglet OnBuch+ de l'app\par}
      \end{minipage}
    \end{tcolorbox}
  \end{minipage}
  \vspace{8pt}
  \popcontenu
  \vfill
  \signatureonbuch
  \restoregeometry
  \newpage
}

% Rangée « Ce cours contient » (page de garde)
\newcommand{\poptile}[3]{% #1 couleur #2 gros texte #3 légende
  \begin{tcolorbox}[enhanced,colback=#1,colframe=popink,boxrule=2pt,arc=9pt,shadow={2.5pt}{-2.5pt}{0pt}{popink},
      left=6pt,right=6pt,top=9pt,bottom=9pt,halign=center,valign=center,height=1.75cm,width=\linewidth,nobeforeafter]
    {\popdisplay\fontsize{12.5}{13}\selectfont\color{white}\MakeUppercase{#2}}\par\vspace{3pt}
    {\centering\popsemi\fontsize{7.8}{9.5}\selectfont\color{white}#3\par}
  \end{tcolorbox}}
\newcommand{\popcontenu}{%
  \par\noindent{\popxboldls\fontsize{7.5}{7.5}\selectfont\color{popmuted}CE COURS SIGNÉ CONTIENT}\par\vspace{7pt}
  \noindent\begin{minipage}[t]{0.235\linewidth}\poptile{popblue}{Cours}{complet, illustré, pas à pas}\end{minipage}\hfill
  \begin{minipage}[t]{0.235\linewidth}\poptile{poporange}{Méthodes}{et exercices résolus}\end{minipage}\hfill
  \begin{minipage}[t]{0.235\linewidth}\poptile{poppink}{Exercices}{type examen + corrigés}\end{minipage}\hfill
  \begin{minipage}[t]{0.235\linewidth}\poptile{popgreen}{Fiche bilan}{l'essentiel à retenir}\end{minipage}\par}

% Sommaire
\newtcolorbox{sommaire}{enhanced,colback=white,colframe=popink,boxrule=2.2pt,arc=10pt,
  drop shadow=popink,left=18pt,right=18pt,top=13pt,bottom=13pt,before skip=6pt,after skip=20pt,
  before upper={\poppill{Sommaire}{poppurple}{white}\par\vspace{9pt}\popsemi\fontsize{10.6}{14.5}\selectfont\color{popink}}}

% Signature de fin de cours — poussée tout en bas de la dernière page
\newcommand{\findecours}{%
  \par\vfill\nopagebreak
  \begin{center}\popsquiggle{poporange}\end{center}\vspace{4pt}
  \signatureonbuch
}
\AtBeginDocument{\def\llbracket{\mathopen{[\mkern-3.5mu[}}\def\rrbracket{\mathclose{]\mkern-3.5mu]}}} % intervalles d'entiers (glyphe ⟦ absent de la police)
% Glyphes absents de Fira Math : redéfinis à partir de glyphes disponibles
\AtBeginDocument{%
  \def\setminus{\mathbin{\backslash}}\let\smallsetminus\setminus
  \def\vdots{\mathord{\vbox{\baselineskip3.2pt\lineskiplimit0pt\kern4pt\hbox{.}\hbox{.}\hbox{.}}}}%
  \def\ddots{\mathinner{\mkern1mu\raise6.5pt\hbox{.}\mkern2mu\raise3.6pt\hbox{.}\mkern2mu\raise0.7pt\hbox{.}\mkern1mu}}%
  \def\triangle{\mathord{\scalebox{0.9}{$\symup{\Delta}$}}}%
  \def\nabla{\mathord{\raisebox{\depth}{\scalebox{1}[-1]{$\symup{\Delta}$}}}}%
  \def\checkmark{\popcheck{popdark}}%
  \def\star{\mathbin{\ast}}\def\bigstar{\mathord{\ast}}%
  \def\diamond{\mathbin{\scalebox{0.6}{$\Diamond$}}}%
}

%% FIGFIX-BEGIN v3 — correctifs globaux des figures (docs/onbuch-generation/tools/figfix)
\usepackage{adjustbox}\usepackage{etoolbox}
% toute figure TikZ/pgfplots plus large que la ligne est réduite à la largeur disponible
\BeforeBeginEnvironment{tikzpicture}{\begin{adjustbox}{max width=\linewidth}}
\AfterEndEnvironment{tikzpicture}{\end{adjustbox}}
% légendes pgfplots lisibles par-dessus les courbes
\pgfplotsset{every axis/.append style={legend style={fill=white,fill opacity=0.92,text opacity=1,draw=black!15,inner sep=3pt}}}
% étiquettes d'arêtes (syntaxe edge["..."]) avec fond blanc
\tikzset{every edge quotes/.append style={fill=white,inner sep=1.2pt}}
% bulles de cartes mentales : texte à la ligne dans la bulle, bulle un peu plus grande
\tikzset{every concept/.append style={text width=2.0cm,align=center,minimum size=2.3cm,inner sep=2pt}}
%% FIGFIX-END
\begin{document}

\pagegarde{Statistiques}{Mathématiques}{Tle A}{Organisation et gestion des données}{A09}{9 h}{Cette leçon présente les outils statistiques à une variable (diagrammes, paramètres de position et de dispersion) puis à deux variables (nuage de points, point moyen, ajustement affine par la méthode de Mayer). Elle aboutit à l'exploitation de données réelles pour décrire, résumer et prévoir, compétences exigées au Baccalauréat A.
\vspace{4pt}
\begin{multicols}{2}\small
\begin{itemize}
\item À la fin de cette leçon, je sais organiser une série statistique en tableau (effectifs, fréquences, effectifs cumulés)
\item À la fin de cette leçon, je sais construire et lire un diagramme circulaire, à bandes, à bâtons et un histogramme
\item À la fin de cette leçon, je sais construire et exploiter le polygone des fréquences cumulées
\item À la fin de cette leçon, je sais calculer le mode, la classe modale, la moyenne et la médiane d'une série
\item À la fin de cette leçon, je sais déterminer les quartiles et les déciles d'une série (sans interpolation linéaire)
\item À la fin de cette leçon, je sais construire un nuage de points et placer le point moyen
\end{itemize}\end{multicols}}

\begin{sommaire}
\begin{multicols}{2}
\begin{enumerate}
\item Vocabulaire et organisation des données
\item Représentations graphiques
\item Paramètres de position : mode, moyenne, médiane
\item Quartiles et déciles
\item Séries à deux caractères : nuage de points et point moyen
\item Ajustement affine par la méthode de Mayer et prévisions
\item Activité d'intégration
\item Méthodes et exercices résolus
\item Exercices
\item Corrigés des exercices
\item Fiche bilan
\end{enumerate}
\end{multicols}
\end{sommaire}

\begin{objectifs}
\begin{itemize}
\item À la fin de cette leçon, je sais organiser une série statistique en tableau (effectifs, fréquences, effectifs cumulés)
\item À la fin de cette leçon, je sais construire et lire un diagramme circulaire, à bandes, à bâtons et un histogramme
\item À la fin de cette leçon, je sais construire et exploiter le polygone des fréquences cumulées
\item À la fin de cette leçon, je sais calculer le mode, la classe modale, la moyenne et la médiane d'une série
\item À la fin de cette leçon, je sais déterminer les quartiles et les déciles d'une série (sans interpolation linéaire)
\item À la fin de cette leçon, je sais construire un nuage de points et placer le point moyen
\item À la fin de cette leçon, je sais déterminer une droite d'ajustement par la méthode de Mayer
\item À la fin de cette leçon, je sais utiliser une droite d'ajustement pour faire une prévision et en discuter la fiabilité
\end{itemize}
\end{objectifs}

\begin{prerequis}
\begin{itemize}
\item Calculs sur les fractions, pourcentages et proportionnalité (Seconde/Première)
\item Lecture et construction de tableaux de données
\item Coordonnées d'un point dans un repère et équation d'une droite y = ax + b (Première)
\item Notion de moyenne simple vue au collège
\item Représentation graphique d'une fonction affine
\end{itemize}
\end{prerequis}

\newpage

\coursec{Vocabulaire et organisation des données}

À la coopérative de quartier de Mballa II, à Yaoundé, l'ambiance est studieuse ce samedi matin. La présidente, Mme Etoundi, tient entre les mains un cahier rempli de chiffres : pendant douze mois, les membres de la coopérative ont relevé les quantités de sacs de manioc vendues et les prix pratiqués par les commerçants du marché. Demain, à l'assemblée générale, il faudra fixer le prix de vente des sacs pour la saison qui vient. Mais comment présenter clairement cette montagne de chiffres ? Quel prix « central » la coopérative doit-elle choisir ? Et surtout — question qui brûle toutes les lèvres — si le prix augmente, peut-on \emph{prévoir} la quantité qui sera vendue le mois prochain ?

Toute cette leçon va donner à la coopérative — et à toi — les outils mathématiques pour répondre à ces questions. Mais avant de calculer quoi que ce soit, il faut d'abord apprendre à \cle{organiser} les données : c'est l'objet de cette première section. Car en statistiques, comme en cuisine, un bon résultat commence par de bons ingrédients bien rangés.

\courssub{Population, individu et caractère}

Imaginons que Mme Etoundi veuille étudier les dépenses journalières des ménages du quartier Mballa II. Elle ne peut pas interroger tous les habitants de Yaoundé : elle choisit un ensemble bien précis de personnes à observer, et pour chacune, elle note une information précise (la dépense du jour). Cette démarche repose sur trois mots de vocabulaire fondamentaux.

\begin{definition}[Population, individu, caractère]
\begin{itemize}
\item La \cle{population} est l'ensemble sur lequel porte l'étude statistique.
\item Un \cle{individu} (ou unité statistique) est un élément de cette population.
\item Un \cle{caractère} (ou variable statistique) est une propriété que l'on observe ou mesure sur chaque individu.
\end{itemize}
\end{definition}

\begin{exemplebox}[Identifier population, individus et caractère]
\begin{itemize}
\item Étude des notes au baccalauréat des élèves du lycée de Nkol-Eton : la population est l'ensemble des candidats du lycée, un individu est un candidat, le caractère est la note obtenue.
\item Étude du nombre de sacs de manioc vendus chaque mois par la coopérative : la population est l'ensemble des douze mois de l'année, un individu est un mois, le caractère est le nombre de sacs vendus.
\item Enquête sur l'opérateur téléphonique préféré (MTN, Orange, Nexttel) des habitants d'un quartier : la population est l'ensemble des habitants, un individu est un habitant, le caractère est l'opérateur préféré.
\end{itemize}
\end{exemplebox}

Tous les caractères ne se ressemblent pas. Répondre « MTN » à une enquête et répondre « \SI{2500}{F} » n'appellent pas les mêmes traitements mathématiques. On distingue donc plusieurs natures de caractères.

\begin{definition}[Nature d'un caractère]
Un caractère est dit :
\begin{itemize}
\item \cle{qualitatif} lorsqu'il n'est pas mesurable par un nombre : il décrit une catégorie, une modalité (couleur, sexe, opérateur téléphonique, mention au bac) ;
\item \cle{quantitatif} lorsqu'il s'exprime par un nombre. On distingue alors :
\begin{itemize}
\item le caractère \cle{quantitatif discret}, qui ne prend que des valeurs isolées, souvent entières (nombre d'enfants par ménage, nombre de sacs vendus, note entière sur 20) ;
\item le caractère \cle{quantitatif continu}, qui peut prendre n'importe quelle valeur dans un intervalle (taille, masse, température, dépense en francs, durée d'un appel).
\end{itemize}
\end{itemize}
\end{definition}

\begin{attention}[Discret ou continu ?]
Un caractère peut être \emph{théoriquement} continu mais relevé de façon discrète. La taille d'un élève est continue, mais si on la mesure au centimètre près, on obtient des valeurs comme 165, 166, 167\ldots{} C'est la \emph{nature théorique} du caractère qui compte : la taille reste un caractère continu. À l'inverse, le nombre d'enfants d'un ménage est fondamentalement discret : on ne peut pas avoir $2{,}7$ enfants !
\end{attention}

\courssub{Effectifs et fréquences}

Une fois les données recueillies, la première chose à faire est de les \emph{dépouiller} : compter combien de fois chaque valeur apparaît.

\begin{definition}[Effectif, effectif total]
Soit une série statistique portant sur une population de taille $N$.
\begin{itemize}
\item L'\cle{effectif} d'une valeur $x_i$ (ou d'une modalité, ou d'une classe), noté $n_i$, est le nombre d'individus pour lesquels le caractère prend cette valeur.
\item L'\cle{effectif total}, noté $N$, est le nombre total d'individus de la population. On a toujours :
\[ N = n_1 + n_2 + \cdots + n_k = \sum_{i=1}^{k} n_i \]
où $k$ est le nombre de valeurs (ou de classes) distinctes.
\end{itemize}
\end{definition}

L'effectif brut ne suffit pas toujours à comparer. Dire que 12 ménages de Mballa II dépensent moins de \SI{2000}{F} par jour, est-ce beaucoup ? Tout dépend du nombre de ménages interrogés : 12 sur 20, c'est énorme ; 12 sur 500, c'est marginal. D'où l'idée de \emph{ramener à la proportion}.

\begin{definition}[Fréquence]
La \cle{fréquence} d'une valeur $x_i$ d'effectif $n_i$ est le quotient :
\[ f_i = \frac{n_i}{N} \]
C'est un nombre compris entre $0$ et $1$. La \cle{fréquence en pourcentage} est $f_i \times 100$, exprimée en \%.
\end{definition}

\begin{propriete}[Somme des fréquences]
La somme des fréquences de toutes les valeurs d'une série statistique vaut $1$ (soit $100\,\%$) :
\[ \sum_{i=1}^{k} f_i = 1. \]
\end{propriete}

\begin{experience}[Démonstration guidée]
Complète la démonstration en une ligne :
\[ \sum_{i=1}^{k} f_i = \sum_{i=1}^{k} \frac{n_i}{N} = \frac{1}{N}\sum_{i=1}^{k} n_i = \frac{N}{N} = \ldots \]
La réponse est bien sûr $1$, puisque la somme des effectifs partiels redonne l'effectif total. Cette propriété, très simple, est un excellent \emph{test de cohérence} : si la somme de tes fréquences ne vaut pas $1$ (à un arrondi près), tu as fait une erreur de calcul.
\end{experience}

\courssub{Effectifs et fréquences cumulés croissants}

Revenons à la coopérative. Mme Etoundi veut répondre à une question de l'assemblée : « Combien de ménages du quartier dépensent \emph{moins de} \SI{4000}{F} par jour ? » Pour répondre à ce type de question — « moins de », « au plus », « en dessous de » —, les effectifs simples ne suffisent pas : il faut \emph{additionner au fur et à mesure}.

\begin{definition}[Effectifs et fréquences cumulés croissants]
On suppose les valeurs du caractère rangées dans l'ordre croissant : $x_1 < x_2 < \cdots < x_k$.
\begin{itemize}
\item L'\cle{effectif cumulé croissant} (ECC) de la valeur $x_i$ est la somme des effectifs des valeurs inférieures ou égales à $x_i$ :
\[ \mathrm{ECC}_i = n_1 + n_2 + \cdots + n_i. \]
Il répond à la question : « combien d'individus ont une valeur du caractère inférieure ou égale à $x_i$ ? »
\item La \cle{fréquence cumulée croissante} (FCC) de la valeur $x_i$ est :
\[ \mathrm{FCC}_i = f_1 + f_2 + \cdots + f_i = \frac{\mathrm{ECC}_i}{N}. \]
\end{itemize}
\end{definition}

\begin{aretenir}[Propriétés des cumuls]
\begin{itemize}
\item Le dernier ECC est toujours égal à l'effectif total $N$.
\item La dernière FCC est toujours égale à $1$ (soit $100\,\%$).
\item Les ECC et FCC sont des suites croissantes.
\end{itemize}
Ces deux dernières valeurs ($N$ et $1$) sont des points de contrôle obligatoires dans tes tableaux.
\end{aretenir}

\begin{attention}[Piège classique : confondre ECC et FCC]
C'est l'erreur la plus fréquente au baccalauréat. Retiens bien :
\begin{itemize}
\item l'ECC est un \emph{nombre d'individus} : il se termine par $N$ (par exemple $50$ ménages) ;
\item la FCC est une \emph{proportion} : elle se termine par $1$ (ou $100\,\%$).
\end{itemize}
Si ta dernière fréquence cumulée vaut $50$, tu as recopié la colonne des effectifs au lieu de diviser par $N$. De même, n'oublie jamais de vérifier que la dernière FCC vaut exactement $1$ (ou $100\,\%$) : c'est la signature d'un tableau correct.
\end{attention}

\courssub{Regroupement en classes}

Lorsque le caractère est continu — ou discret avec trop de valeurs différentes —, dresser la liste de toutes les valeurs serait illisible. Imagine le relevé des dépenses de 50 ménages : 50 nombres presque tous différents ! On regroupe alors les valeurs dans des intervalles appelés \cle{classes}.

\begin{definition}[Classe, amplitude, centre]
\begin{itemize}
\item Une \cle{classe} est un intervalle $[a\,;\,b[$ dans lequel on regroupe les valeurs du caractère. Par convention, la borne de gauche est incluse, la borne de droite exclue, sauf pour la dernière classe qui est fermée des deux côtés.
\item L'\cle{amplitude} de la classe $[a\,;\,b[$ est le nombre $b - a$.
\item Le \cle{centre} de la classe est :
\[ c = \frac{a + b}{2}. \]
\end{itemize}
\end{definition}

\begin{attention}[Le rôle du centre de classe]
Une fois les données regroupées, on \emph{perd} les valeurs exactes : on sait que 14 ménages dépensent entre \SI{2000}{F} et \SI{4000}{F}, mais pas combien exactement chacun dépense. Pour tous les calculs ultérieurs (moyenne, écart-type), on fait alors l'hypothèse que les valeurs d'une classe se répartissent autour de son \cle{centre} : on remplace chaque classe par son centre. Les résultats obtenus sont donc des \emph{valeurs approchées} — c'est le prix du regroupement.
\end{attention}

\begin{methode}[Construire un tableau à classes d'amplitudes égales]
\begin{enumerate}
\item Repérer la valeur minimale $x_{\min}$ et la valeur maximale $x_{\max}$ ; calculer l'étendue $e = x_{\max} - x_{\min}$.
\item Choisir le nombre de classes : en pratique, entre $4$ et $8$ classes donnent un tableau lisible. (Règle indicative : environ $\sqrt{N}$ classes.)
\item En déduire l'amplitude commune : $a \approx \dfrac{e}{\text{nombre de classes}}$, arrondie à une valeur simple.
\item Écrire les classes $[a_1\,;\,a_2[$, $[a_2\,;\,a_3[$, \ldots{} en veillant à ce qu'elles recouvrent toutes les données sans se chevaucher.
\item Compter les effectifs classe par classe (dépouillement), puis calculer les centres, les ECC, les fréquences et les FCC.
\item Vérifier : somme des effectifs $= N$, dernier ECC $= N$, dernière FCC $= 1$.
\end{enumerate}
\end{methode}

\courssub{Exemple chiffré complet : les notes d'une classe de Terminale A}

\begin{exemplebox}[Notes de 40 élèves de Tle A]
Voici les notes sur 20 obtenues par les 40 élèves d'une classe de Terminale A à un devoir de mathématiques, regroupées en classes $[0\,;\,5[$, $[5\,;\,10[$, $[10\,;\,15[$, $[15\,;\,20]$ :
\begin{center}
\begin{tabularx}{\linewidth}{|l|*{5}{>{\centering\arraybackslash}X|}}
\hline
\rowcolor{popblueL} \textbf{Classe} & \textbf{Centre} $c_i$ & \textbf{Effectif} $n_i$ & \textbf{ECC} & \textbf{Fréquence} $f_i$ & \textbf{FCC} \\
\hline
$[0\,;\,5[$ & $2{,}5$ & $4$ & $4$ & $0{,}10$ & $0{,}10$ \\
\hline
$[5\,;\,10[$ & $7{,}5$ & $11$ & $15$ & $0{,}275$ & $0{,}375$ \\
\hline
$[10\,;\,15[$ & $12{,}5$ & $18$ & $33$ & $0{,}45$ & $0{,}825$ \\
\hline
$[15\,;\,20]$ & $17{,}5$ & $7$ & $40$ & $0{,}175$ & $1$ \\
\hline
\rowcolor{popgoldL} \textbf{Total} & & $40$ & & $1$ & \\
\hline
\end{tabularx}
\end{center}
Lisons ce tableau ensemble :
\begin{itemize}
\item Les centres : $c_1 = \dfrac{0+5}{2} = 2{,}5$ ; $c_2 = \dfrac{5+10}{2} = 7{,}5$ ; etc.
\item Les ECC : $4$ ; puis $4 + 11 = 15$ ; puis $15 + 18 = 33$ ; enfin $33 + 7 = 40 = N$. \checkmark
\item Les fréquences : $f_2 = \dfrac{11}{40} = 0{,}275$, soit $27{,}5\,\%$.
\item Les FCC : $0{,}10$ ; $0{,}10 + 0{,}275 = 0{,}375$ ; $0{,}375 + 0{,}45 = 0{,}825$ ; $0{,}825 + 0{,}175 = 1$. \checkmark
\end{itemize}
Interprétation : $33$ élèves sur $40$, soit $82{,}5\,\%$ de la classe, ont une note inférieure à $15$. Le professeur constate que seuls $17{,}5\,\%$ des élèves atteignent $15$ ou plus : le devoir était difficile !
\end{exemplebox}

\courssub{Un tableau statistique complet en situation réelle}

Terminons par l'exemple qui servira de fil conducteur dans toute la leçon : les dépenses journalières (en centaines de francs) de $50$ ménages du quartier Mballa II, relevées par la coopérative.

\begin{popfigure}
\begin{center}
\begin{tabularx}{\linewidth}{|l|*{5}{>{\centering\arraybackslash}X|}}
\hline
\rowcolor{popblueL} \textbf{Dépense (cent. F)} & \textbf{Centre} $c_i$ & \textbf{Effectif} $n_i$ & \textbf{ECC} & \textbf{Fréquence} $f_i$ & \textbf{FCC} \\
\hline
$[10\,;\,20[$ & $15$ & $6$ & $6$ & $0{,}12$ & $0{,}12$ \\
\hline
$[20\,;\,30[$ & $25$ & $14$ & $20$ & $0{,}28$ & $0{,}40$ \\
\hline
$[30\,;\,40[$ & $35$ & $17$ & $37$ & $0{,}34$ & $0{,}74$ \\
\hline
$[40\,;\,50[$ & $45$ & $9$ & $46$ & $0{,}18$ & $0{,}92$ \\
\hline
$[50\,;\,60]$ & $55$ & $4$ & $50$ & $0{,}08$ & $1$ \\
\hline
\rowcolor{popgoldL} \textbf{Total} & & $50$ & & $1$ & \\
\hline
\end{tabularx}
\end{center}
\legende{Dépenses journalières de 50 ménages de Mballa II (en centaines de francs). Vérifications : somme des effectifs $= 50 = N$ ; dernier ECC $= 50$ ; dernière FCC $= 1$.}
\end{popfigure}

Grâce à ce tableau, Mme Etoundi peut répondre instantanément à de nombreuses questions : $40\,\%$ des ménages dépensent moins de \SI{3000}{F} par jour ; $37$ ménages sur $50$ dépensent moins de \SI{4000}{F}. Les données, autrefois informes, sont devenues \emph{parlantes}.

\begin{exoresolu}[Construire un tableau statistique complet]
Énoncé. Une boutique de crédit téléphonique de Bafoussam a relevé le nombre de recharges vendues chaque jour pendant 30 jours : les ventes vont de 12 à 58 recharges. On décide de regrouper les données en 5 classes d'amplitude 10 : $[10\,;\,20[$, $[20\,;\,30[$, $[30\,;\,40[$, $[40\,;\,50[$, $[50\,;\,60]$. Le dépouillement donne les effectifs respectifs : $3$, $7$, $11$, $6$, $3$. Construire le tableau complet (centres, ECC, fréquences, FCC) et déterminer le pourcentage de jours où la boutique a vendu moins de 40 recharges.
\tcblower
Solution. L'effectif total est $N = 3 + 7 + 11 + 6 + 3 = 30$. Les centres sont $15$, $25$, $35$, $45$, $55$. On calcule ensuite les cumuls et les fréquences :
\begin{align*}
\text{ECC} &: \quad 3\,;\quad 3+7=10\,;\quad 10+11=21\,;\quad 21+6=27\,;\quad 27+3=30 \checkmark\\
f_i &: \quad \frac{3}{30}=0{,}10\,;\quad \frac{7}{30}\approx 0{,}233\,;\quad \frac{11}{30}\approx 0{,}367\,;\quad \frac{6}{30}=0{,}20\,;\quad \frac{3}{30}=0{,}10\\
\text{FCC} &: \quad 0{,}10\,;\quad 0{,}333\,;\quad 0{,}70\,;\quad 0{,}90\,;\quad 1 \checkmark
\end{align*}
Le nombre de jours avec moins de 40 recharges est l'ECC de la classe $[30\,;\,40[$, soit $21$ jours, c'est-à-dire une FCC de $0{,}70$ : la boutique a vendu moins de 40 recharges pendant $70\,\%$ des jours.
\end{exoresolu}

\begin{savaistu}[Les statistiques au service du développement]
Au Cameroun, l'Institut National de la Statistique (INS) réalise régulièrement des enquêtes auprès des ménages (enquêtes ECAM) pour mesurer la pauvreté, la consommation ou l'accès à l'électricité. Chaque enquête repose exactement sur les notions de cette section : une population (les ménages camerounais), des individus, des caractères (revenu, taille du ménage, source d'énergie), des effectifs et des fréquences. Les décisions publiques — où construire une école, comment subventionner le manioc de la coopérative de Mballa II — s'appuient sur ces tableaux statistiques. Maîtriser ce vocabulaire, c'est déjà comprendre comment un pays se décrit lui-même.
\end{savaistu}

\begin{exercice}[À toi de jouer]
On a relevé les tailles (en cm) de 36 élèves d'un lycée de Douala, regroupées en classes : $[150\,;\,160[$ : $5$ élèves ; $[160\,;\,170[$ : $13$ élèves ; $[170\,;\,180[$ : $12$ élèves ; $[180\,;\,190]$ : $6$ élèves.
\begin{enumerate}
\item Préciser la population, l'individu et la nature du caractère étudié.
\item Construire le tableau complet : centres, effectifs, ECC, fréquences (en \%, arrondies à $0{,}1$), FCC.
\item Combien d'élèves mesurent moins de \SI{170}{cm} ? Quel pourcentage cela représente-t-il ?
\end{enumerate}
\end{exercice}

\begin{corrige}[Exercice 1]
1. Population : les 36 élèves du lycée ; individu : un élève ; caractère : la taille, quantitative continue.

2. $N = 5 + 13 + 12 + 6 = 36$. Centres : $155$, $165$, $175$, $185$.
\begin{center}
\begin{tabularx}{\linewidth}{|l|*{5}{>{\centering\arraybackslash}X|}}
\hline
\rowcolor{popblueL} \textbf{Classe} & \textbf{Centre} & $n_i$ & \textbf{ECC} & $f_i$ (\%) & \textbf{FCC} \\
\hline
$[150\,;\,160[$ & $155$ & $5$ & $5$ & $13{,}9$ & $0{,}139$ \\
\hline
$[160\,;\,170[$ & $165$ & $13$ & $18$ & $36{,}1$ & $0{,}50$ \\
\hline
$[170\,;\,180[$ & $175$ & $12$ & $30$ & $33{,}3$ & $0{,}833$ \\
\hline
$[180\,;\,190]$ & $185$ & $6$ & $36$ & $16{,}7$ & $1$ \\
\hline
\rowcolor{popgoldL} \textbf{Total} & & $36$ & & $100$ & \\
\hline
\end{tabularx}
\end{center}
Vérifications : dernier ECC $= 36 = N$ ; dernière FCC $= 1$. \checkmark

3. Les élèves mesurant moins de \SI{170}{cm} sont ceux des deux premières classes : ECC $= 18$ élèves, soit $50\,\%$ de l'effectif total.
\end{corrige}

\begin{aretenir}[L'essentiel de la section]
\begin{itemize}
\item Une étude statistique porte sur une \cle{population}, composée d'\cle{individus}, dont on observe un \cle{caractère} (qualitatif, quantitatif discret ou continu).
\item L'\cle{effectif} $n_i$ compte, la \cle{fréquence} $f_i = \dfrac{n_i}{N}$ proportionne ; la somme des fréquences vaut $1$.
\item Les \cle{ECC} et \cle{FCC} répondent aux questions « moins de » : le dernier ECC vaut $N$, la dernière FCC vaut $1$.
\item Pour un caractère continu, on regroupe en \cle{classes} d'amplitudes égales et on remplace chaque classe par son \cle{centre} $c = \dfrac{a+b}{2}$ pour les calculs.
\end{itemize}
\end{aretenir}

Nos données sont désormais bien organisées. Mais un tableau de chiffres, même parfait, reste austère pour une assemblée générale. Dans la prochaine section, nous apprendrons à \emph{dessiner} les données : diagrammes circulaires, à bandes, à bâtons, histogrammes et polygones des fréquences cumulées. Mme Etoundi pourra alors présenter ses chiffres en un coup d'œil.

\coursec{Représentations graphiques}

Dans la section précédente, nous avons appris à organiser des données brutes dans des tableaux : effectifs, effectifs cumulés, fréquences, fréquences cumulées. Un tableau est précis, mais il parle peu à l'œil. Or, dans la vie quotidienne — résultats du Baccalauréat publiés dans les journaux, répartition du budget de l'État, parts de marché de MTN et Orange —, ce sont les \cle{graphiques} qui convainquent et qui font comprendre une situation en un coup d'œil. Cette section est consacrée aux cinq représentations graphiques au programme : le diagramme circulaire, le diagramme à bandes, le diagramme à bâtons, l'histogramme et le polygone des fréquences cumulées. Pour chacun, nous répondrons à trois questions : à quelle situation s'applique-t-il ? comment le construit-on ? que peut-on y lire ?

\courssub{Le diagramme circulaire}

\textbf{Intuition.} Imagine un gâteau d'anniversaire à partager entre quatre personnes selon leur appétit : celui qui a le plus faim reçoit la plus grande part. Le diagramme circulaire fonctionne exactement ainsi : le disque entier représente l'effectif total, et chaque modalité reçoit un \emph{secteur angulaire} dont l'angle est proportionnel à son effectif.

\begin{definition}[Diagramme circulaire]
Le \cle{diagramme circulaire} d'une série statistique est un disque partagé en secteurs angulaires tels que l'angle de chaque secteur est proportionnel à l'effectif (ou à la fréquence) de la modalité correspondante. Si $f_i$ est la fréquence d'une modalité d'effectif $n_i$, l'angle de son secteur est :
\[ \alpha_i = f_i \times 360^{\circ} = \frac{n_i}{N}\times 360^{\circ}, \]
où $N$ désigne l'effectif total.
\end{definition}

\begin{methode}[Construire un diagramme circulaire]
\begin{enumerate}
\item Compléter le tableau avec la colonne des fréquences $f_i = \dfrac{n_i}{N}$.
\item Calculer chaque angle : $\alpha_i = f_i \times 360^{\circ}$ (arrondir au degré près si nécessaire, en veillant à ce que la somme reste $360^{\circ}$).
\item Tracer un cercle, puis reporter chaque angle au rapporteur, les secteurs les uns à la suite des autres.
\item Colorier les secteurs et ajouter une légende claire.
\end{enumerate}
\end{methode}

\begin{exemplebox}[Budget mensuel d'une famille de Yaoundé]
Une famille répartit son budget mensuel de \SI{200000}{F} ainsi : nourriture \SI{90000}{F}, loyer \SI{50000}{F}, scolarité \SI{35000}{F}, transport \SI{25000}{F}. Calculons les angles :
\begin{align*}
\text{Nourriture : } & \frac{90000}{200000}\times 360^{\circ} = 162^{\circ}, &
\text{Loyer : } & \frac{50000}{200000}\times 360^{\circ} = 90^{\circ},\\
\text{Scolarité : } & \frac{35000}{200000}\times 360^{\circ} = 63^{\circ}, &
\text{Transport : } & \frac{25000}{200000}\times 360^{\circ} = 45^{\circ}.
\end{align*}
Vérification : $162 + 90 + 63 + 45 = 360^{\circ}$. Le diagramme est tracé ci-dessous.
\end{exemplebox}

\begin{popfigure}
\begin{center}
\begin{tikzpicture}[scale=1.1]
% Nourriture 162°, Loyer 90°, Scolarité 63°, Transport 45°
\fill[popblueL] (0,0) -- (90:2.4) arc (90:252:2.4) -- cycle;
\fill[poporangeL] (0,0) -- (252:2.4) arc (252:342:2.4) -- cycle;
\fill[popgreenL] (0,0) -- (342:2.4) arc (342:405:2.4) -- cycle;
\fill[poppinkL] (0,0) -- (405:2.4) arc (405:450:2.4) -- cycle;
\draw[thick] (0,0) circle (2.4);
\draw[thick] (0,0) -- (90:2.4);
\draw[thick] (0,0) -- (252:2.4);
\draw[thick] (0,0) -- (342:2.4);
\draw[thick] (0,0) -- (405:2.4);
\node at (171:1.55) {\small Nourriture};
\node at (171:1.05) {\small $162^{\circ}$};
\node at (297:1.6) {\small Loyer};
\node at (297:1.15) {\small $90^{\circ}$};
\node at (13.5:1.65) {\small Scolarité};
\node at (13.5:1.1) {\small $63^{\circ}$};
\node at (67.5:1.65) {\small Transport};
\node at (67.5:1.1) {\small $45^{\circ}$};
\end{tikzpicture}
\end{center}
\legende{Diagramme circulaire de la répartition du budget familial mensuel.}
\end{popfigure}

\begin{attention}[Vérification obligatoire]
La somme des angles doit être exactement $360^{\circ}$. Si les arrondis donnent $359^{\circ}$ ou $361^{\circ}$, ajuste le plus grand secteur d'un degré. Autre erreur fréquente : calculer l'angle avec l'effectif au lieu de la fréquence — c'est bien $\alpha_i = f_i \times 360^{\circ}$, et non $n_i \times 360^{\circ}$ !
\end{attention}

\courssub{Le diagramme à bandes}

Le diagramme à bandes repose sur la même idée de proportionnalité, mais le « gâteau » est remplacé par une bande rectangulaire (ou par un demi-disque pour le diagramme semi-circulaire).

\begin{definition}[Diagramme à bandes]
Dans un \cle{diagramme à bandes rectangulaires}, on partage un rectangle de longueur fixe $L$ (souvent \SI{10}{cm}) en bandes dont les longueurs sont proportionnelles aux effectifs : la longueur de la bande de la modalité $i$ est $\ell_i = f_i \times L$. Dans un \cle{diagramme semi-circulaire}, les angles sont proportionnels aux effectifs avec un total de $180^{\circ}$ : $\alpha_i = f_i \times 180^{\circ}$.
\end{definition}

\begin{exemplebox}[Répartition des élèves d'une Terminale A]
Une classe de 60 élèves compte 24 élèves ayant choisi l'option Littérature, 21 l'option Histoire-Géographie et 15 l'option Philosophie renforcée. Avec une bande de \SI{12}{cm} :
\[ \ell_1 = \frac{24}{60}\times 12 = \SI{4,8}{cm}, \quad \ell_2 = \frac{21}{60}\times 12 = \SI{4,2}{cm}, \quad \ell_3 = \frac{15}{60}\times 12 = \SI{3}{cm}. \]
Vérification : $4{,}8 + 4{,}2 + 3 = 12$ cm. On trace un rectangle de \SI{12}{cm} de long et l'on découpe les trois bandes côte à côte, avec une légende.
\end{exemplebox}

\courssub{Le diagramme à bâtons : le roi des séries discrètes}

\textbf{Intuition.} Quand le caractère est \emph{quantitatif discret} (nombre d'enfants, nombre de buts marqués, note entière), chaque valeur est isolée sur l'axe des abscisses. On dresse alors au-dessus de chaque valeur un « bâton » vertical dont la hauteur est l'effectif : plus le bâton est haut, plus la valeur est fréquente.

\begin{definition}[Diagramme à bâtons]
Le \cle{diagramme à bâtons} d'une série discrète est obtenu en plaçant les valeurs $x_i$ du caractère en abscisse et en traçant, au-dessus de chaque $x_i$, un segment vertical de hauteur égale à l'effectif $n_i$ (ou à la fréquence $f_i$).
\end{definition}

\begin{exemplebox}[Nombre d'enfants par ménage]
Une enquête dans un quartier de Douala portant sur 50 ménages donne :
\begin{center}
\begin{tabularx}{\linewidth}{|l|*{6}{X|}}\hline
\rowcolor{popblueL} Nombre d'enfants $x_i$ & 0 & 1 & 2 & 3 & 4 & 5 \\\hline
Effectif $n_i$ & 5 & 9 & 14 & 12 & 7 & 3 \\\hline
\end{tabularx}
\end{center}
Vérifions l'effectif total : $5+9+14+12+7+3 = 50$. Le diagramme à bâtons correspondant est tracé ci-dessous. On lit immédiatement le mode : $x = 2$ enfants (le bâton le plus haut, d'effectif 14).
\end{exemplebox}

\begin{popfigure}
\begin{center}
\begin{tikzpicture}
\begin{axis}[
    width=0.85\linewidth, height=6cm,
    xlabel={Nombre d'enfants}, ylabel={Effectif},
    xmin=-0.5, xmax=5.5, ymin=0, ymax=16,
    xtick={0,1,2,3,4,5}, ytick={0,2,4,6,8,10,12,14,16},
    ymajorgrids=true, grid style={popline!40},
    axis lines=left, tick label style={font=\small}, label style={font=\small}
]
\addplot+[ycomb, popblue, very thick, mark=*, mark options={fill=popblue}] coordinates {(0,5) (1,9) (2,14) (3,12) (4,7) (5,3)};
\end{axis}
\end{tikzpicture}
\end{center}
\legende{Diagramme à bâtons du nombre d'enfants par ménage (50 ménages).}
\end{popfigure}

\courssub{L'histogramme : le roi des séries continues}

\textbf{Intuition.} Quand le caractère est \emph{continu} (taille, masse, durée d'appel), les données sont regroupées en \emph{classes} $[a\,;b[$. On ne peut plus tracer un bâton par valeur : on dessine des \emph{rectangles juxtaposés}, un par classe. Mais attention, une subtilité capitale apparaît : ce n'est pas la hauteur qui compte, c'est l'\emph{aire}.

\begin{definition}[Histogramme]
L'\cle{histogramme} d'une série groupée en classes est constitué de rectangles juxtaposés dont les bases sont les amplitudes des classes et dont les \textbf{aires sont proportionnelles aux effectifs} (ou aux fréquences). La hauteur du rectangle d'une classe est donc proportionnelle à la \emph{densité d'effectif} :
\[ h_i = k\times\frac{n_i}{a_i}, \]
où $a_i$ est l'amplitude de la classe et $k$ un coefficient d'échelle commun à toutes les classes. \textbf{Cas particulier fondamental :} si toutes les classes ont la \emph{même amplitude}, les hauteurs sont proportionnelles aux effectifs et l'on peut prendre directement $h_i = n_i$.
\end{definition}

\begin{attention}[Le piège classique du Bac]
Dans un histogramme, ce sont les \cle{aires}, et non les hauteurs, qui sont proportionnelles aux effectifs. Si une classe $[20\,;40[$ a une amplitude double de la classe $[0\,;10[$, son rectangle doit être \emph{deux fois moins haut} à effectif égal. Oublier cette règle quand les amplitudes sont inégales est l'erreur la plus fréquente — et la plus sanctionnée. En revanche, si toutes les amplitudes sont égales, prendre les hauteurs égales aux effectifs convient parfaitement.
\end{attention}

\begin{exemplebox}[Durées d'appels chez un opérateur]
On a relevé la durée (en minutes) de 100 appels passés depuis une agence :
\begin{center}
\begin{tabularx}{\linewidth}{|l|*{5}{X|}}\hline
\rowcolor{popblueL} Durée (min) & $[0\,;2[$ & $[2\,;4[$ & $[4\,;6[$ & $[6\,;8[$ & $[8\,;10[$ \\\hline
Effectif & 12 & 25 & 33 & 20 & 10 \\\hline
\end{tabularx}
\end{center}
Vérification : $12+25+33+20+10 = 100$. Toutes les classes ont la même amplitude (2 min) : les hauteurs des rectangles sont donc égales aux effectifs.
\end{exemplebox}

\begin{popfigure}
\begin{center}
\begin{tikzpicture}
\begin{axis}[
    width=0.85\linewidth, height=6cm,
    xlabel={Durée (min)}, ylabel={Effectif},
    xmin=0, xmax=10, ymin=0, ymax=38,
    xtick={0,2,4,6,8,10}, ytick={0,5,10,15,20,25,30,35},
    ymajorgrids=true, grid style={popline!40},
    axis lines=left, tick label style={font=\small}, label style={font=\small}
]
\addplot+[ybar interval, fill=popblue!45, draw=popblue, thick] coordinates {(0,12) (2,25) (4,33) (6,20) (8,10) (10,0)};
\end{axis}
\end{tikzpicture}
\end{center}
\legende{Histogramme des durées d'appels (classes d'amplitudes égales).}
\end{popfigure}

\begin{savaistu}[Florence Nightingale, infirmière et statisticienne]
L'histogramme et ses variantes ont été popularisés au XIX\textsuperscript{e} siècle par Florence Nightingale. Pendant la guerre de Crimée, cette infirmière britannique représenta graphiquement les causes de mortalité des soldats et démontra, chiffres à l'appui, que la plupart des décès étaient dus aux mauvaises conditions d'hygiène et non aux blessures de guerre. Ses diagrammes convainquirent les autorités sanitaires de réformer les hôpitaux militaires : c'est l'un des premiers grands triomphes de la statistique au service de la santé publique. Elle fut d'ailleurs la première femme élue membre de la Royal Statistical Society.
\end{savaistu}

\courssub{Le polygone des fréquences cumulées croissantes}

\textbf{Intuition.} Jusqu'ici, nos graphiques répondaient à la question « combien ? ». Le polygone cumulé répond à une autre question, essentielle : « combien de données sont \emph{inférieures} à telle valeur ? ». C'est l'outil qui permettra de lire graphiquement la médiane et les quartiles, sans aucun calcul.

\begin{definition}[Polygone des fréquences cumulées croissantes]
Pour une série groupée en classes, le \cle{polygone des fréquences cumulées croissantes} (FCC) est obtenu en plaçant dans un repère les points de coordonnées $(b_i\,;\,F_i)$, où $b_i$ est la \textbf{borne supérieure} de la classe numéro $i$ et $F_i$ la fréquence cumulée croissante correspondante, puis en reliant ces points par des segments de droite. On ajoute le point $(b_0\,;\,0)$, où $b_0$ est la borne inférieure de la première classe. (Le polygone des effectifs cumulés croissants, ECC, se construit de même avec les effectifs cumulés.)
\end{definition}

\begin{methode}[Construire le polygone des FCC]
\begin{enumerate}
\item Dresser le tableau des classes avec les fréquences cumulées croissantes $F_i$.
\item Placer le point de départ $(b_0\,;\,0)$ : borne inférieure de la première classe, fréquence nulle.
\item Pour chaque classe, placer le point $(\text{borne supérieure}\,;\,F_i)$.
\item Relier les points par des segments, dans l'ordre croissant des abscisses.
\item Vérifier que la courbe obtenue est croissante et se termine à l'ordonnée $1$ (ou $100\,\%$).
\end{enumerate}
\end{methode}

\begin{exoresolu}[Polygone des FCC et lecture de la médiane]
\textbf{Énoncé.} On reprend les durées d'appels de l'exemple précédent :
$[0\,;2[$ : 12 ; $[2\,;4[$ : 25 ; $[4\,;6[$ : 33 ; $[6\,;8[$ : 20 ; $[8\,;10[$ : 10 (total $N = 100$).
\begin{enumerate}
\item Dresser le tableau des fréquences cumulées croissantes.
\item Tracer le polygone des FCC.
\item Lire graphiquement une valeur approchée de la médiane.
\end{enumerate}
\tcblower
\textbf{Solution.}
\textbf{1.} Les fréquences sont $0{,}12$ ; $0{,}25$ ; $0{,}33$ ; $0{,}20$ ; $0{,}10$. On cumule : $0{,}12$ ; $0{,}12+0{,}25 = 0{,}37$ ; $0{,}37+0{,}33 = 0{,}70$ ; $0{,}70+0{,}20 = 0{,}90$ ; $0{,}90+0{,}10 = 1$. D'où le tableau :
\begin{center}
\begin{tabularx}{\linewidth}{|l|*{5}{X|}}\hline
\rowcolor{popblueL} Borne supérieure $b_i$ & 2 & 4 & 6 & 8 & 10 \\\hline
FCC $F_i$ & 0,12 & 0,37 & 0,70 & 0,90 & 1,00 \\\hline
\end{tabularx}
\end{center}
\textbf{2.} On place les points $(0\,;\,0)$, $(2\,;\,0{,}12)$, $(4\,;\,0{,}37)$, $(6\,;\,0{,}70)$, $(8\,;\,0{,}90)$, $(10\,;\,1)$ et on les relie par des segments (figure ci-dessous).

\textbf{3.} La médiane est l'abscisse du point du polygone d'ordonnée $0{,}5$. On trace la droite horizontale d'équation $y = 0{,}5$ ; elle coupe le polygone sur le segment joignant $(4\,;\,0{,}37)$ à $(6\,;\,0{,}70)$. La verticale abaissée depuis ce point d'intersection donne $Me \approx 4{,}8$ min. Interprétation : environ la moitié des appels dure moins de 4,8 minutes.
\end{exoresolu}

\begin{popfigure}
\begin{center}
\begin{tikzpicture}
\begin{axis}[
    width=0.88\linewidth, height=7cm,
    xlabel={Durée (min)}, ylabel={Fréquence cumulée croissante},
    xmin=0, xmax=10, ymin=0, ymax=1.05,
    xtick={0,2,4,6,8,10}, ytick={0,0.1,0.2,0.3,0.4,0.5,0.6,0.7,0.8,0.9,1},
    ymajorgrids=true, grid style={popline!40},
    axis lines=left, tick label style={font=\small}, label style={font=\small}
]
\addplot[popblue, very thick, mark=*, mark options={fill=popblue}] coordinates {(0,0) (2,0.12) (4,0.37) (6,0.70) (8,0.90) (10,1)};
\addplot[poporange, dashed, thick] coordinates {(0,0.5) (4.8,0.5)};
\addplot[poporange, dashed, thick] coordinates {(4.8,0) (4.8,0.5)};
\addplot[mark=*, mark options={fill=poporange}, only marks, poporange] coordinates {(4.8,0.5)};
\node[poporange, font=\small, anchor=north] at (axis cs:4.8,0) {$Me \approx 4{,}8$};
\node[poporange, font=\small, anchor=south east] at (axis cs:0,0.5) {$0{,}5$};
\end{axis}
\end{tikzpicture}
\end{center}
\legende{Polygone des FCC et lecture graphique de la médiane.}
\end{popfigure}

\begin{propriete}[Lecture graphique de la médiane et des quartiles — admise]
Sur le polygone des fréquences cumulées croissantes :
\begin{itemize}
\item la \cle{médiane} $Me$ est l'abscisse du point d'ordonnée $0{,}5$ ;
\item le \cle{premier quartile} $Q_1$ est l'abscisse du point d'ordonnée $0{,}25$ ;
\item le \cle{troisième quartile} $Q_3$ est l'abscisse du point d'ordonnée $0{,}75$.
\end{itemize}
Cette propriété est admise ; elle se justifie graphiquement : le polygone des FCC indique, pour chaque valeur $x$, la proportion de données inférieures à $x$. Chercher la médiane, c'est chercher la valeur qui partage la série en deux moitiés, c'est-à-dire la valeur en dessous de laquelle se trouvent exactement $50\,\%$ des données : c'est précisément l'abscisse correspondant à l'ordonnée $0{,}5$. Le même raisonnement vaut pour $Q_1$ ($25\,\%$) et $Q_3$ ($75\,\%$). De même, les déciles se liraient aux ordonnées $0{,}1$ ; $0{,}2$ ; \ldots ; $0{,}9$.
\end{propriete}

\begin{attention}[Borne supérieure, pas centre de classe !]
L'erreur la plus fréquente est de placer les points au \emph{centre} des classes. Pour un polygone \emph{cumulé croissant}, on place toujours les points aux \textbf{bornes supérieures} des classes, car la fréquence cumulée $F_i$ signifie « proportion de données strictement inférieures à $b_i$ ». Et n'oublie pas le point de départ $(b_0\,;\,0)$, sans lequel la lecture graphique de $Q_1$ peut devenir impossible.
\end{attention}

\courssub{Bien choisir son diagramme}

Chaque diagramme a son terrain de prédilection. Le tableau suivant résume le choix à faire selon la nature du caractère et le message à faire passer.

\begin{center}
\begin{tabularx}{\linewidth}{|l|X|X|}\hline
\rowcolor{popblueL} \textbf{Diagramme} & \textbf{Nature du caractère} & \textbf{Ce qu'il montre le mieux} \\\hline
Circulaire / semi-circulaire & Qualitatif ou discret (peu de modalités) & Les parts d'un tout (répartition en \%) \\\hline
À bandes & Qualitatif ou discret & Comparaison de parts entre plusieurs groupes \\\hline
À bâtons & Quantitatif discret & Comparaison des effectifs, lecture du mode \\\hline
Histogramme & Quantitatif continu (classes) & Forme de la répartition, concentration des données \\\hline
Polygone des FCC & Quantitatif (surtout continu) & Lecture de la médiane, des quartiles, des déciles \\\hline
\end{tabularx}
\end{center}

\begin{aretenir}[L'essentiel de la section]
\begin{itemize}
\item Diagramme circulaire : angle $= f_i \times 360^{\circ}$ ; la somme des angles vaut $360^{\circ}$.
\item Diagramme à bâtons (série discrète) : hauteur du bâton $=$ effectif.
\item Histogramme (série continue) : \textbf{aires} proportionnelles aux effectifs ; si les amplitudes sont égales, les hauteurs égales aux effectifs conviennent.
\item Polygone des FCC : points $(\text{borne supérieure}\,;\,F_i)$ reliés, en partant de $(b_0\,;\,0)$ ; la médiane se lit à l'ordonnée $0{,}5$, $Q_1$ à $0{,}25$, $Q_3$ à $0{,}75$.
\item Le choix du diagramme dépend de la nature du caractère et de l'information recherchée.
\end{itemize}
\end{aretenir}

\begin{exercice}[Pour s'entraîner]
Les notes de mathématiques de 80 élèves de Terminale A sont regroupées ainsi : $[0\,;5[$ : 8 ; $[5\,;10[$ : 22 ; $[10\,;15[$ : 38 ; $[15\,;20]$ : 12.
\begin{enumerate}
\item Quel diagramme choisir pour représenter cette série ? Justifier.
\item Calculer les fréquences cumulées croissantes et tracer le polygone des FCC.
\item Déterminer graphiquement la médiane, $Q_1$ et $Q_3$, puis interpréter $Q_3$.
\item Calculer les angles du diagramme circulaire associé à ces quatre classes.
\end{enumerate}
\end{exercice}

\begin{corrige}[Exercice]
\textbf{1.} Le caractère est quantitatif continu, regroupé en classes d'amplitudes égales : l'\emph{histogramme} est le diagramme adapté (hauteurs $=$ effectifs).

\textbf{2.} Fréquences : $\dfrac{8}{80} = 0{,}10$ ; $\dfrac{22}{80} = 0{,}275$ ; $\dfrac{38}{80} = 0{,}475$ ; $\dfrac{12}{80} = 0{,}15$. FCC : $F(5) = 0{,}10$ ; $F(10) = 0{,}10 + 0{,}275 = 0{,}375$ ; $F(15) = 0{,}375 + 0{,}475 = 0{,}85$ ; $F(20) = 1$. On place les points $(0\,;\,0)$, $(5\,;\,0{,}10)$, $(10\,;\,0{,}375)$, $(15\,;\,0{,}85)$, $(20\,;\,1)$ et on les relie par des segments.

\textbf{3.} Lecture à l'ordonnée $0{,}5$ (sur le segment joignant $(10\,;\,0{,}375)$ à $(15\,;\,0{,}85)$) : $Me \approx 11{,}3$. À l'ordonnée $0{,}25$ (segment de $(5\,;\,0{,}10)$ à $(10\,;\,0{,}375)$) : $Q_1 \approx 7{,}7$. À l'ordonnée $0{,}75$ : $Q_3 \approx 13{,}9$. Interprétation de $Q_3$ : environ $75\,\%$ des élèves ont une note inférieure à $13{,}9$ (et donc environ $25\,\%$ ont une note supérieure à $13{,}9$).

\textbf{4.} Angles : $0{,}10 \times 360^{\circ} = 36^{\circ}$ ; $0{,}275 \times 360^{\circ} = 99^{\circ}$ ; $0{,}475 \times 360^{\circ} = 171^{\circ}$ ; $0{,}15 \times 360^{\circ} = 54^{\circ}$. Vérification : $36 + 99 + 171 + 54 = 360^{\circ}$.
\end{corrige}

Maintenant que tu sais représenter une série et lire ses grandes tendances sur un graphique, il est temps de résumer cette série par des nombres précis : c'est l'objet de la section suivante, consacrée aux paramètres de position — mode, moyenne et médiane.

\coursec{Paramètres de position : mode, moyenne, médiane}

Dans la section précédente, nous avons appris à \emph{voir} une série statistique : diagrammes à bâtons, histogrammes, polygones des fréquences cumulées. Mais un dessin, aussi parlant soit-il, ne suffit pas toujours. Quand le proviseur d'un lycée de Yaoundé veut comparer les résultats de deux classes de Terminale, ou quand une entreprise de Douala veut annoncer « le salaire typique » de ses employés, il faut résumer toute la série par \textbf{un seul nombre}. C'est exactement le rôle des \cle{paramètres de position} : le \cle{mode}, la \cle{moyenne} et la \cle{médiane}. Chacun répond à une question différente, et choisir le bon est déjà une compétence.

\courssub{Le mode : la valeur la plus fréquente}

\textbf{Intuition.} Un vendeur de chaussures au marché Mokolo ne commande pas la pointure « moyenne » de ses clients : il commande la pointure qu'on lui demande \emph{le plus souvent}. Cette valeur-là, c'est le mode.

\begin{definition}[Mode et classe modale]
Le \cle{mode} d'une série statistique discrète est la valeur du caractère qui possède l'\textbf{effectif le plus élevé} (ou, de façon équivalente, la plus grande fréquence). Une série peut avoir plusieurs modes (on dit qu'elle est \emph{bimodale} si deux valeurs dominent à égalité). Pour une série \textbf{groupée en classes}, on ne parle plus de mode mais de \cle{classe modale} : c'est la classe d'effectif maximal. Si les classes n'ont pas la même amplitude, c'est la classe dont la \textbf{densité} (effectif divisé par l'amplitude) est maximale qui est considérée comme la classe modale (celle dont le rectangle est le plus haut sur l'histogramme).
\end{definition}

\begin{exemplebox}[Lecture immédiate du mode]
On a relevé le nombre d'enfants de 50 ménages d'un quartier de Bafoussam :
\begin{center}
\begin{tabular}{|l|c|c|c|c|c|c|}
\hline
\rowcolor{popblueL} Nombre d'enfants $x_i$ & 0 & 1 & 2 & 3 & 4 & 5 \\
\hline
Effectif $n_i$ & 4 & 9 & 15 & 12 & 7 & 3 \\
\hline
\end{tabular}
\end{center}
L'effectif maximal est $15$, atteint pour la valeur $x_i = 2$. Le mode est donc $\boxed{2 \text{ enfants}}$. Sur le diagramme à bâtons, c'est tout simplement le \textbf{bâton le plus haut}.
\end{exemplebox}

\begin{attention}[Mode et classe modale : ne pas confondre]
Pour une série groupée (par exemple les tailles regroupées en classes $[150;160[$, $[160;170[$\ldots), on ne peut pas donner une valeur exacte du mode : on donne la \textbf{classe modale}, celle dont l'effectif est le plus grand (ou la densité la plus grande si amplitudes inégales). Écrire « le mode est $165$ cm » alors que les données sont groupées est une erreur : on ne connaît pas les valeurs individuelles à l'intérieur de la classe.
\end{attention}

\courssub{La moyenne : le partage équitable}

\textbf{Intuition.} Trois amis gagnent respectivement \num{2000}, \num{3000} et \num{7000} FCFA à un petit commerce. S'ils décident de tout mettre en commun puis de partager également, chacun reçoit $\frac{2000+3000+7000}{3} = \num{4000}$ FCFA. La moyenne, c'est exactement ce \emph{partage équitable} : la somme totale est redistribuée à parts égales.

\begin{definition}[Moyenne d'une série discrète]
Soit une série statistique discrète de valeurs $x_1, x_2, \ldots, x_p$, d'effectifs respectifs $n_1, n_2, \ldots, n_p$, et d'effectif total $N = n_1 + n_2 + \cdots + n_p$. La \cle{moyenne} de la série, notée $\bar{x}$, est :
\[ \bar{x} = \frac{n_1 x_1 + n_2 x_2 + \cdots + n_p x_p}{N} = \frac{\displaystyle\sum_{i=1}^{p} n_i x_i}{N}. \]
\end{definition}

\begin{attention}[Ne pas oublier les effectifs]
L'erreur classique consiste à calculer $\frac{x_1 + x_2 + \cdots + x_p}{p}$ en additionnant simplement les valeurs du caractère. C'est \textbf{faux} dès que les effectifs ne sont pas tous égaux : chaque valeur doit être \emph{pondérée} par son effectif. La moyenne d'une série statistique est une \cle{moyenne pondérée}.
\end{attention}

\begin{exemplebox}[Moyenne pondérée]
Reprenons les 50 ménages de Bafoussam :
\[ \bar{x} = \frac{4\times 0 + 9\times 1 + 15\times 2 + 12\times 3 + 7\times 4 + 3\times 5}{50} = \frac{0+9+30+36+28+15}{50} = \frac{118}{50} = 2{,}36. \]
En moyenne, un ménage compte $2{,}36$ enfants. Remarque que la moyenne n'est pas nécessairement une valeur de la série : aucun ménage n'a $2{,}36$ enfant, et c'est parfaitement normal.
\end{exemplebox}

\textbf{Cas d'une série groupée en classes.} Quand les données sont regroupées en classes $[a_i ; b_i[$, on ne connaît plus les valeurs exactes. On fait alors l'hypothèse que les valeurs de chaque classe sont \emph{réparties autour de son centre}.

\begin{methode}[Calculer la moyenne d'une série groupée]
\begin{enumerate}
\item Ajouter au tableau une colonne des \cle{centres de classes} : $c_i = \dfrac{a_i + b_i}{2}$.
\item Ajouter une colonne des produits $n_i c_i$.
\item Additionner les produits : $\sum n_i c_i$.
\item Diviser par l'effectif total : $\bar{x} = \dfrac{\sum n_i c_i}{N}$.
\end{enumerate}
La valeur obtenue est une \textbf{valeur approchée} de la moyenne réelle.
\end{methode}

\begin{attention}[Bornes ou centres ? Le piège du Bac]
L'erreur la plus fréquente consiste à calculer la moyenne d'une série groupée avec les \textbf{bornes} des classes (par exemple en utilisant $150$, $160$, $170$\ldots) au lieu des \textbf{centres} ($155$, $165$, $175$\ldots). Les bornes ne représentent aucune valeur réelle de la série ; seul le centre $c_i = \frac{a_i+b_i}{2}$ est un représentant légitime de la classe. Au Baccalauréat, la colonne des centres doit apparaître explicitement dans ton tableau : elle fait partie de la rédaction attendue.
\end{attention}

\begin{propriete}[Linéarité de la moyenne]
Soit une série statistique de valeurs $x_i$, d'effectifs $n_i$ et de moyenne $\bar{x}$. Si l'on transforme chaque valeur par une fonction affine $y_i = a x_i + b$ (avec $a$ et $b$ réels), alors la moyenne de la nouvelle série est :
\[ \bar{y} = a\,\bar{x} + b. \]
Cette propriété est \textbf{exigible} : sa démonstration peut être demandée au Baccalauréat.
\end{propriete}

\begin{experience}[Démonstration guidée de la linéarité de la moyenne]
On part de la définition de la moyenne de la série des $y_i$, et on progresse pas à pas.
\begin{enumerate}
\item \textbf{Écris la définition} : $\bar{y} = \dfrac{\sum_{i=1}^{p} n_i y_i}{N}$.
\item \textbf{Remplace} $y_i$ par $a x_i + b$ : $\bar{y} = \dfrac{\sum_{i=1}^{p} n_i (a x_i + b)}{N}$.
\item \textbf{Développe} le numérateur : $\sum n_i(ax_i + b) = a\sum n_i x_i + b\sum n_i$.
\item \textbf{Utilise deux faits clés} : $\sum_{i=1}^{p} n_i = N$ (effectif total) et $\sum_{i=1}^{p} n_i x_i = N\bar{x}$ (définition de $\bar{x}$).
\item \textbf{Conclus} :
\[ \bar{y} = \frac{a\,N\bar{x} + b\,N}{N} = \frac{N(a\bar{x} + b)}{N} = a\bar{x} + b. \qquad \blacksquare \]
\end{enumerate}
Retiens la démarche : partir de la définition, développer, puis faire apparaître $\sum n_i = N$ et $\sum n_i x_i = N\bar{x}$.
\end{experience}

\begin{exemplebox}[Application : conversion et prime]
Les notes sur 20 d'une classe ont une moyenne $\bar{x} = 11{,}4$. Le professeur décide de transformer chaque note par $y = 0{,}5x + 4$ pour harmoniser avec une autre épreuve. La nouvelle moyenne est immédiatement :
\[ \bar{y} = 0{,}5 \times 11{,}4 + 4 = 5{,}7 + 4 = 9{,}7. \]
Inutile de recalculer toute la série ! De même, si chaque employé d'une entreprise reçoit une prime de \num{10000} FCFA, le salaire moyen augmente exactement de \num{10000} FCFA.
\end{exemplebox}

\courssub{La médiane : la valeur qui partage en deux}

\textbf{Intuition.} Rangeons les élèves d'une classe du plus petit au plus grand. La médiane est la taille de l'élève « du milieu » : il y a autant d'élèves plus petits que de plus grands. Contrairement à la moyenne, elle ne dépend pas de la \emph{grandeur} des valeurs extrêmes, seulement de leur \emph{rang}.

\begin{definition}[Médiane d'une série discrète]
La \cle{médiane} d'une série statistique, notée $Me$, est une valeur du caractère telle qu'au moins $50\,\%$ des valeurs lui sont inférieures ou égales et au moins $50\,\%$ lui sont supérieures ou égales. En pratique, après avoir rangé les $N$ valeurs dans l'\textbf{ordre croissant} :
\begin{itemize}
\item si $N$ est \textbf{impair} : $Me$ est la valeur de rang $\dfrac{N+1}{2}$ ;
\item si $N$ est \textbf{pair} : $Me$ est la demi-somme des valeurs de rangs $\dfrac{N}{2}$ et $\dfrac{N}{2}+1$.
\end{itemize}
\end{definition}

\begin{exemplebox}[Médiane selon la parité de $N$]
\textbf{Cas impair.} Notes rangées de $N = 9$ élèves : $5 ; 7 ; 8 ; 10 ; 11 ; 12 ; 14 ; 15 ; 18$. Le rang cherché est $\frac{9+1}{2} = 5$ : $Me = 11$ (la 5\up{e} valeur). Il y a bien 4 notes en dessous et 4 au-dessus.

\textbf{Cas pair.} Notes rangées de $N = 10$ élèves : $5 ; 7 ; 8 ; 10 ; 11 ; 12 ; 14 ; 15 ; 17 ; 18$. On prend les valeurs de rangs $\frac{10}{2} = 5$ et $6$, soit $11$ et $12$ : $Me = \dfrac{11+12}{2} = 11{,}5$.
\end{exemplebox}

\begin{attention}[Ranger avant tout !]
La médiane n'a de sens qu'après avoir \textbf{rangé les valeurs dans l'ordre croissant}. Prendre « la valeur du milieu » d'une liste non triée est l'erreur la plus fréquente — et la plus sanctionnée. Pour une série présentée par effectifs, utilise les \cle{effectifs cumulés croissants} pour repérer le rang cherché.
\end{attention}

\textbf{Cas d'une série groupée en classes.} On ne connaît plus les valeurs individuelles : on détermine d'abord la \cle{classe médiane}, c'est-à-dire la première classe dont l'effectif cumulé croissant atteint ou dépasse $\frac{N}{2}$. Ensuite, le programme de Terminale A prévoit une \textbf{lecture graphique} de la médiane sur le polygone des fréquences cumulées croissantes (FCC) : on lit l'abscisse du point d'ordonnée $50\,\%$. L'interpolation linéaire par le calcul est \emph{hors programme}.

\begin{exoresolu}[Salaires dans une PME de Douala : moyenne contre médiane]
Une PME de Douala emploie 25 salariés. Les salaires mensuels (en milliers de FCFA), rangés par classes, sont :
\begin{center}
\begin{tabular}{|l|c|c|c|c|c|}
\hline
\rowcolor{popblueL} Salaire (milliers FCFA) & $[50;100[$ & $[100;150[$ & $[150;200[$ & $[200;250[$ & $[250;500[$ \\
\hline
Effectif $n_i$ & 8 & 9 & 5 & 2 & 1 \\
\hline
\end{tabular}
\end{center}
\begin{enumerate}
\item Calculer le salaire moyen.
\item Déterminer la classe médiane, puis estimer la médiane par lecture graphique sur le polygone des FCC.
\item Le salaire de la dernière classe est celui du directeur (\num{450000} FCFA). Commenter.
\end{enumerate}
\tcblower
\textbf{1. Salaire moyen.} On ajoute la colonne des centres :
\begin{center}
\begin{tabular}{|l|c|c|c|}
\hline
\rowcolor{popblueL} Classe & Centre $c_i$ & $n_i$ & $n_i c_i$ \\
\hline
$[50;100[$ & 75 & 8 & 600 \\
$[100;150[$ & 125 & 9 & \num{1125} \\
$[150;200[$ & 175 & 5 & 875 \\
$[200;250[$ & 225 & 2 & 450 \\
$[250;500[$ & 375 & 1 & 375 \\
\hline
\textbf{Total} & & $N = 25$ & $\sum n_i c_i = \num{3425}$ \\
\hline
\end{tabular}
\end{center}
\[ \bar{x} = \frac{\num{3425}}{25} = 137 \quad \text{soit environ } \num{137000} \text{ FCFA}. \]

\textbf{2. Médiane.} Effectifs cumulés croissants : $8$ ; $17$ ; $22$ ; $24$ ; $25$. On a $\frac{N}{2} = 12{,}5$. Le premier cumul qui atteint $12{,}5$ est $17$ : la \textbf{classe médiane} est $[100;150[$. Sur le polygone des FCC (fréquences cumulées en \% : $32\,\%$ à la borne 100 ; $68\,\%$ à la borne 150 ; $88\,\%$ à la borne 200 ; $96\,\%$ à la borne 250 ; $100\,\%$ à la borne 500), l'abscisse du point d'ordonnée $50\,\%$ se lit sur le segment reliant $(100;32)$ à $(150;68)$, aux environs de $125$ : $Me \approx \num{125000}$ FCFA (voir figure ci-dessous).

\textbf{3. Commentaire.} La moyenne ($\num{137000}$ FCFA) est nettement supérieure à la médiane ($\approx \num{125000}$ FCFA) : le salaire très élevé du directeur \emph{tire la moyenne vers le haut}, alors que la médiane n'en est pas affectée. En réalité, 17 salariés sur 25 gagnent moins de \num{150000} FCFA : la médiane décrit mieux le « salaire typique » de cette entreprise.
\end{exoresolu}

\begin{popfigure}
\begin{center}
\begin{tikzpicture}[x=0.028cm,y=0.055cm]
\draw[->,thick,popdark] (0,0) -- (520,0) node[below left,font=\small]{salaire (milliers FCFA)};
\draw[->,thick,popdark] (0,0) -- (0,108) node[below left,font=\small]{FCC (\%)};
\foreach \x in {50,100,150,200,250,500} \draw (\x,1.5) -- (\x,-1.5) node[below,font=\scriptsize]{\x};
\foreach \y in {20,40,60,80,100} \draw (3,\y) -- (-3,\y) node[left,font=\scriptsize]{\y};
% Polygone des FCC : points aux bornes supérieures des classes (et borne inférieure de la 1ère à 0%)
\draw[very thick,popblue] (50,0) -- (100,32) -- (150,68) -- (200,88) -- (250,96) -- (500,100);
\fill[popblue] (50,0) circle (2.5) (100,32) circle (2.5) (150,68) circle (2.5) (200,88) circle (2.5) (250,96) circle (2.5) (500,100) circle (2.5);
% Lecture graphique de la médiane (y=50)
\draw[dashed,thick,poporange] (0,50) -- (125,50) -- (125,0);
\fill[poporange] (125,50) circle (3);
\node[below,poporange,font=\small] at (125,-6) {$Me \approx 125$};
\node[left,poporange,font=\small] at (-4,50) {$50$};
\end{tikzpicture}
\end{center}
\legende{Lecture graphique de la médiane sur le polygone des fréquences cumulées croissantes : on trace l'horizontale d'ordonnée $50\,\%$, puis on lit l'abscisse du point d'intersection avec le polygone. Ici $Me \approx 125$ (soit \num{125000} FCFA).}
\end{popfigure}

\courssub{Comparer les trois paramètres sur une même série}

Résumons sur une série discrète complète. On a relevé le nombre de buts marqués par match lors de 20 journées d'un championnat local : $0 ; 0 ; 1 ; 1 ; 1 ; 1 ; 2 ; 2 ; 2 ; 2 ; 2 ; 2 ; 3 ; 3 ; 3 ; 3 ; 4 ; 4 ; 5 ; 9$.

\begin{itemize}
\item \textbf{Mode} : la valeur $2$ apparaît 6 fois (effectif maximal) : $Mo = 2$.
\item \textbf{Médiane} : $N = 20$ (pair), valeurs de rangs 10 et 11 : toutes deux valent $2$, donc $Me = 2$.
\item \textbf{Moyenne} : $\bar{x} = \dfrac{0+0+4\times 1 + 6\times 2 + 4\times 3 + 2\times 4 + 5 + 9}{20} = \dfrac{50}{20} = 2{,}5$.
\end{itemize}

\begin{popfigure}
\begin{center}
\begin{tikzpicture}[x=1.05cm,y=0.62cm]
\draw[->,thick,popdark] (-0.3,0) -- (9.6,0) node[below left,font=\small]{buts};
\draw[->,thick,popdark] (0,-0.3) -- (0,7) node[below left,font=\small]{effectif};
\foreach \x in {0,1,2,3,4,5,6,7,8,9} \draw (\x,0.12) -- (\x,-0.12) node[below,font=\scriptsize]{\x};
\foreach \y in {1,2,3,4,5,6} \draw (0.08,\y) -- (-0.08,\y) node[left,font=\scriptsize]{\y};
\draw[very thick,popblue] (0,0) -- (0,2); \fill[popblue] (0,2) circle (2.2);
\draw[very thick,popblue] (1,0) -- (1,4); \fill[popblue] (1,4) circle (2.2);
\draw[very thick,popblue] (2,0) -- (2,6); \fill[popblue] (2,6) circle (2.2);
\draw[very thick,popblue] (3,0) -- (3,4); \fill[popblue] (3,4) circle (2.2);
\draw[very thick,popblue] (4,0) -- (4,2); \fill[popblue] (4,2) circle (2.2);
\draw[very thick,popblue] (5,0) -- (5,1); \fill[popblue] (5,1) circle (2.2);
\draw[very thick,popblue] (9,0) -- (9,1); \fill[popblue] (9,1) circle (2.2);
\draw[dashed,thick,popgreen] (2,0) -- (2,6.6) node[above,font=\small]{mode $=$ médiane $= 2$};
\draw[dashed,thick,poporange] (2.5,0) -- (2.5,6.6) node[above,font=\small,xshift=34pt]{moyenne $= 2{,}5$};
\end{tikzpicture}
\end{center}
\legende{Diagramme à bâtons de la série des buts. Le bâton le plus haut donne le mode ($2$) ; la médiane ($2$) partage l'effectif en deux ; la moyenne ($2{,}5$) est décalée vers la droite à cause du match à $9$ buts.}
\end{popfigure}

\begin{aretenir}[Mode, moyenne, médiane : quel résumé choisir ?]
\begin{itemize}
\item Le \cle{mode} répond à : « quelle est la valeur la plus fréquente ? » — utile pour un stock, une pointure, une note dominante.
\item La \cle{moyenne} répond à : « que vaut chaque part si l'on partage le total équitablement ? » — mais elle est \textbf{sensible aux valeurs extrêmes} : une seule valeur très grande la déforme.
\item La \cle{médiane} répond à : « quelle est la valeur centrale ? » — elle est \textbf{robuste} : les valeurs extrêmes ne la modifient pas (ou très peu).
\end{itemize}
Quand la série contient des valeurs exceptionnelles (salaire du directeur, match à 9 buts), la médiane est souvent un meilleur résumé que la moyenne.
\end{aretenir}

\begin{savaistu}[Le milliardaire du village]
Imaginons un village de 100 habitants où chacun gagne \num{50000} FCFA par mois. Le revenu moyen est de \num{50000} FCFA, et la médiane aussi. Un jour, un milliardaire s'installe dans le village avec un revenu mensuel de 100 millions de FCFA. La moyenne grimpe à environ $\frac{100\times \num{50000} + \num{100000000}}{101} \approx \num{1040000}$ FCFA : on pourrait titrer « le revenu moyen du village dépasse le million ! » alors qu'aucun villageois n'a vu sa vie changer. La médiane, elle, reste sagement à \num{50000} FCFA. C'est pourquoi les instituts de statistique (comme l'INS au Cameroun) publient souvent le revenu \emph{médian} en plus du revenu moyen : la moyenne parle du total des richesses, la médiane parle de la réalité vécue par la moitié de la population.
\end{savaistu}

\begin{exercice}[À toi de jouer]
Les tailles (en cm) de 40 élèves d'une Terminale A de Garoua sont regroupées ainsi :
\begin{center}
\begin{tabular}{|l|c|c|c|c|c|}
\hline
\rowcolor{popblueL} Classe & $[150;160[$ & $[160;165[$ & $[165;170[$ & $[170;180[$ & $[180;190[$ \\
\hline
Effectif & 6 & 10 & 12 & 9 & 3 \\
\hline
\end{tabular}
\end{center}
\begin{enumerate}
\item Déterminer la classe modale et la classe médiane.
\item Calculer la taille moyenne (on construira la colonne des centres).
\item Le directeur des études affirme : « la moitié des élèves mesurent moins de $167$ cm ». Cette affirmation est-elle compatible avec la classe médiane ?
\end{enumerate}
\end{exercice}

\begin{corrige}[Exercice]
\textbf{1.} L'effectif maximal est $12$ : la classe modale est $[165;170[$. Effectifs cumulés croissants : $6$ ; $16$ ; $28$ ; $37$ ; $40$. On a $\frac{N}{2} = 20$ ; le premier cumul atteignant $20$ est $28$ : la classe médiane est aussi $[165;170[$.

\textbf{2.} Centres : $155$ ; $162{,}5$ ; $167{,}5$ ; $175$ ; $185$. Alors :
\[ \bar{x} = \frac{6\times 155 + 10\times 162{,}5 + 12\times 167{,}5 + 9\times 175 + 3\times 185}{40} = \frac{930 + \num{1625} + \num{2010} + \num{1575} + 555}{40} = \frac{\num{6695}}{40} \approx 167{,}4 \text{ cm}. \]

\textbf{3.} La médiane appartient à la classe $[165;170[$ : une lecture graphique sur le polygone des FCC donnerait une valeur voisine de $167$ cm (le cumul atteint $40\,\%$ à la borne $165$ cm et $70\,\%$ à la borne $170$ cm, et $50\,\%$ se situe entre les deux). L'affirmation du directeur est donc \textbf{compatible} avec les données, mais seule une lecture graphique précise permettrait de la confirmer.
\end{corrige}

\coursec{Quartiles et déciles}

Dans la section précédente, nous avons appris à résumer une série statistique par un petit nombre de valeurs : le mode, la moyenne et surtout la médiane, qui partage la série ordonnée en deux moitiés. Mais deux séries peuvent avoir la même médiane et être pourtant très différentes : dans un lycée de Yaoundé, une classe où toutes les notes sont proches de 10 et une classe où la moitié des élèves a 2 et l'autre moitié 18 peuvent avoir la même médiane 10. Pour distinguer ces situations, on raffine le découpage : au lieu de couper la série en deux, on la coupe en \emph{quatre} (les \cle{quartiles}) ou en \emph{dix} (les \cle{déciles}). C'est l'objet de cette section.

\courssub{Les quartiles : partager la série en quatre parts égales}

\textbf{Intuition.}\par
Imagine les 40 candidats d'un concours classés par ordre de mérite. La médiane sépare la moitié admise « dans la bonne moitié » de l'autre. Mais le jury veut plus de détails : repérer le quart des meilleurs candidats (pour une bourse d'excellence) et le quart des plus faibles (pour une session de rattrapage). Il lui faut donc deux nouvelles « frontières » : une valeur qui laisse 25\,\% des candidats en dessous, et une valeur qui laisse 75\,\% des candidats en dessous. Ce sont les quartiles $Q_1$ et $Q_3$.

\begin{definition}[Quartiles]
Soit une série statistique de $N$ valeurs, rangées par ordre croissant.
\begin{itemize}
\item Le \cle{premier quartile} $Q_1$ est la plus petite valeur de la série telle qu'au moins $25\,\%$ des données lui sont inférieures ou égales.
\item Le \cle{deuxième quartile} $Q_2$ n'est autre que la \cle{médiane} $Me$ : au moins $50\,\%$ des données lui sont inférieures ou égales.
\item Le \cle{troisième quartile} $Q_3$ est la plus petite valeur de la série telle qu'au moins $75\,\%$ des données lui sont inférieures ou égales.
\end{itemize}
Ainsi $Q_1$, $Me$ et $Q_3$ partagent la série ordonnée en quatre parts contenant chacune environ $25\,\%$ des données.
\end{definition}

\begin{aretenir}[Lecture concrète des quartiles]
\begin{itemize}
\item $Q_1$ : au moins un quart des valeurs sont inférieures ou égales à $Q_1$ (et donc au plus trois quarts lui sont supérieures).
\item $Q_3$ : au moins trois quarts des valeurs sont inférieures ou égales à $Q_3$ (et donc au plus un quart lui sont supérieures).
\item La moitié centrale des données, environ $50\,\%$, est comprise entre $Q_1$ et $Q_3$.
\end{itemize}
\end{aretenir}

\courssub{Détermination des quartiles dans une série discrète : la méthode des rangs}

Pour une série discrète ordonnée, le programme de Terminale A retient une règle simple et sans ambiguïté : on repère le \emph{rang} du quartile dans la liste ordonnée. La définition « plus petite valeur telle qu'au moins 25\,\%... » conduit à prendre le \emph{plafond} de $N/4$ (c'est-à-dire l'entier supérieur si $N/4$ n'est pas entier, sinon $N/4$ lui-même).

\begin{methode}[Déterminer $Q_1$ et $Q_3$ par les rangs]
Soit une série de $N$ valeurs rangées par ordre croissant : $x_1 \le x_2 \le \dots \le x_N$.
\begin{enumerate}
\item \textbf{Écrire la série dans l'ordre croissant.} C'est l'étape que l'on oublie le plus souvent : les quartiles n'ont de sens que sur une série ordonnée.
\item \textbf{Calculer le rang de $Q_1$.} On calcule $\dfrac{N}{4}$. On arrondit \emph{à l'entier supérieur} (plafond) ; $Q_1$ est alors la valeur de ce rang.
\item \textbf{Calculer le rang de $Q_3$.} On calcule $\dfrac{3N}{4}$, on arrondit de la même façon à l'entier supérieur, et $Q_3$ est la valeur de ce rang.
\item \textbf{Conclure par une phrase d'interprétation} : « au moins 25\,\% des valeurs sont inférieures ou égales à $Q_1$ », etc.
\end{enumerate}
\end{methode}

\begin{exemplebox}[Petit exemple de mise en route]
Voici les notes d'un devoir de 11 élèves, déjà rangées :
\[ 4 \;;\; 6 \;;\; 7 \;;\; 8 \;;\; 9 \;;\; 10 \;;\; 11 \;;\; 12 \;;\; 13 \;;\; 15 \;;\; 17. \]
On a $N = 11$.
\begin{itemize}
\item Rang de $Q_1$ : $\dfrac{11}{4} = 2,75$, arrondi à $3$. Donc $Q_1 = x_3 = 7$.
\item Médiane : rang $\dfrac{11+1}{2} = 6$, donc $Me = x_6 = 10$.
\item Rang de $Q_3$ : $\dfrac{3 \times 11}{4} = 8,25$, arrondi à $9$. Donc $Q_3 = x_9 = 13$.
\end{itemize}
Interprétation : au moins 25\,\% des élèves ont une note inférieure ou égale à 7, et au moins 75\,\% ont une note inférieure ou égale à 13. La moitié des notes se situe entre 7 et 13.
\end{exemplebox}

\courssub{Les déciles : partager la série en dix parts égales}

Pour un découpage encore plus fin, on partage la série ordonnée en \emph{dix} parts d'environ $10\,\%$ chacune. Les neuf « frontières » ainsi obtenues sont les déciles $D_1, D_2, \dots, D_9$.

\begin{definition}[Déciles]
Soit une série statistique de $N$ valeurs rangées par ordre croissant. Pour tout entier $k$ compris entre 1 et 9, le \cle{décile} $D_k$ est la plus petite valeur de la série telle qu'au moins $10 \times k\,\%$ des données lui sont inférieures ou égales.
En particulier :
\begin{itemize}
\item $D_1$ laisse au moins $10\,\%$ des données en dessous ou égales ;
\item $D_5$ laisse au moins $50\,\%$ des données en dessous ou égales : c'est donc la \cle{médiane}, $D_5 = Me = Q_2$ ;
\item $D_9$ laisse au moins $90\,\%$ des données en dessous ou égales.
\end{itemize}
\end{definition}

\begin{methode}[Déterminer un décile par le rang]
Pour une série discrète ordonnée de $N$ valeurs :
\begin{enumerate}
\item Calculer le rang $\dfrac{k \times N}{10}$.
\item Arrondir à l'entier supérieur (plafond).
\item $D_k$ est la valeur de ce rang dans la série ordonnée.
\end{enumerate}
Par exemple, pour $N = 30$ : $D_1$ est la valeur de rang $\frac{30}{10} = 3$, $D_3$ celle de rang $\frac{3 \times 30}{10} = 9$, et $D_9$ celle de rang $\frac{9 \times 30}{10} = 27$.
\end{methode}

La frise ci-dessous visualise ce partage en dix parts : chaque zone colorée contient environ 10\,\% des données, et l'on y lit la position des quartiles (qui coïncident avec certains déciles « intermédiaires » : $Q_1$ est entre $D_2$ et $D_3$, $Me = D_5$, $Q_3$ est entre $D_7$ et $D_8$).

\begin{popfigure}
\begin{center}
\begin{tikzpicture}[x=1.05cm]
% Frise : 10 segments
\foreach \i in {0,...,9} {
  \fill[popblue!\the\numexpr 15+7*\i\relax] (\i,0) rectangle (\i+1,0.9);
  \draw[white,thick] (\i,0) rectangle (\i+1,0.9);
}
\draw[popdark,thick] (0,0) rectangle (10,0.9);
% Étiquettes des déciles
\foreach \i/\nom in {1/D_1,2/D_2,3/D_3,4/D_4,5/D_5,6/D_6,7/D_7,8/D_8,9/D_9} {
  \draw[popdark] (\i,0) -- (\i,-0.18);
  \node[below,popdark,font=\small] at (\i,-0.18) {$\nom$};
}
% Pourcentages
\foreach \i/\p in {0/0\,\%,1/10\,\%,2/20\,\%,3/30\,\%,4/40\,\%,5/50\,\%,6/60\,\%,7/70\,\%,8/80\,\%,9/90\,\%,10/100\,\%} {
  \node[above,popmuted,font=\scriptsize] at (\i,0.9) {\p};
}
% Quartiles
\draw[poporange,very thick] (2.5,0) -- (2.5,1.25);
\node[above,poporange,font=\small\bfseries] at (2.5,1.25) {$Q_1$};
\draw[popink,very thick] (5,0) -- (5,1.25);
\node[above,popink,font=\small\bfseries] at (5,1.25) {$Me = Q_2 = D_5$};
\draw[poporange,very thick] (7.5,0) -- (7.5,1.25);
\node[above,poporange,font=\small\bfseries] at (7.5,1.25) {$Q_3$};
% Zones (remplacement des accolades par des traits et texte)
\draw[popgreen,thick] (0,-0.9) -- (2.5,-0.9);
\draw[popgreen,thick] (0,-0.95) -- (0,-0.85);
\draw[popgreen,thick] (2.5,-0.95) -- (2.5,-0.85);
\node[below,popgreen,font=\small] at (1.25,-1.1) {25\,\% des données};
\draw[popgreen,thick] (2.5,-0.9) -- (7.5,-0.9);
\draw[popgreen,thick] (2.5,-0.95) -- (2.5,-0.85);
\draw[popgreen,thick] (7.5,-0.95) -- (7.5,-0.85);
\node[below,popgreen,font=\small] at (5,-1.1) {50\,\% des données (zone centrale)};
\draw[popgreen,thick] (7.5,-0.9) -- (10,-0.9);
\draw[popgreen,thick] (7.5,-0.95) -- (7.5,-0.85);
\draw[popgreen,thick] (10,-0.95) -- (10,-0.85);
\node[below,popgreen,font=\small] at (8.75,-1.1) {25\,\% des données};
\end{tikzpicture}
\end{center}
\legende{Partage d'une série ordonnée par les déciles (dix parts de 10\,\%) et position des quartiles.}
\end{popfigure}

\courssub{Exemple chiffré complet : les tailles de 30 élèves}

\begin{exoresolu}[Tailles des élèves d'une classe de Terminale A]
On a mesuré les tailles (en cm) des 30 élèves d'une classe de Terminale A d'un lycée de Douala. Voici les valeurs rangées par ordre croissant :
\begin{center}
\begin{tabular}{|c|cccccccccc|}
\hline
\rowcolor{popblueL} Rang & 1 & 2 & 3 & 4 & 5 & 6 & 7 & 8 & 9 & 10 \\
\hline
Taille & 152 & 154 & 155 & 156 & 157 & 158 & 159 & 160 & 160 & 161 \\
\hline
\rowcolor{popblueL} Rang & 11 & 12 & 13 & 14 & 15 & 16 & 17 & 18 & 19 & 20 \\
\hline
Taille & 162 & 163 & 163 & 164 & 165 & 165 & 166 & 167 & 168 & 168 \\
\hline
\rowcolor{popblueL} Rang & 21 & 22 & 23 & 24 & 25 & 26 & 27 & 28 & 29 & 30 \\
\hline
Taille & 169 & 170 & 170 & 171 & 172 & 173 & 174 & 175 & 177 & 180 \\
\hline
\end{tabular}
\end{center}
Déterminer $Q_1$, $Me$, $Q_3$ ainsi que $D_1$ et $D_9$, puis interpréter les résultats.
\tcblower
\textbf{Solution détaillée.} La série est déjà ordonnée et $N = 30$.

\textbf{Premier quartile.} $\dfrac{N}{4} = \dfrac{30}{4} = 7,5$, arrondi à $8$. Donc $Q_1 = x_8 = 160$ cm.

\textbf{Médiane.} $N = 30$ est pair : la médiane est la demi-somme des valeurs de rangs 15 et 16 :
\[ Me = \dfrac{x_{15} + x_{16}}{2} = \dfrac{165 + 165}{2} = 165 \text{ cm}. \]

\textbf{Troisième quartile.} $\dfrac{3N}{4} = \dfrac{90}{4} = 22,5$, arrondi à $23$. Donc $Q_3 = x_{23} = 170$ cm.

\textbf{Déciles $D_1$ et $D_9$.} $\dfrac{30}{10} = 3$, donc $D_1 = x_3 = 155$ cm ; $\dfrac{9 \times 30}{10} = 27$, donc $D_9 = x_{27} = 174$ cm.

\textbf{Interprétations.}
\begin{itemize}
\item Au moins 25\,\% des élèves mesurent moins de 160 cm (ou exactement 160 cm).
\item Au moins la moitié des élèves mesurent 165 cm ou moins.
\item Au moins 75\,\% des élèves mesurent 170 cm ou moins ; autrement dit, environ un quart des élèves mesurent plus de 170 cm.
\item La moitié des élèves de la classe a une taille comprise entre 160 cm et 170 cm.
\item Environ 10\,\% des élèves mesurent 155 cm ou moins, et environ 10\,\% mesurent plus de 174 cm.
\end{itemize}
\end{exoresolu}

\courssub{Lecture graphique sur le polygone des fréquences cumulées croissantes}

Lorsque la série est \emph{groupée en classes}, on ne connaît plus les valeurs individuelles : la méthode des rangs ne s'applique plus. Au programme de Terminale A, on détermine alors les quartiles et les déciles \emph{graphiquement}, grâce au \cle{polygone des fréquences cumulées croissantes} (FCC).

Rappelons le principe : pour chaque classe $[a\,;\,b[$, on porte en abscisse la borne supérieure $b$ et en ordonnée la fréquence cumulée croissante (comprise entre 0 et 1). On relie les points par des segments. Le polygone obtenu est croissant : à une fréquence cumulée donnée sur l'axe des ordonnées correspond une unique valeur sur l'axe des abscisses.

\begin{methode}[Lire quartiles et déciles sur le polygone des FCC]
\begin{enumerate}
\item Construire le tableau des fréquences cumulées croissantes et tracer le polygone des FCC (bornes supérieures des classes en abscisse, FCC en ordonnée, de 0 à 1).
\item Pour $Q_1$ : tracer la droite horizontale d'ordonnée $0,25$. Elle coupe le polygone en un point ; l'abscisse de ce point est $Q_1$.
\item Pour la médiane : même lecture avec l'ordonnée $0,5$.
\item Pour $Q_3$ : même lecture avec l'ordonnée $0,75$.
\item Pour un décile $D_k$ : même lecture avec l'ordonnée $\dfrac{k}{10}$ (par exemple $0,3$ pour $D_3$, $0,9$ pour $D_9$).
\item Conclure en précisant qu'il s'agit de valeurs \emph{lues graphiquement}, donc approchées.
\end{enumerate}
\end{methode}

\begin{exemplebox}[Chiffre d'affaires de boutiques]
Une étude porte sur le chiffre d'affaires mensuel (en centaines de milliers de FCFA) de 60 boutiques du marché central d'une ville. Le tableau des FCC est le suivant :
\begin{center}
\begin{tabularx}{\linewidth}{|l|X|X|X|X|X|}
\hline
\rowcolor{popblueL} Classe & $[2\,;\,6[$ & $[6\,;\,10[$ & $[10\,;\,14[$ & $[14\,;\,18[$ & $[18\,;\,22[$ \\
\hline
Effectif & 6 & 15 & 21 & 12 & 6 \\
\hline
FCC & 0,10 & 0,35 & 0,70 & 0,90 & 1 \\
\hline
\end{tabularx}
\end{center}
Le polygone des FCC est tracé ci-dessous. On lit :
\begin{itemize}
\item $Q_1 \approx 8$ (la droite $y = 0,25$ coupe le polygone au point d'abscisse environ 8) ;
\item $Me \approx 11,4$ (lecture à $y = 0,5$) ;
\item $Q_3 \approx 15$ (lecture à $y = 0,75$).
\end{itemize}
Interprétation : environ un quart des boutiques réalisent un chiffre d'affaires mensuel inférieur à \num{800000} FCFA, la moitié en réalisent moins d'environ \num{1140000} FCFA, et environ un quart dépassent \num{1500000} FCFA.
\end{exemplebox}

\begin{popfigure}
\begin{center}
\begin{tikzpicture}
\begin{axis}[
  width=0.85\linewidth, height=7cm,
  xlabel={Chiffre d'affaires (centaines de milliers de FCFA)},
  ylabel={Fréquences cumulées croissantes},
  xmin=0, xmax=23, ymin=0, ymax=1.05,
  xtick={2,6,10,14,18,22},
  ytick={0,0.25,0.5,0.75,1},
  grid=both, grid style={popline!40},
  axis line style={popdark},
]
% Polygone des FCC
\addplot[popcurve,mark=*,mark options={fill=popblue}] coordinates {(2,0) (6,0.10) (10,0.35) (14,0.70) (18,0.90) (22,1)};
% Lecture Q1
\addplot[poporange,dashed,thick] coordinates {(0,0.25) (8,0.25)};
\addplot[poporange,dashed,thick] coordinates {(8,0) (8,0.25)};
\node[poporange,below,font=\small] at (axis cs:8,0) {$Q_1 \approx 8$};
% Lecture Me
\addplot[popink,dashed,thick] coordinates {(0,0.5) (11.43,0.5)};
\addplot[popink,dashed,thick] coordinates {(11.43,0) (11.43,0.5)};
\node[popink,below,font=\small] at (axis cs:11.43,0) {$Me \approx 11{,}4$};
% Lecture Q3
\addplot[popgreen,dashed,thick] coordinates {(0,0.75) (15,0.75)};
\addplot[popgreen,dashed,thick] coordinates {(15,0) (15,0.75)};
\node[popgreen,below,font=\small] at (axis cs:15,0) {$Q_3 \approx 15$};
\end{axis}
\end{tikzpicture}
\end{center}
\legende{Polygone des fréquences cumulées croissantes et lectures graphiques de $Q_1$, $Me$ et $Q_3$ aux ordonnées $0,25$, $0,5$ et $0,75$.}
\end{popfigure}

\courssub{L'écart interquartile : une mesure simple de dispersion}

Les quartiles ne servent pas qu'à décrire la position : ils permettent aussi de mesurer l'\emph{étalement} de la série.

\begin{definition}[Écart interquartile]
L'\cle{écart interquartile} d'une série statistique est le nombre
\[ Q_3 - Q_1. \]
L'intervalle $[Q_1\,;\,Q_3]$ est appelé \cle{intervalle interquartile} : il contient environ $50\,\%$ des données, les valeurs « centrales » de la série.
\end{definition}

Plus l'écart interquartile est petit, plus les valeurs centrales sont concentrées autour de la médiane ; plus il est grand, plus la série est dispersée. Dans l'exemple des tailles des 30 élèves, $Q_3 - Q_1 = 170 - 160 = 10$ cm : la moitié des élèves se tient dans une fourchette de seulement 10 cm, la classe est assez homogène en taille.

\begin{savaistu}[Pourquoi préférer parfois l'écart interquartile à l'étendue ?]
L'étendue (maximum moins minimum) est très sensible aux valeurs extrêmes : un seul élève de 1,95 m dans une classe la fait exploser. L'écart interquartile, lui, ignore le quart des plus petites valeurs et le quart des plus grandes : il est dit \emph{robuste}. C'est pourquoi les instituts de statistique — comme l'INS du Cameroun dans ses enquêtes sur les revenus des ménages — utilisent massivement les quartiles et les déciles : le rapport $\frac{D_9}{D_1}$, par exemple, est un indicateur classique des inégalités de revenu.
\end{savaistu}

\courssub{Pièges fréquents et erreurs à éviter}

\begin{attention}[L'interpolation linéaire est hors programme en Terminale A]
Dans certains ouvrages, on calcule $Q_1$ ou la médiane d'une série groupée par \emph{interpolation linéaire} à l'intérieur de la classe contenant le quartile. \textbf{Cette méthode n'est pas au programme de la série A.} Au Baccalauréat A :
\begin{itemize}
\item pour une série \emph{discrète}, on détermine les quartiles et déciles \textbf{par les rangs} dans la série ordonnée ;
\item pour une série \emph{groupée en classes}, on les détermine \textbf{par lecture graphique} sur le polygone des fréquences cumulées croissantes, en traçant les lignes horizontales à $0,25$, $0,5$, $0,75$ (ou $\frac{k}{10}$ pour les déciles).
\end{itemize}
Toute valeur lue graphiquement doit être présentée comme approchée (on écrit $Q_1 \approx \dots$) et les traits de lecture doivent figurer sur le graphique : ils justifient ta réponse devant le correcteur.
\end{attention}

\begin{attention}[Autres erreurs classiques]
\begin{itemize}
\item \textbf{Oublier d'ordonner la série.} Les quartiles se déterminent sur la série rangée par ordre croissant ; sur une série brute, le rang ne signifie rien.
\item \textbf{Confondre le rang et la valeur.} Si le rang de $Q_1$ est 8, il faut donner la \emph{valeur} $x_8$, pas le nombre 8.
\item \textbf{Croire que $Q_1$ est « le quart de l'effectif ».} $Q_1$ est une \emph{valeur du caractère} (une taille, une note, un montant en FCFA), jamais un effectif ni un pourcentage.
\item \textbf{Oublier que $Q_2 = D_5 = Me$.} Ces trois notations désignent le même partage en deux moitiés ; les reconnaître dans un énoncé fait gagner un temps précieux.
\item \textbf{Arrondir le rang à l'entier inférieur.} La définition exige « au moins 25\,\% » : on arrondit donc toujours le rang \emph{à l'entier supérieur} (plafond) lorsqu'il n'est pas entier.
\end{itemize}
\end{attention}

\begin{exercice}[À toi de jouer]
Les dépenses hebdomadaires (en milliers de FCFA) de 20 ménages d'un quartier de Bafoussam, rangées par ordre croissant, sont :
\[ 8 \;;\; 9 \;;\; 10 \;;\; 11 \;;\; 12 \;;\; 12 \;;\; 13 \;;\; 14 \;;\; 15 \;;\; 15 \;;\; 16 \;;\; 17 \;;\; 18 \;;\; 19 \;;\; 20 \;;\; 21 \;;\; 23 \;;\; 25 \;;\; 28 \;;\; 32. \]
\begin{enumerate}
\item Déterminer $Q_1$, $Me$ et $Q_3$.
\item Déterminer $D_1$ et $D_9$.
\item Calculer l'écart interquartile et interpréter le résultat par une phrase.
\end{enumerate}
\end{exercice}

\begin{corrige}[Exercice : dépenses des ménages]
On a $N = 20$ et la série est déjà ordonnée.
\begin{enumerate}
\item Rang de $Q_1$ : $\dfrac{20}{4} = 5$, donc $Q_1 = x_5 = 12$ milliers de FCFA. Médiane : demi-somme des valeurs de rangs 10 et 11 : $Me = \dfrac{15 + 16}{2} = 15,5$ milliers de FCFA. Rang de $Q_3$ : $\dfrac{3 \times 20}{4} = 15$, donc $Q_3 = x_{15} = 20$ milliers de FCFA.
\item $D_1$ : rang $\dfrac{20}{10} = 2$, donc $D_1 = x_2 = 9$. $D_9$ : rang $\dfrac{9 \times 20}{10} = 18$, donc $D_9 = x_{18} = 25$ milliers de FCFA.
\item $Q_3 - Q_1 = 20 - 12 = 8$ milliers de FCFA : la moitié des ménages dépensent entre \num{12000} et \num{20000} FCFA par semaine, ce qui traduit une dispersion modérée autour de la médiane.
\end{enumerate}
\end{corrige}

\begin{aretenir}[L'essentiel de la section]
\begin{itemize}
\item Les \cle{quartiles} $Q_1$, $Q_2 = Me$, $Q_3$ partagent la série ordonnée en quatre parts de 25\,\% ; les \cle{déciles} $D_1, \dots, D_9$ la partagent en dix parts de 10\,\%, avec $D_5 = Me$.
\item Série discrète : $Q_1$ est la valeur de rang $\frac{N}{4}$ (arrondi au supérieur), $Q_3$ celle de rang $\frac{3N}{4}$, $D_k$ celle de rang $\frac{kN}{10}$.
\item Série groupée : lecture \emph{graphique} sur le polygone des FCC aux ordonnées $0,25$ ; $0,5$ ; $0,75$ (et $\frac{k}{10}$ pour les déciles). Pas d'interpolation linéaire au Bac A.
\item L'\cle{écart interquartile} $Q_3 - Q_1$ mesure la dispersion de la moitié centrale des données : c'est un complément précieux de la médiane.
\end{itemize}
\end{aretenir}

\coursec{Séries à deux caractères : nuage de points et point moyen}

Dans les sections précédentes, nous avons appris à résumer une série statistique à \emph{un seul} caractère : la médiane, les quartiles et les déciles nous donnent une photographie précise de la dispersion d'une seule variable (les notes d'une classe, les tailles des élèves, les prix d'un produit). Mais dans la vie réelle, on s'intéresse très souvent à \emph{deux grandeurs à la fois} et surtout au \emph{lien} qui peut exister entre elles. Le commerçant du marché Mokolo se demande : « quand j'augmente le prix du kilogramme de tomates, de combien diminue la quantité vendue ? » Le médecin se demande : « le poids d'un nouveau-né dépend-il de son âge gestationnel ? » L'opérateur téléphonique se demande : « la consommation de données augmente-t-elle avec le revenu des abonnés ? »

Pour répondre à ces questions, il faut étudier les deux caractères \emph{simultanément}. C'est l'objet de cette section : construire un outil visuel, le \cle{nuage de points}, puis le résumer par un point remarquable, le \cle{point moyen}, qui servira de fondation à la section suivante (l'ajustement affine par la méthode de Mayer).

\courssub{Série statistique double : les couples $(x_i\,;\,y_i)$}

\begin{definition}[Série statistique double]
On dit qu'on a une \cle{série statistique double} (ou série à deux caractères) lorsque, sur une même population de $n$ individus, on relève \emph{simultanément} deux caractères quantitatifs $X$ et $Y$. La série est alors la liste des $n$ couples :
\[
(x_1\,;\,y_1),\ (x_2\,;\,y_2),\ \dots,\ (x_n\,;\,y_n)
\]
où $x_i$ et $y_i$ sont les valeurs des deux caractères relevées sur le $i$-ème individu.
\end{definition}

\begin{attention}[Les deux valeurs d'un couple sont inséparables]
Dans une série double, $x_i$ et $y_i$ correspondent au \emph{même individu} (ou à la même observation). On ne peut pas mélanger les lignes du tableau ! Si le jour $n^{\circ}3$ le prix était de \SI{200}{FCFA} et la quantité vendue de \SI{70}{kg}, le couple $(200\,;\,70)$ forme un tout : permuter les $y_i$ entre eux détruirait complètement l'information sur le lien entre les deux caractères.
\end{attention}

\begin{exemplebox}[Prix et quantité vendue au marché Mokolo]
Une vendeuse de tomates au marché Mokolo (Yaoundé) relève, pendant 8 jours de marché, le prix de vente du kilogramme de tomates et la quantité totale vendue dans la journée :

\begin{center}
\begin{tabularx}{\linewidth}{|l|X|X|X|X|X|X|X|X|}
\hline
\rowcolor{popblueL} Jour $i$ & 1 & 2 & 3 & 4 & 5 & 6 & 7 & 8 \\
\hline
Prix $x_i$ (FCFA/kg) & 100 & 150 & 200 & 250 & 300 & 350 & 400 & 450 \\
\hline
Quantité $y_i$ (kg) & 80 & 75 & 70 & 60 & 55 & 45 & 40 & 30 \\
\hline
\end{tabularx}
\end{center}

On obtient une série double de 8 couples : $(100\,;\,80)$, $(150\,;\,75)$, \dots, $(450\,;\,30)$. La question naturelle est : \emph{comment la quantité vendue évolue-t-elle quand le prix augmente ?}
\end{exemplebox}

\courssub{Le nuage de points : voir le lien d'un coup d'œil}

L'idée est d'une simplicité géniale : chaque couple $(x_i\,;\,y_i)$ n'est rien d'autre qu'un couple de coordonnées. On peut donc le représenter par un \emph{point} dans un repère du plan. L'ensemble de tous ces points forme un « nuage » dont la \emph{forme} révèle le lien entre les deux caractères.

\begin{definition}[Nuage de points]
Dans un repère orthogonal $(O\,;\,\vec{\imath},\vec{\jmath})$, le \cle{nuage de points} associé à la série double $(x_i\,;\,y_i)_{1\le i\le n}$ est l'ensemble des $n$ points $M_i$ de coordonnées $(x_i\,;\,y_i)$.
\end{definition}

\begin{methode}[Construire un nuage de points lisible]
\begin{enumerate}
\item Repérer les valeurs \emph{minimales et maximales} de $x$ et de $y$ (ici $x$ varie de 100 à 450, $y$ de 30 à 80).
\item Choisir les \emph{unités} sur chaque axe de sorte que le nuage occupe bien toute la feuille : les deux axes n'ont pas besoin d'avoir la même échelle (ce ne sont pas des grandeurs de même nature).
\item \emph{Décaler l'origine} si nécessaire : si toutes les valeurs de $x$ sont comprises entre 100 et 450, il est inutile de faire commencer l'axe à 0 ; on peut graduer à partir de 100 (on le signale parfois par un trait de coupure).
\item Placer chaque point $M_i(x_i\,;\,y_i)$ par une croix ou un point bien visible, sans les relier.
\end{enumerate}
\end{methode}

\begin{attention}[Erreur fréquente : des graduations qui écrasent le nuage]
L'erreur la plus fréquente au Baccalauréat est de mal choisir les graduations. Si tu gradues l'axe des ordonnées de 0 à 80 alors que toutes les valeurs de $y$ sont entre 30 et 80, ton nuage sera tassé dans la moitié supérieure du graphique et sa forme sera illisible. \textbf{Réflexe systématique} : regarder les valeurs extrêmes de chaque caractère, décaler l'origine de chaque axe juste en dessous du minimum, et choisir un pas de graduation régulier et commode (10, 25, 50\dots). Un nuage bien construit occupe l'essentiel de la zone graphique.
\end{attention}

\begin{popfigure}
\begin{center}
\begin{tikzpicture}
\begin{axis}[
    width=0.9\linewidth, height=7.5cm,
    xlabel={Prix $x$ (FCFA/kg)}, ylabel={Quantité $y$ (kg)},
    xmin=50, xmax=500, ymin=20, ymax=90,
    xtick={100,150,200,250,300,350,400,450},
    ytick={30,40,50,60,70,80},
    grid=both, grid style={popline!40},
    legend style={at={(0.02,0.02)}, anchor=south west, font=\small}
]
\addplot[only marks, mark=*, mark size=2.5pt, popblue] coordinates {
    (100,80) (150,75) (200,70) (250,60) (300,55) (350,45) (400,40) (450,30)
};
\addlegendentry{Points du nuage $M_i$}
\addplot[only marks, mark=*, mark size=3.5pt, popink] coordinates {(275,56.875)};
\addlegendentry{Point moyen $G(275\,;\,56{,}875)$}
\node[popink, anchor=south west, font=\small\bfseries] at (axis cs:280,58) {$G$};
\end{axis}
\end{tikzpicture}
\end{center}
\legende{Nuage de points de la série (prix ; quantité vendue) au marché Mokolo. Les points semblent s'aligner le long d'une droite descendante : un ajustement affine paraît pertinent. Le point moyen $G$ est marqué en rouge.}
\end{popfigure}

\courssub{Le point moyen $G$ : le centre du nuage}

Le nuage donne une information visuelle, mais on a besoin d'un point de référence précis, qui jouera pour le nuage le rôle que la moyenne joue pour une série simple. C'est le \cle{point moyen}.

\begin{definition}[Point moyen]
Le \cle{point moyen} d'un nuage de $n$ points $M_i(x_i\,;\,y_i)$ est le point $G$ dont les coordonnées sont les moyennes arithmétiques des coordonnées des points du nuage :
\[
G\left(\bar{x}\,;\,\bar{y}\right) \qquad \text{avec} \qquad \bar{x}=\frac{1}{n}\sum_{i=1}^{n} x_i \quad \text{et} \quad \bar{y}=\frac{1}{n}\sum_{i=1}^{n} y_i .
\]
\end{definition}

Autrement dit : on calcule \emph{séparément} la moyenne des $x_i$ et la moyenne des $y_i$, comme si on avait deux séries simples indépendantes, et ces deux moyennes sont les coordonnées de $G$.

\begin{methode}[Calculer et placer le point moyen $G$]
\begin{enumerate}
\item Calculer la somme $\sum x_i$ de toutes les valeurs du premier caractère, puis diviser par $n$ : c'est $\bar{x}$.
\item Calculer la somme $\sum y_i$ de toutes les valeurs du second caractère, puis diviser par $n$ : c'est $\bar{y}$.
\item Conclure : $G(\bar{x}\,;\,\bar{y})$.
\item Placer $G$ sur le graphique avec une couleur ou un symbole différent des points du nuage, et le nommer explicitement.
\end{enumerate}
\end{methode}

\begin{exemplebox}[Calcul du point moyen pour le marché Mokolo]
Reprenons nos 8 relevés. Calculons les deux sommes :
\begin{align*}
\sum x_i &= 100+150+200+250+300+350+400+450 = 2200 \\
\sum y_i &= 80+75+70+60+55+45+40+30 = 455
\end{align*}
On en déduit :
\[
\bar{x}=\frac{2200}{8}=275 \qquad \text{et} \qquad \bar{y}=\frac{455}{8}=56{,}875.
\]
Le point moyen du nuage est donc $G(275\,;\,56{,}875)$ : c'est le point rouge de la figure précédente. Interprétation concrète : en moyenne, sur ces 8 jours, le prix pratiqué était de \SI{275}{FCFA/kg} et la quantité vendue d'environ \SI{56,9}{kg} par jour.
\end{exemplebox}

\begin{attention}[Pièges à éviter avec le point moyen]
\begin{itemize}
\item Ne confonds pas $\bar{x}$ et $\bar{y}$ : le point moyen s'écrit $G(\bar{x}\,;\,\bar{y})$, dans le même ordre que les couples $(x_i\,;\,y_i)$.
\item Le point moyen n'est \emph{pas}, en général, un point du nuage : ici $G(275\,;\,56{,}875)$ ne correspond à aucun jour observé. C'est un point « synthétique ».
\item Si les valeurs sont données avec des \emph{effectifs} (tableau d'effectifs par couple), il faut utiliser les moyennes pondérées : $\bar{x}=\dfrac{\sum n_i x_i}{\sum n_i}$.
\end{itemize}
\end{attention}

\begin{savaistu}[Le point moyen, un « centre de gravité »]
Le point moyen $G$ est exactement le \emph{centre de gravité} (ou barycentre) du nuage, au sens de la physique : si tu découpes le nuage dans une plaque de carton en attribuant à chaque point la même masse, la plaque resterait en équilibre sur une pointe placée sous $G$. C'est d'ailleurs la même notion que le barycentre de points affectés de coefficients égaux que tu as rencontré en géométrie. Ce n'est pas un hasard si la lettre $G$ est utilisée : $G$ comme « gravité » ! En statistiques, $G$ est le point autour duquel le nuage s'équilibre, et c'est par lui que passera la droite d'ajustement de Mayer que nous construirons dans la prochaine section.
\end{savaistu}

\courssub{Forme du nuage : quand un ajustement affine est-il envisageable ?}

Tout l'enjeu de la représentation graphique est de répondre à la question : \emph{les deux caractères semblent-ils liés par une relation affine} $y = ax+b$ ?

\begin{aretenir}[Lecture de la forme d'un nuage]
\begin{itemize}
\item Si le nuage est \cle{allongé autour d'une droite} (les points semblent presque alignés), on dit qu'il existe une \emph{corrélation affine} entre les deux caractères : un \emph{ajustement affine} est pertinent. Si la droite imaginée « monte », le lien est croissant ; si elle « descend » (comme prix et quantité vendue), le lien est décroissant.
\item Si le nuage est \cle{très dispersé}, sans direction privilégiée (forme « ronde » ou éclatée), il n'y a pas de lien affine apparent : chercher une droite d'ajustement n'a \emph{aucun sens}.
\item Si le nuage épouse une \cle{courbe} (parabole, exponentielle\dots), un lien existe mais il n'est \emph{pas affine} : l'ajustement par une droite est inadapté.
\end{itemize}
\end{aretenir}

\begin{popfigure}
\begin{center}
\begin{tikzpicture}
\begin{axis}[
    name=gauche,
    width=0.48\linewidth, height=6cm,
    title={\small Nuage allongé : ajustement pertinent},
    title style={font=\small},
    xmin=0, xmax=10, ymin=0, ymax=10,
    xtick=\empty, ytick=\empty,
    grid=both, grid style={popline!30}
]
\addplot[only marks, mark=*, mark size=2pt, popblue] coordinates {
    (1,1.6) (2,2.7) (3,3.2) (4,4.4) (5,5.1) (6,5.8) (7,7.1) (8,7.6) (9,8.8)
};
\addplot[popcurve, domain=0.5:9.5, samples=2, poporange, dashed] {0.9*x+0.6};
\end{axis}
\begin{axis}[
    at={($(gauche.east)+(1.2cm,0)$)}, anchor=west,
    width=0.48\linewidth, height=6cm,
    title={\small Nuage dispersé : pas d'ajustement affine},
    title style={font=\small},
    xmin=0, xmax=10, ymin=0, ymax=10,
    xtick=\empty, ytick=\empty,
    grid=both, grid style={popline!30}
]
\addplot[only marks, mark=*, mark size=2pt, poppurple] coordinates {
    (1,3) (1.5,7.5) (2.5,2) (3,6) (4,8.5) (4.5,1.5) (5.5,5) (6,9) (7,2.5) (8,6.5) (9,4) (9.5,8)
};
\end{axis}
\end{tikzpicture}
\end{center}
\legende{À gauche, un nuage allongé autour d'une direction : on peut envisager une droite d'ajustement. À droite, un nuage sans structure : aucune droite ne résume le nuage, l'ajustement affine n'est pas pertinent.}
\end{popfigure}

\begin{exoresolu}[Nuage de points et point moyen]
Le tableau suivant donne, pour 6 communes du Cameroun, le taux d'électrification $x_i$ (en \%) et le nombre moyen $y_i$ de téléviseurs pour 10 ménages :
\begin{center}
\begin{tabular}{|l|c|c|c|c|c|c|}
\hline
\rowcolor{popblueL} Commune & A & B & C & D & E & F \\
\hline
$x_i$ (en \%) & 20 & 35 & 50 & 65 & 80 & 90 \\
\hline
$y_i$ & 2 & 3 & 5 & 6 & 8 & 9 \\
\hline
\end{tabular}
\end{center}
\begin{enumerate}
\item Représenter le nuage de points (on décalera l'origine des axes).
\item Calculer les coordonnées du point moyen $G$ et le placer sur le graphique.
\item Un ajustement affine de ce nuage semble-t-il pertinent ? Justifier.
\end{enumerate}
\tcblower
\textbf{Solution.}
\begin{enumerate}
\item Les valeurs de $x$ vont de 20 à 90 : on gradue l'axe des abscisses à partir de 20 (par exemple 1 cm pour 10 \%). Les valeurs de $y$ vont de 2 à 9 : on gradue l'axe des ordonnées à partir de 0 ou de 2 (1 cm pour 1 téléviseur). On place les six points $(20\,;\,2)$, $(35\,;\,3)$, $(50\,;\,5)$, $(65\,;\,6)$, $(80\,;\,8)$, $(90\,;\,9)$ : ils sont presque alignés le long d'une droite montante.
\item Calculons les moyennes :
\begin{align*}
\bar{x} &= \frac{20+35+50+65+80+90}{6} = \frac{340}{6} = \frac{170}{3} \approx 56{,}7 \\
\bar{y} &= \frac{2+3+5+6+8+9}{6} = \frac{33}{6} = 5{,}5
\end{align*}
Donc $G\left(\dfrac{170}{3}\,;\,5{,}5\right)$, soit approximativement $G(56{,}7\,;\,5{,}5)$. On le place sur le graphique avec un symbole distinct (croix rouge, par exemple).
\item Le nuage est nettement allongé autour d'une direction ascendante : les points sont presque alignés. Un ajustement affine est donc \emph{pertinent}. Concrètement : plus le taux d'électrification est élevé, plus l'équipement en téléviseurs est important, et ce lien semble de nature affine.
\end{enumerate}
\end{exoresolu}

\begin{exercice}[À toi de jouer]
Une station-service relève pendant 7 jours le prix du litre d'essence $x_i$ (en FCFA) et le volume vendu $y_i$ (en centaines de litres) :
\begin{center}
\begin{tabular}{|l|c|c|c|c|c|c|c|}
\hline
\rowcolor{popblueL} Jour & 1 & 2 & 3 & 4 & 5 & 6 & 7 \\
\hline
$x_i$ & 630 & 640 & 650 & 660 & 670 & 680 & 690 \\
\hline
$y_i$ & 52 & 50 & 47 & 46 & 43 & 41 & 39 \\
\hline
\end{tabular}
\end{center}
\begin{enumerate}
\item Construire le nuage de points en choisissant judicieusement l'origine et les unités des axes.
\item Déterminer les coordonnées exactes du point moyen $G$ et le placer.
\item La forme du nuage suggère-t-elle un ajustement affine ? De quel sens (croissant ou décroissant) ?
\end{enumerate}
\end{exercice}

\begin{corrige}[Exercice]
\begin{enumerate}
\item $x$ varie de 630 à 690 : on décale l'origine à 630 (1 cm pour 10 FCFA). $y$ varie de 39 à 52 : on décale l'origine à 38 par exemple (1 cm pour 2 centaines de litres). On place les 7 points.
\item $\bar{x}=\dfrac{630+640+\cdots+690}{7}=\dfrac{4620}{7}=660$ et $\bar{y}=\dfrac{52+50+47+46+43+41+39}{7}=\dfrac{318}{7}\approx 45{,}43$. Donc $G(660\,;\,\frac{318}{7})$.
\item Le nuage est allongé autour d'une droite descendante : un ajustement affine est pertinent, avec un coefficient directeur \emph{négatif} (quand le prix augmente, le volume vendu diminue).
\end{enumerate}
\end{corrige}

\begin{aretenir}[L'essentiel de la section]
\begin{itemize}
\item Une \cle{série double} est une liste de couples $(x_i\,;\,y_i)$ relevés sur les mêmes individus ; chaque couple est inséparable.
\item Le \cle{nuage de points} est l'ensemble des points $M_i(x_i\,;\,y_i)$ dans un repère : soigne le choix des unités et décale l'origine si besoin.
\item Le \cle{point moyen} est $G(\bar{x}\,;\,\bar{y})$ avec $\bar{x}=\dfrac{\sum x_i}{n}$ et $\bar{y}=\dfrac{\sum y_i}{n}$ ; c'est le « centre de gravité » du nuage.
\item Un nuage \emph{allongé autour d'une droite} autorise un ajustement affine ; un nuage dispersé ou courbé le rend non pertinent.
\end{itemize}
\end{aretenir}

Maintenant que nous savons représenter un nuage, le résumer par son point moyen et juger de la pertinence d'un ajustement, il reste à construire \emph{effectivement} la droite qui résume le mieux le nuage : c'est l'objet de la section suivante, consacrée à l'ajustement affine par la méthode de Mayer et à son utilisation pour faire des prévisions.

\coursec{Ajustement affine par la méthode de Mayer et prévisions}

Dans la section précédente, nous avons appris à représenter une série statistique à deux caractères par un \cle{nuage de points} et à résumer ce nuage par son \cle{point moyen} $G$. Nous avons constaté que, lorsque les points du nuage semblent s'aligner le long d'une direction privilégiée, on peut résumer toute l'information du nuage par une \emph{droite}. Reste à construire cette droite de manière rigoureuse et reproductible : c'est exactement l'objet de cette dernière section. Une fois cette droite obtenue, elle devient un formidable outil de \cle{prévision} : on pourra estimer une valeur inconnue de $y$ à partir d'une valeur de $x$.

\courssub{L'idée de l'ajustement affine}

\courssub{Pourquoi ajuster un nuage par une droite ?}

Imaginons la gérante d'une boutique d'airtime à Yaoundé qui note chaque mois son chiffre d'affaires. Les montants varient d'un mois à l'autre, mais globalement ils \emph{augmentent}. Si elle reporte ses données sur un graphique, elle obtient un nuage de points qui « monte » de la gauche vers la droite. Elle aimerait répondre à des questions comme : \emph{« Si la tendance se poursuit, quel chiffre d'affaires puis-je espérer dans deux mois ? »}. Pour cela, il lui faut remplacer le nuage par un modèle simple : une droite qui passe « au mieux » au milieu des points.

\begin{definition}[Ajustement affine]
Soit une série statistique à deux caractères $x$ et $y$, représentée par un nuage de $n$ points $M_i(x_i\,;\,y_i)$. Réaliser un \cle{ajustement affine} du nuage, c'est déterminer une droite $\mathcal{D}$ d'équation $y = ax + b$ qui passe « le plus près possible » de l'ensemble des points du nuage. Cette droite est appelée \cle{droite d'ajustement} du nuage.
\end{definition}

\begin{attention}[Quand l'ajustement affine a-t-il un sens ?]
On ne cherche une droite d'ajustement que lorsque le nuage a une \emph{forme allongée}, c'est-à-dire lorsque les points semblent se répartir autour d'une direction rectiligne. Si le nuage est rond, ou s'il dessine une courbe (parabole, exponentielle...), un ajustement affine n'a pas de sens et les prévisions qu'on en tirerait seraient trompeuses.
\end{attention}

Il existe plusieurs méthodes pour construire une droite d'ajustement. Au programme de Terminale A figure la \cle{méthode de Mayer}, simple, rapide et entièrement calculable à la main.

\courssub{La méthode de Mayer}

\courssub{Principe : deux sous-nuages, deux points moyens}

L'idée de Mayer est d'une grande élégance. Plutôt que de chercher une droite passant « au mieux » parmi tous les points, on va :
\begin{itemize}
\item découper le nuage en \textbf{deux sous-nuages} de même effectif (ou d'effectifs aussi proches que possible), en rangeant les points par valeurs \emph{croissantes} de $x$ ;
\item remplacer chaque sous-nuage par son \textbf{point moyen} : $G_1$ pour le premier sous-nuage, $G_2$ pour le second ;
\item tracer la droite $(G_1G_2)$ : c'est la droite de Mayer.
\end{itemize}

\begin{definition}[Droite de Mayer]
Soit un nuage de $n$ points rangés par valeurs croissantes de $x$. On partage le nuage en deux sous-nuages de même effectif (ou d'effectifs différant de $1$ si $n$ est impair) : le premier sous-nuage contient les points d'abscisses les plus petites, le second ceux d'abscisses les plus grandes. On note $G_1(\bar{x}_1\,;\,\bar{y}_1)$ et $G_2(\bar{x}_2\,;\,\bar{y}_2)$ les points moyens respectifs de ces deux sous-nuages. La \cle{droite de Mayer} du nuage est la droite $(G_1G_2)$. Son équation $y = ax + b$ est obtenue par :
\[ a = \frac{\bar{y}_2 - \bar{y}_1}{\bar{x}_2 - \bar{x}_1} \qquad \text{puis} \qquad b = \bar{y}_1 - a\,\bar{x}_1. \]
\end{definition}

\begin{methode}[Méthode de Mayer en 5 étapes]
\begin{enumerate}
\item \textbf{Ordonner} : ranger les points du nuage par valeurs croissantes du caractère $x$.
\item \textbf{Partager} : découper la série en deux sous-nuages de même effectif (si $n$ est pair) ou d'effectifs différant de $1$ (si $n$ est impair).
\item \textbf{Calculer $G_1$ et $G_2$} : pour chaque sous-nuage, calculer la moyenne des abscisses et la moyenne des ordonnées :
\[ G_1\left(\bar{x}_1 = \frac{1}{n_1}\sum x_i \ ;\ \bar{y}_1 = \frac{1}{n_1}\sum y_i\right), \qquad G_2\left(\bar{x}_2 \ ;\ \bar{y}_2\right). \]
\item \textbf{Écrire l'équation} : calculer $a = \dfrac{\bar{y}_2 - \bar{y}_1}{\bar{x}_2 - \bar{x}_1}$, puis $b = \bar{y}_1 - a\bar{x}_1$. La droite de Mayer a pour équation $y = ax + b$.
\item \textbf{Tracer} : placer $G_1$ et $G_2$ sur le graphique et tracer la droite $(G_1G_2)$ à travers le nuage.
\end{enumerate}
\end{methode}

\begin{attention}[L'erreur classique : oublier d'ordonner]
Le piège numéro 1 de la méthode de Mayer est de partager le tableau \emph{tel qu'il est présenté dans l'énoncé}, sans vérifier que les valeurs de $x$ sont rangées dans l'ordre croissant. Le partage se fait toujours \textbf{après avoir rangé les points par abscisses croissantes}. Si les données ne sont pas ordonnées, les deux sous-nuages ne correspondent plus aux « petites » et « grandes » valeurs de $x$, et la droite obtenue est fausse. Prends l'habitude de réécrire le tableau dans l'ordre avant tout calcul.
\end{attention}

\begin{propriete}[La droite de Mayer passe par le point moyen $G$ (admise)]
La droite de Mayer d'un nuage passe toujours par le \cle{point moyen} $G(\bar{x}\,;\,\bar{y})$ du nuage complet. Cette propriété est admise ; elle fournit un excellent moyen de \emph{vérifier} ses calculs : les coordonnées de $G$ doivent satisfaire l'équation $\bar{y} = a\bar{x} + b$.
\end{propriete}

\begin{experience}[Vérification guidée de la propriété]
Reprenons des notations simples : le premier sous-nuage a $n_1$ points, le second $n_2$ points, avec $n = n_1 + n_2$.
\begin{enumerate}
\item Rappelle la définition du point moyen du nuage complet : $\bar{x} = \dfrac{1}{n}\displaystyle\sum_{i=1}^{n} x_i$.
\item Regroupe la somme en deux morceaux : les points du premier sous-nuage, puis ceux du second. Montre que
\[ \bar{x} = \frac{n_1\,\bar{x}_1 + n_2\,\bar{x}_2}{n_1 + n_2} \qquad \text{et de même} \qquad \bar{y} = \frac{n_1\,\bar{y}_1 + n_2\,\bar{y}_2}{n_1 + n_2}. \]
\item Dans le cas fréquent où $n_1 = n_2$, en déduis que $G$ est le \emph{milieu} du segment $[G_1G_2]$.
\item Conclus : $G$ appartient à la droite $(G_1G_2)$, c'est-à-dire à la droite de Mayer.
\end{enumerate}
Cette vérification sera menée concrètement sur l'exemple chiffré ci-dessous.
\end{experience}

\courssub{Exemple complet : la boutique d'airtime}

\begin{exemplebox}[Chiffre d'affaires d'une boutique d'airtime sur 8 mois]
Madame Essomba tient une boutique d'airtime (MTN et Orange) au marché de Mokolo. Elle relève son chiffre d'affaires mensuel $y$ (en dizaines de milliers de FCFA) en fonction du rang $x$ du mois :

\begin{center}
\begin{tabular}{|c|*{8}{c}|}
\hline
\rowcolor{popblueL} Mois $x_i$ & 1 & 2 & 3 & 4 & 5 & 6 & 7 & 8 \\
\hline
CA $y_i$ & 120 & 135 & 128 & 150 & 160 & 155 & 175 & 190 \\
\hline
\end{tabular}
\end{center}

\textbf{Étape 1 — Ordonner.} Les valeurs de $x$ sont déjà croissantes : rien à faire.

\textbf{Étape 2 — Partager.} Il y a $n = 8$ points : deux sous-nuages de $4$ points chacun.
\begin{itemize}
\item Sous-nuage 1 : $(1\,;\,120)$, $(2\,;\,135)$, $(3\,;\,128)$, $(4\,;\,150)$ ;
\item Sous-nuage 2 : $(5\,;\,160)$, $(6\,;\,155)$, $(7\,;\,175)$, $(8\,;\,190)$.
\end{itemize}

\textbf{Étape 3 — Points moyens.}
\begin{align*}
\bar{x}_1 &= \frac{1+2+3+4}{4} = 2{,}5, & \bar{y}_1 &= \frac{120+135+128+150}{4} = \frac{533}{4} = 133{,}25, \\
\bar{x}_2 &= \frac{5+6+7+8}{4} = 6{,}5, & \bar{y}_2 &= \frac{160+155+175+190}{4} = \frac{680}{4} = 170.
\end{align*}
Donc $G_1(2{,}5\,;\,133{,}25)$ et $G_2(6{,}5\,;\,170)$.

\textbf{Étape 4 — Équation de la droite.}
\[ a = \frac{\bar{y}_2 - \bar{y}_1}{\bar{x}_2 - \bar{x}_1} = \frac{170 - 133{,}25}{6{,}5 - 2{,}5} = \frac{36{,}75}{4} = 9{,}1875, \]
\[ b = \bar{y}_1 - a\,\bar{x}_1 = 133{,}25 - 9{,}1875 \times 2{,}5 = 133{,}25 - 22{,}96875 = 110{,}28125. \]
La droite de Mayer a pour équation : $\boxed{y = 9{,}1875\,x + 110{,}28125}$.

\textbf{Vérification avec le point moyen $G$.} Le point moyen du nuage complet est :
\[ \bar{x} = \frac{1+2+\cdots+8}{8} = 4{,}5, \qquad \bar{y} = \frac{533 + 680}{8} = \frac{1213}{8} = 151{,}625. \]
Or $a\bar{x} + b = 9{,}1875 \times 4{,}5 + 110{,}28125 = 41{,}34375 + 110{,}28125 = 151{,}625 = \bar{y}$. La propriété est vérifiée : nos calculs sont cohérents.

\textbf{Étape 5 — Tracé.} On place $G_1$, $G_2$ et on trace la droite $(G_1G_2)$ (voir figure ci-dessous).
\end{exemplebox}

\begin{popfigure}
\begin{tikzpicture}
\begin{axis}[
    width=\linewidth, height=8cm,
    xlabel={Mois $x$}, ylabel={CA $y$ (dizaines de milliers de FCFA)},
    xmin=0, xmax=11, ymin=90, ymax=230,
    grid=both, grid style={popline!40},
    xtick={1,2,3,4,5,6,7,8,9,10},
    legend style={at={(0.02,0.98)}, anchor=north west, font=\small}
]
\addplot[only marks, mark=*, mark size=2.2pt, popblue] coordinates {(1,120) (2,135) (3,128) (4,150) (5,160) (6,155) (7,175) (8,190)};
\addlegendentry{Points du nuage}
\addplot[popcurve, poporange, domain=0.5:10.5, samples=2]{9.1875*x + 110.28125};
\addlegendentry{Droite de Mayer}
\addplot[only marks, mark=square*, mark size=3pt, popgreen] coordinates {(2.5,133.25) (6.5,170)};
\addlegendentry{$G_1$ et $G_2$}
\addplot[only marks, mark=diamond*, mark size=3.2pt, poppurple] coordinates {(4.5,151.625)};
\addlegendentry{Point moyen $G$}
\addplot[dashed, poppink, thick] coordinates {(10,90) (10,202.15625)};
\addplot[dashed, poppink, thick] coordinates {(0,202.15625) (10,202.15625)};
\addplot[only marks, mark=*, mark size=2.5pt, poppink] coordinates {(10,202.15625)};
\node[font=\small, poppink, anchor=west] at (axis cs:10.1,202) {prévision};
\node[font=\footnotesize, popgreen, anchor=south east] at (axis cs:2.5,133.25) {$G_1$};
\node[font=\footnotesize, popgreen, anchor=south east] at (axis cs:6.5,170) {$G_2$};
\end{axis}
\end{tikzpicture}
\legende{Nuage de points, points moyens $G_1$ et $G_2$, point moyen $G$ et droite de Mayer. La prévision pour le mois 10 est matérialisée par les pointillés roses.}
\end{popfigure}

\courssub{Utiliser la droite de Mayer pour prévoir}

Une fois l'équation $y = ax + b$ obtenue, la prévision est un simple calcul de substitution.

\begin{methode}[Prévision à l'aide de la droite d'ajustement]
\begin{enumerate}
\item \textbf{Prévoir $y$ connaissant $x$} : remplacer $x$ par sa valeur dans l'équation $y = ax + b$ et calculer $y$.
\item \textbf{Prévoir $x$ connaissant $y$} (prévision réciproque) : résoudre l'équation $ax + b = y_0$, c'est-à-dire calculer $x = \dfrac{y_0 - b}{a}$.
\item \textbf{Interpréter} le résultat dans le contexte (unités, arrondi cohérent avec la situation).
\end{enumerate}
\end{methode}

\begin{exoresolu}[Prévisions pour la boutique d'airtime]
\textbf{Énoncé.} On reprend la série de Madame Essomba et la droite de Mayer $y = 9{,}1875\,x + 110{,}28125$ ($y$ en dizaines de milliers de FCFA, $x$ rang du mois).
\begin{enumerate}
\item Estimer le chiffre d'affaires prévisible au mois 10.
\item À partir de quel mois le chiffre d'affaires mensuel dépasserait-il \num{2500000} FCFA si la tendance se poursuivait ?
\end{enumerate}
\tcblower
\textbf{Solution.}
\begin{enumerate}
\item On remplace $x$ par $10$ :
\[ y = 9{,}1875 \times 10 + 110{,}28125 = 91{,}875 + 110{,}28125 = 202{,}15625. \]
Le chiffre d'affaires prévisible au mois 10 est d'environ $202{,}2$ dizaines de milliers de FCFA, soit environ \num{2022000} FCFA.
\item \num{2500000} FCFA correspondent à $y = 250$ (dizaines de milliers). On résout :
\[ 9{,}1875\,x + 110{,}28125 = 250 \iff x = \frac{250 - 110{,}28125}{9{,}1875} = \frac{139{,}71875}{9{,}1875} \approx 15{,}2. \]
Le seuil serait franchi au cours du mois de rang 16 (le premier rang entier supérieur à $15{,}2$).
\end{enumerate}
\end{exoresolu}

\begin{attention}[Prudence avec l'extrapolation]
Une prévision n'est fiable que si la valeur de $x$ reste \emph{proche} du nuage (interpolation ou extrapolation courte). Prévoir le mois 10 à partir de données allant jusqu'au mois 8 est raisonnable ; prévoir le mois 30 avec la même droite serait hasardeux : la tendance peut changer (concurrence, saturation du marché, saisonnalité...). La droite de Mayer décrit une \emph{tendance observée}, pas une loi éternelle. Au Bac, signale toujours, en une phrase, qu'une extrapolation lointaine est à manier avec prudence.
\end{attention}

\courssub{Synthèse visuelle de la méthode}

\begin{popfigure}
\begin{tikzpicture}[
    node distance=0.55cm,
    etape/.style={rectangle, rounded corners, draw=popblue, fill=popblueL, text width=6.2cm, align=center, font=\small, inner sep=5pt},
    res/.style={rectangle, rounded corners, draw=poporange, fill=poporangeL, text width=6.2cm, align=center, font=\small\bfseries, inner sep=5pt},
    fleche/.style={-{Stealth[length=2.5mm]}, thick, popdark}
]
\node[etape] (e1) {\textbf{1. Ordonner} : ranger les points par valeurs croissantes de $x$};
\node[etape, below=of e1] (e2) {\textbf{2. Partager} : deux sous-nuages de même effectif (ou presque)};
\node[etape, below=of e2] (e3) {\textbf{3. Calculer} les points moyens $G_1(\bar{x}_1;\bar{y}_1)$ et $G_2(\bar{x}_2;\bar{y}_2)$};
\node[etape, below=of e3] (e4) {\textbf{4. Équation} : $a=\dfrac{\bar{y}_2-\bar{y}_1}{\bar{x}_2-\bar{x}_1}$, puis $b=\bar{y}_1-a\bar{x}_1$};
\node[etape, below=of e4] (e5) {\textbf{5. Tracer} la droite $(G_1G_2)$ sur le nuage};
\node[res, below=of e5] (r) {Droite de Mayer : $y=ax+b$ \; $\Rightarrow$ \; prévisions};
\draw[fleche] (e1) -- (e2);
\draw[fleche] (e2) -- (e3);
\draw[fleche] (e3) -- (e4);
\draw[fleche] (e4) -- (e5);
\draw[fleche] (e5) -- (r);
\end{tikzpicture}
\legende{Organigramme récapitulatif des cinq étapes de la méthode de Mayer.}
\end{popfigure}

\courssub{Pour aller un peu plus loin}

\begin{savaistu}[La méthode des moindres carrés (hors programme)]
Il existe une autre méthode d'ajustement affine, plus précise : la \cle{méthode des moindres carrés}. Elle consiste à chercher la droite qui rend \emph{minimale} la somme des carrés des écarts verticaux entre les points du nuage et la droite. C'est elle qu'utilisent les calculatrices et les tableurs lorsque tu demandes une « régression linéaire ». Elle est \textbf{hors programme} en Terminale A, mais il est utile d'en connaître l'existence : au Bac, sauf indication contraire de l'énoncé, c'est la \textbf{méthode de Mayer} qui est attendue, avec ses cinq étapes clairement rédigées. La droite des moindres carrés, comme celle de Mayer, passe par le point moyen $G$.
\end{savaistu}

\begin{aretenir}[L'essentiel de la section]
\begin{itemize}
\item Un \cle{ajustement affine} remplace un nuage allongé par une droite $y = ax + b$.
\item La \cle{droite de Mayer} s'obtient en partageant le nuage \textbf{ordonné par $x$ croissants} en deux sous-nuages, puis en traçant $(G_1G_2)$, où $G_1$ et $G_2$ sont les points moyens des sous-nuages.
\item $a = \dfrac{\bar{y}_2 - \bar{y}_1}{\bar{x}_2 - \bar{x}_1}$ et $b = \bar{y}_1 - a\bar{x}_1$.
\item La droite de Mayer passe par le point moyen $G$ du nuage complet : c'est un test de vérification de tes calculs.
\item Prévision : on substitue dans l'équation (dans un sens ou dans l'autre) ; toute extrapolation loin du nuage doit être annoncée avec prudence.
\end{itemize}
\end{aretenir}

\begin{exercice}[À toi de jouer]
Une coopérative agricole de l'Ouest relève la quantité de maïs vendue $y$ (en tonnes) en fonction du rang $x$ de la semaine :
\begin{center}
\begin{tabular}{|c|*{6}{c}|}
\hline
\rowcolor{popblueL} Semaine $x_i$ & 1 & 2 & 3 & 4 & 5 & 6 \\
\hline
Quantité $y_i$ & 12 & 14 & 13 & 17 & 19 & 21 \\
\hline
\end{tabular}
\end{center}
\begin{enumerate}
\item Représenter le nuage de points.
\item Déterminer l'équation de la droite de Mayer.
\item Vérifier que le point moyen $G$ appartient à cette droite.
\item Estimer la quantité vendue à la semaine 8, puis déterminer à partir de quelle semaine les ventes dépasseraient 30 tonnes.
\end{enumerate}
\end{exercice}

\begin{corrige}[Exercice : coopérative agricole]
\begin{enumerate}
\item Nuage : points $(1;12)$, $(2;14)$, $(3;13)$, $(4;17)$, $(5;19)$, $(6;21)$, allongé et croissant : l'ajustement affine est justifié.
\item Les $x$ sont déjà croissants ; $n = 6$, deux sous-nuages de 3 points.
\[ G_1 : \bar{x}_1 = \frac{1+2+3}{3} = 2, \quad \bar{y}_1 = \frac{12+14+13}{3} = 13 ; \qquad G_2 : \bar{x}_2 = \frac{4+5+6}{3} = 5, \quad \bar{y}_2 = \frac{17+19+21}{3} = 19. \]
\[ a = \frac{19-13}{5-2} = \frac{6}{3} = 2, \qquad b = 13 - 2 \times 2 = 9. \]
Droite de Mayer : $y = 2x + 9$.
\item $G : \bar{x} = \dfrac{1+2+3+4+5+6}{6} = 3{,}5$ et $\bar{y} = \dfrac{12+14+13+17+19+21}{6} = \dfrac{96}{6} = 16$. Or $2 \times 3{,}5 + 9 = 16 = \bar{y}$ : $G$ appartient bien à la droite.
\item Semaine 8 : $y = 2 \times 8 + 9 = 25$ tonnes. Seuil de 30 tonnes : $2x + 9 = 30 \iff x = 10{,}5$ ; les ventes dépasseraient 30 tonnes à partir de la semaine 11 (prévision lointaine, à interpréter avec prudence).
\end{enumerate}
\end{corrige}

\newpage

\coursec{Activité d'intégration}

\coursec{Activité d'intégration : Gestion statistique de la coopérative Nyambé}

\courssub{Objectifs de l'activité}

À l'issue de cette activité d'intégration, tu seras capable de mobiliser \emph{simultanément} l'ensemble des savoir-faire de la leçon dans une situation authentique de gestion d'une coopérative agricole :

\begin{itemize}
\item organiser une série statistique brute en classes et construire le tableau des effectifs et des fréquences ;
\item construire et interpréter un histogramme et un polygone des fréquences cumulées croissantes ;
\item calculer et interpréter la moyenne, la médiane et les quartiles d'une série groupée ;
\item construire un nuage de points, déterminer le point moyen et ajuster le nuage par la méthode de Mayer ;
\item utiliser l'ajustement pour formuler une prévision et une recommandation argumentée, en discutant sa fiabilité.
\end{itemize}

Au-delà des calculs, cette activité développe des savoir-être essentiels : travail en équipe, rigueur dans la présentation des données, esprit critique face à une prévision, et capacité à convaincre une assemblée en 5 minutes.

\courssub{Situation}

La coopérative agricole \og{}Nyambé\fg{} du quartier Mballa II (Yaoundé) prépare son assemblée générale annuelle. Le bureau doit prendre deux décisions importantes et te sollicite, toi et ton groupe, comme statisticiens-conseils.

\textbf{Premier dossier : le panier vivrier des ménages.} Pour négocier des achats groupés de vivres (manioc, plantain, macabo), la coopérative a enquêté sur la dépense mensuelle en vivres (en milliers de FCFA) de 40 ménages du quartier. Voici les données brutes :

\begin{center}
\begin{tabular}{|c|c|c|c|c|c|c|c|c|c|}
\hline
22 & 35 & 18 & 41 & 27 & 33 & 38 & 25 & 31 & 45 \\
\hline
29 & 36 & 24 & 39 & 32 & 28 & 34 & 21 & 37 & 30 \\
\hline
26 & 33 & 42 & 19 & 35 & 31 & 28 & 36 & 24 & 38 \\
\hline
34 & 29 & 32 & 27 & 44 & 36 & 52 & 30 & 35 & 57 \\
\hline
\end{tabular}
\end{center}

Le bureau souhaite un \textbf{court rapport} comprenant un histogramme, le polygone des fréquences cumulées croissantes, ainsi que la moyenne, la médiane et les quartiles, avec une phrase d'interprétation pour chaque paramètre.

\textbf{Deuxième dossier : le prix du sac de manioc.} La coopérative hésite sur le prix de vente du sac de manioc de \SI{50}{kg}. Sur les 10 derniers mois, elle a relevé le prix pratiqué $x$ (en milliers de FCFA) et le nombre de sacs vendus $y$ :

\begin{center}
\begin{tabular}{|l|c|c|c|c|c|c|c|c|c|c|}
\hline
Mois & 1 & 2 & 3 & 4 & 5 & 6 & 7 & 8 & 9 & 10 \\
\hline
Prix $x_i$ & 10 & 11 & 12 & 13 & 14 & 15 & 16 & 17 & 18 & 19 \\
\hline
Ventes $y_i$ & 60 & 57 & 55 & 52 & 50 & 47 & 45 & 42 & 40 & 37 \\
\hline
\end{tabular}
\end{center}

Le bureau envisage de fixer le prix à \num{17000} FCFA. Il te demande de prévoir les ventes à ce prix et de formuler une recommandation argumentée.

\courssub{Consignes}

Travail en groupes de 4. Restitution orale de 5 minutes par groupe.

\textbf{Dossier 1 — Le panier vivrier}

\begin{enumerate}
\item Regrouper les 40 données en classes d'amplitude 10 : $[10\,;\,20[$, $[20\,;\,30[$, $[30\,;\,40[$, $[40\,;\,50[$, $[50\,;\,60[$. Dresser le tableau complet : effectifs, fréquences, effectifs cumulés croissants, fréquences cumulées croissantes.
\item Construire l'histogramme des effectifs. Justifier le choix de ce diagramme plutôt qu'un diagramme à bâtons.
\item Construire le polygone des fréquences cumulées croissantes (en \%).
\item Calculer la moyenne de la série groupée. Interpréter le résultat en une phrase destinée à l'assemblée.
\item Déterminer la médiane et les quartiles $Q_1$ et $Q_3$ (on prendra le centre de la classe contenant le rang visé). Interpréter la médiane et l'écart interquartile.
\end{enumerate}

\textbf{Dossier 2 — Le prix du sac de manioc}

\begin{enumerate}
\setcounter{enumi}{5}
\item Représenter le nuage de points $M_i(x_i\,;\,y_i)$ et placer le point moyen $G$.
\item Déterminer une droite d'ajustement affine par la méthode de Mayer (partage du nuage en deux sous-nuages de 5 points).
\item Prévoir les ventes pour un prix de \num{17000} FCFA, puis pour \num{20000} FCFA.
\item Formuler une recommandation de prix au bureau, en discutant la fiabilité de la prévision (qualité de l'alignement, taille de l'échantillon, facteurs non pris en compte).
\end{enumerate}

\courssub{Ressources à mobiliser}

\begin{aretenir}[Boîte à outils de la leçon]
\begin{itemize}
\item Fréquence d'une classe : $f_i = \dfrac{n_i}{N}$ ; fréquence cumulée croissante : somme des fréquences des classes situées à gauche.
\item Moyenne d'une série groupée : $\bar{x} = \dfrac{\sum n_i c_i}{N}$, où $c_i$ est le centre de la classe.
\item Médiane : valeur qui partage la série ordonnée en deux effectifs égaux ; pour une série groupée, on retient le centre de la \cle{classe médiane} (celle où l'effectif cumulé atteint $N/2$). Même principe pour $Q_1$ (rang $N/4$) et $Q_3$ (rang $3N/4$).
\item Point moyen : $G\left(\bar{x}\,;\,\bar{y}\right)$ avec $\bar{x} = \dfrac{1}{n}\sum x_i$ et $\bar{y} = \dfrac{1}{n}\sum y_i$.
\item Méthode de Mayer : partager le nuage en deux sous-nuages, calculer les points moyens $G_1$ et $G_2$ ; la droite de Mayer est la droite $(G_1G_2)$, d'équation $y = ax + b$ avec $a = \dfrac{\bar{y}_2 - \bar{y}_1}{\bar{x}_2 - \bar{x}_1}$.
\end{itemize}
\end{aretenir}

\courssub{Démarche et corrigé commenté}

\begin{exoresolu}[Corrigé commenté de l'activité]
\textbf{Dossier 1.} 
\begin{enumerate}
\item Le dépouillement donne : classe $[10;20[$ : 2 données (18, 19) ; $[20;30[$ : 12 données ; $[30;40[$ : 20 données ; $[40;50[$ : 4 données (41, 42, 44, 45) ; $[50;60[$ : 2 données (52, 57). Total : $2+12+20+4+2 = 40$. Le tableau complet :

\begin{center}
\begin{tabular}{|l|c|c|c|c|c|}
\hline
\rowcolor{popblueL} Classe & $[10;20[$ & $[20;30[$ & $[30;40[$ & $[40;50[$ & $[50;60[$ \\
\hline
Effectif $n_i$ & 2 & 12 & 20 & 4 & 2 \\
\hline
Fréquence $f_i$ & 0,05 & 0,30 & 0,50 & 0,10 & 0,05 \\
\hline
ECC & 2 & 14 & 34 & 38 & 40 \\
\hline
FCC (en \%) & 5 & 35 & 85 & 95 & 100 \\
\hline
\end{tabular}
\end{center}

\item L'histogramme représente une variable quantitative \emph{continue} regroupée en classes : des rectangles \emph{juxtaposés} dont les bases sont les classes et les hauteurs les effectifs (2 ; 12 ; 20 ; 4 ; 2). Un diagramme à bâtons serait incorrect car il est réservé aux séries discrètes non groupées.

\begin{popfigure}
\begin{tikzpicture}
\begin{axis}[
    width=\linewidth,
    height=7cm,
    ybar interval,
    xtick={10,20,30,40,50,60},
    xticklabels={10,20,30,40,50,60},
    xlabel={Dépense mensuelle (en milliers de FCFA)},
    ylabel={Effectif},
    ymin=0, ymax=22,
    axis lines=left,
    tick align=outside,
    bar width=10,
    fill=popblue!70,
    draw=popdark,
]
\addplot coordinates {
    (10,2) (20,12) (30,20) (40,4) (50,2) (60,0)
};
\end{axis}
\end{tikzpicture}
\legende{Histogramme des effectifs de la dépense mensuelle en vivres (40 ménages).}
\end{popfigure}

\item Le polygone des fréquences cumulées croissantes relie les points $(10\,;\,0)$, $(20\,;\,5)$, $(30\,;\,35)$, $(40\,;\,85)$, $(50\,;\,95)$, $(60\,;\,100)$ : on cumule aux \emph{bornes droites} des classes.

\begin{popfigure}
\begin{tikzpicture}
\begin{axis}[
    width=\linewidth,
    height=7cm,
    xlabel={Dépense mensuelle (en milliers de FCFA)},
    ylabel={Fréquence cumulée (\%)},
    xmin=10, xmax=60,
    ymin=0, ymax=105,
    axis lines=left,
    tick align=outside,
    grid=major,
    grid style={popline},
]
\addplot[poporange, mark=*, thick] coordinates {
    (10,0) (20,5) (30,35) (40,85) (50,95) (60,100)
};
\end{axis}
\end{tikzpicture}
\legende{Polygone des fréquences cumulées croissantes (en \%).}
\end{popfigure}

\item Moyenne : centres des classes $c_i$ : 15 ; 25 ; 35 ; 45 ; 55.
\[
\bar{x} = \frac{2\times 15 + 12\times 25 + 20\times 35 + 4\times 45 + 2\times 55}{40} = \frac{30+300+700+180+110}{40} = \frac{\num{1320}}{40} = 33.
\]
Interprétation : \og{}Un ménage du quartier dépense en moyenne environ \num{33000} FCFA par mois en vivres.\fg{}

\item Médiane : $N/2 = 20$ ; l'ECC atteint 20 dans la classe $[30;40[$ (classe médiane) ; on retient $Me \approx 35$ (milliers de FCFA). 
$Q_1$ : rang $N/4 = 10$, atteint dans $[20;30[$, donc $Q_1 \approx 25$. 
$Q_3$ : rang $3N/4 = 30$, atteint dans $[30;40[$, donc $Q_3 \approx 35$. 
Interprétation : La moitié des ménages dépense moins de \num{35000} FCFA ; 50 \% des ménages dépensent entre \num{25000} et \num{35000} FCFA (écart interquartile $Q_3 - Q_1 = 10$ milliers de FCFA, dispersion modérée).
\end{enumerate}

\tcblower

\textbf{Dossier 2.} 
\begin{enumerate}
\setcounter{enumi}{5}
\item Le nuage de points $M_i(x_i\,;\,y_i)$ est quasi rectiligne et décroissant : un ajustement affine est pertinent. 
Point moyen : $\bar{x} = \dfrac{10+11+\cdots+19}{10} = \dfrac{145}{10} = 14,5$ ; $\bar{y} = \dfrac{60+57+\cdots+37}{10} = \dfrac{485}{10} = 48,5$. Donc $G(14,5\,;\,48,5)$.

\begin{popfigure}
\begin{tikzpicture}
\begin{axis}[
    width=\linewidth,
    height=8cm,
    xlabel={Prix $x$ (en milliers de FCFA)},
    ylabel={Ventes $y$ (nombre de sacs)},
    xmin=9, xmax=21,
    ymin=30, ymax=65,
    axis lines=left,
    tick align=outside,
    grid=major,
    grid style={popline},
]
% Nuage de points
\addplot[only marks, mark=*, mark size=2.5, popgreen] coordinates {
    (10,60) (11,57) (12,55) (13,52) (14,50) (15,47) (16,45) (17,42) (18,40) (19,37)
};
% Point moyen G
\addplot[only marks, mark=*, mark size=3, poppink] coordinates {(14.5,48.5)};
\node[poppink, above right] at (axis cs:14.5,48.5) {$G(14,5;48,5)$};
% Points moyens partiels G1, G2
\addplot[only marks, mark=square*, mark size=3, poporange] coordinates {(12,54.8) (17,42.2)};
\node[poporange, above left] at (axis cs:12,54.8) {$G_1(12;54,8)$};
\node[poporange, below right] at (axis cs:17,42.2) {$G_2(17;42,2)$};
% Droite de Mayer
\addplot[popcurve, domain=9:21, samples=2, thick, popblue] {-2.52*x + 85.04};
\node[popblue, above left] at (axis cs:10,60) {$y = -2,52x + 85,04$};
\end{axis}
\end{tikzpicture}
\legende{Nuage de points, point moyen $G$, points moyens partiels $G_1$, $G_2$ et droite de Mayer.}
\end{popfigure}

\item Méthode de Mayer : partage en deux sous-nuages de 5 points.
Sous-nuage 1 (mois 1 à 5) : 
$G_1\left(\dfrac{10+11+12+13+14}{5}\,;\,\dfrac{60+57+55+52+50}{5}\right) = G_1(12\,;\,54,8)$.
Sous-nuage 2 (mois 6 à 10) : 
$G_2\left(\dfrac{15+16+17+18+19}{5}\,;\,\dfrac{47+45+42+40+37}{5}\right) = G_2(17\,;\,42,2)$.
\[
a = \frac{42,2 - 54,8}{17 - 12} = \frac{-12,6}{5} = -2,52 \quad ; \quad b = 54,8 - (-2,52)\times 12 = 54,8 + 30,24 = 85,04.
\]
Droite de Mayer : $y = -2,52x + 85,04$. Vérification : $G_2$ est bien sur la droite : $-2,52\times 17 + 85,04 = 42,2$.

\item Prévisions : 
Pour $x = 17$ (soit \num{17000} FCFA) : $y = -2,52\times 17 + 85,04 = 42,2$, soit environ \textbf{42 sacs vendus}. 
Pour $x = 20$ (soit \num{20000} FCFA) : $y = -2,52\times 20 + 85,04 = 34,64$, soit environ 35 sacs.

\item Recommandation possible : À \num{17000} FCFA, la coopérative peut espérer vendre environ 42 sacs, chiffre cohérent avec le relevé du mois 8 (42 sacs à 17). Le nuage étant très bien aligné, la prévision est assez fiable \emph{à court terme} et \emph{à l'intérieur de l'intervalle d'observation} $[10;19]$ ; toutefois, 10 relevés seulement, la saisonnalité (fêtes de fin d'année, rentrée scolaire) et la concurrence du marché Mokolo ne sont pas pris en compte : il faut rester prudent, surtout pour une extrapolation au-delà de \num{19000} FCFA.
\end{enumerate}
\end{exoresolu}

\courssub{Bilan de l'activité}

Cette activité a permis de mettre en évidence plusieurs idées force de la leçon :

\begin{itemize}
\item le \cle{regroupement en classes} rend lisible une masse de données brutes, au prix d'une perte d'information (on ne calcule plus que des valeurs approchées : moyenne groupée 33 alors que la moyenne exacte des données brutes vaut 32,8) ;
\item le choix du \cle{diagramme} n'est pas anodin : histogramme pour une variable continue groupée, polygone des FCC pour lire rapidement médiane et quartiles ;
\item les \cle{paramètres de position} se complètent : la moyenne résume, la médiane partage, les quartiles mesurent la dispersion ;
\item l'\cle{ajustement de Mayer} transforme des observations en outil de décision : la coopérative peut désormais simuler l'effet d'un prix sur les ventes ;
\item enfin, une prévision statistique n'est jamais une certitude : sa fiabilité dépend de l'alignement du nuage, de la taille de l'échantillon et de la stabilité du contexte. Un bon statisticien sait calculer, mais aussi \emph{douter} et argumenter.
\end{itemize}

\begin{attention}[Pour la restitution orale]
En 5 minutes, va à l'essentiel : un tableau, un graphique, trois paramètres, une droite, une recommandation. Chaque chiffre annoncé doit être accompagné d'une phrase d'interprétation en langage courant : l'assemblée de Mballa II n'est pas composée de mathématiciens, mais de décideurs !
\end{attention}

\newpage

\coursec{Méthodes et exercices résolus}

\coursec{Vocabulaire et organisation des données}

\begin{definition}[Série statistique à une variable]
On appelle \cle{série statistique} l'ensemble des valeurs prises par un \cle{caractère} étudié sur une \cle{population} d'individus (ou d'objets).
\begin{itemize}
\item Le caractère est \cle{qualitatif} s'il exprime une qualité (couleur, marque, oui/non).
\item Le caractère est \cle{quantitatif discret} s'il ne prend que des valeurs isolées, souvent entières (nombre d'enfants, note sur 20).
\item Le caractère est \cle{quantitatif continu} s'il peut prendre toute valeur dans un intervalle (taille, masse, durée, recette). On le regroupe alors en \cle{classes} d'amplitude $a_i = b_i - a_i$.
\end{itemize}
Pour chaque \cle{modalité} $x_i$ (ou chaque classe), on note :
\begin{itemize}
\item $n_i$ : l'\cle{effectif} (nombre d'individus) ;
\item $N = \sum n_i$ : l'\cle{effectif total} ;
\item $f_i = \dfrac{n_i}{N}$ : la \cle{fréquence} (souvent en \%) ; $\sum f_i = 1$ ;
\item Les \cle{effectifs cumulés croissants} (ECC) : somme des effectifs des modalités $\le x_i$ ;
\item Les \cle{effectifs cumulés décroissants} (ECD) : somme des effectifs des modalités $\ge x_i$.
\end{itemize}
\end{definition}

\begin{aretenir}[Tableau de valeurs et tableau de classes]
\begin{itemize}
\item \textbf{Série discrète} : on liste chaque valeur distincte $x_i$ avec son effectif $n_i$.
\item \textbf{Série continue (groupée en classes)} : on choisit un nombre raisonnable de classes (généralement 5 à 15) de même amplitude si possible, qui partitionnent l'étendue des valeurs. On note pour chaque classe $[a_i\,;\,b_i[$ son centre $c_i = \dfrac{a_i+b_i}{2}$ et son effectif $n_i$.
\end{itemize}
\textbf{Vérification systématique} : avant tout calcul, contrôlez que $\sum n_i = N$ et $\sum f_i = 1$ (ou $100\,\%$).
\end{aretenir}

\coursec{Représentations graphiques}

\begin{methode}[Construire et lire les diagrammes usuels]
\begin{enumerate}
\item \textbf{Identifier le type de caractère} :
\begin{itemize}
\item Qualitatif ou discret à peu de modalités $\rightarrow$ diagramme à \cle{bâtons}, à \cle{bandes} ou \cle{circulaire}.
\item Quantitatif continu (classes) $\rightarrow$ \cle{histogramme} (aires proportionnelles aux effectifs) ou \cle{polygone des fréquences cumulées croissantes}.
\end{itemize}
\item \textbf{Calculer les fréquences} $f_i = \dfrac{n_i}{N}$.
\item \textbf{Adapter l'échelle} :
\begin{itemize}
\item \emph{Circulaire} : angle $\alpha_i = 360^\circ \times f_i$.
\item \emph{Bâtons / Bandes} : hauteur (ou longueur) proportionnelle à $n_i$ ou $f_i$.
\item \emph{Histogramme} : \textbf{aire} du rectangle $=$ $k \times n_i$. Si amplitudes égales, hauteur $\propto n_i$ ; sinon hauteur $h_i = \dfrac{n_i}{\text{amplitude}_i}$ (densité d'effectif).
\item \emph{Polygone des fréquences cumulées} : on marque les points $(b_i\,;\, \text{ECC}_i)$ et on les relie par des segments (on ajoute l'origine de la première classe avec ECC $=0$).
\end{itemize}
\item \textbf{Lire et comparer} : toujours légender axes, unités, titre. Pour comparer deux populations de tailles différentes, utilisez les \emph{fréquences} (ou pourcentages), jamais les seuls effectifs.
\end{enumerate}
\end{methode}

\begin{exoresolu}[Diagramme circulaire : budget d'un ménage à Yaoundé]
Un ménage répartit son budget mensuel de \num{200000} FCFA ainsi : alimentation \num{80000}, transport \num{40000}, scolarité \num{30000}, communication (MTN/Orange) \num{20000}, divers \num{30000}. Construire le tableau des fréquences et angles.
\tcblower
\textbf{Solution.} $N = \num{200000}$.
\begin{center}
\begin{tabular}{|l|c|c|c|}
\hline
\rowcolor{popblueL} Poste & Dépense (FCFA) & Fréquence $f_i$ & Angle $\alpha_i = 360 f_i$ \\
\hline
Alimentation & \num{80000} & $0{,}40$ & $144^\circ$ \\
\hline
Transport & \num{40000} & $0{,}20$ & $72^\circ$ \\
\hline
Scolarité & \num{30000} & $0{,}15$ & $54^\circ$ \\
\hline
Communication & \num{20000} & $0{,}10$ & $36^\circ$ \\
\hline
Divers & \num{30000} & $0{,}15$ & $54^\circ$ \\
\hline
\textbf{Total} & \num{200000} & $1$ & $360^\circ$ \\
\hline
\end{tabular}
\end{center}
\textbf{Vérification} : $\sum f_i = 1$, $\sum \alpha_i = 360^\circ$. On trace les secteurs successifs au rapporteur.
\end{exoresolu}

\begin{exoresolu}[Histogramme : recettes journalières au marché Mokolo]
Un commerçant a relevé ses recettes (en milliers de FCFA) sur $50$ jours :
\begin{center}
\begin{tabular}{|l|c|c|c|c|}
\hline
\rowcolor{popblueL} Recette (kFCFA) & $[10\,;\,20[$ & $[20\,;\,30[$ & $[30\,;\,40[$ & $[40\,;\,50[$ \\
\hline
Effectif $n_i$ & 8 & 15 & 18 & 9 \\
\hline
\end{tabular}
\end{center}
Construire l'histogramme. Les classes ont même amplitude ($10$).
\tcblower
\textbf{Solution.} Amplitude constante $= 10$. Hauteur des rectangles $\propto$ effectifs. On peut prendre 1 cm pour 1 unité d'effectif (ou densité $n_i/10$).
\begin{popfigure}
\begin{tikzpicture}[scale=0.7]
\begin{axis}[
    ybar,
    bar width=1.5cm,
    width=\linewidth,
    height=8cm,
    xlabel={Recettes (kFCFA)},
    ylabel={Effectif},
    xtick={15,25,35,45},
    xticklabels={$[10;20[$, $[20;30[$, $[30;40[$, $[40;50[$},
    ymin=0, ymax=20,
    ymajorgrids=true,
    grid style=dashed,
    popblue,
    fill=popblueL,
    enlarge x limits=0.15,
    enlarge y limits={upper=0.15},
]
\addplot coordinates {(15,8) (25,15) (35,18) (45,9)};
\end{axis}
\end{tikzpicture}
\legende{Histogramme des recettes journalières (amplitudes égales).}
\end{popfigure}
\textbf{Note} : si les amplitudes différaient, on calculerait la densité $h_i = n_i / \text{amplitude}_i$ pour que l'aire $= h_i \times \text{amplitude}_i = n_i$.
\end{exoresolu}

\coursec{Paramètres de position : mode, moyenne, médiane}

\begin{propriete}[Mode, Moyenne, Médiane]
\begin{itemize}
\item \textbf{Mode} : modalité (ou classe \emph{modale}) d'effectif maximal. Une série peut être multimodale.
\item \textbf{Moyenne} $\bar{x}$ :
\begin{itemize}
\item Discrète : $\bar{x} = \dfrac{1}{N}\sum n_i x_i$.
\item En classes : $\bar{x} = \dfrac{1}{N}\sum n_i c_i$ où $c_i = \dfrac{a_i+b_i}{2}$ est le centre de la classe.
\end{itemize}
Propriété : $\sum n_i(x_i - \bar{x}) = 0$ (centre de gravité).
\item \textbf{Médiane} $Me$ (valeur centrale, sans interpolation linéaire au programme) :
\begin{itemize}
\item Discrète : valeurs rangées. Si $N$ impair, $Me =$ valeur de rang $\frac{N+1}{2}$. Si $N$ pair, $Me =$ demi-somme des valeurs de rangs $\frac{N}{2}$ et $\frac{N}{2}+1$ (lecture via ECC).
\item En classes : on repère la \emph{classe médiane} où l'ECC atteint ou dépasse $\frac{N}{2}$. La médiane appartient à cette classe ; par commodité, on retient souvent son centre $c_{med}$ comme valeur approchée.
\end{itemize}
\end{itemize}
\end{propriete}

\begin{methode}[Calculer mode, moyenne, médiane]
\begin{enumerate}
\item Vérifier $N = \sum n_i$.
\item \textbf{Mode} : repérer $n_{max}$.
\item \textbf{Moyenne} : calculer $\sum n_i x_i$ (ou $\sum n_i c_i$) puis diviser par $N$. Garder la forme fractionnaire ou 2 décimales.
\item \textbf{Médiane} : construire la ligne des ECC. Chercher le(s) rang(s) médian(s). Lire la valeur correspondante.
\item \textbf{Conclure} par une phrase interprétée avec l'unité.
\end{enumerate}
\end{methode}

\begin{exoresolu}[Notes de Tle A (série discrète)]
$40$ élèves, notes sur $20$ :
\begin{center}
\begin{tabular}{|l|c|c|c|c|c|}
\hline
\rowcolor{popblueL} Note $x_i$ & 6 & 9 & 11 & 13 & 16 \\
\hline
Effectif $n_i$ & 5 & 8 & 12 & 10 & 5 \\
\hline
\end{tabular}
\end{center}
Déterminer mode, moyenne, médiane.
\tcblower
\textbf{Solution.} $N = 40$.
\begin{enumerate}
\item \textbf{Mode} : $n_{max}=12$ pour $x=11$. Mode $= 11/20$.
\item \textbf{Moyenne} : $\sum n_i x_i = 5\!\times\!6 + 8\!\times\!9 + 12\!\times\!11 + 10\!\times\!13 + 5\!\times\!16 = 444$. $\bar{x} = 444/40 = 11{,}1/20$.
\item \textbf{Médiane} : $N$ pair $\rightarrow$ rangs $20$ et $21$.
ECC : $5$ (note 6), $13$ (note 9), $25$ (note 11), $35$ (note 13), $40$ (note 16).
Rangs $20$ et $21$ sont dans la tranche ECC $=25$ (note 11). $Me = \frac{11+11}{2} = 11/20$.
\end{enumerate}
\textbf{Conclusion} : la classe a pour note moyenne $11{,}1$, médiane $11$ et mode $11$. La série est assez symétrique.
\end{exoresolu}

\begin{exoresolu}[Série en classes : recettes journalières (suite)]
Même tableau que ci-dessus ($N=50$). Calculer la moyenne et la classe médiane.
\tcblower
\textbf{Solution.}
\textbf{Moyenne.} Centres : $15,\ 25,\ 35,\ 45$.
$\sum n_i c_i = 8\!\times\!15 + 15\!\times\!25 + 18\!\times\!35 + 9\!\times\!45 = 120 + 375 + 630 + 405 = 1530$.
$\bar{x} = 1530/50 = 30{,}6$ kFCFA $= \num{30600}$ FCFA.

\textbf{Classe médiane.} $N/2 = 25$. ECC : $8,\ 23,\ 41,\ 50$. Premier ECC $\ge 25$ est $41$ pour la classe $[30\,;\,40[$. 
La médiane appartient à $[30\,;\,40[$ kFCFA. (Centre $= 35$ kFCFA comme valeur approchée).
\end{exoresolu}

\coursec{Quartiles et déciles}

\begin{definition}[Quartiles et déciles (série discrète)]
Soit une série de $N$ valeurs rangées dans l'ordre croissant avec ECC.
\begin{itemize}
\item \textbf{Premier quartile} $Q_1$ : plus petite valeur telle qu'au moins $25\,\%$ des valeurs lui soient $\le$. Rang $r_1 = N/4$ (si entier, valeur de ce rang ; sinon valeur de rang $\lceil r_1 \rceil$).
\item \textbf{Troisième quartile} $Q_3$ : même principe pour $75\,\%$ (rang $r_3 = 3N/4$).
\item \textbf{Déciles} $D_k$ ($k=1..9$) : $10k\,\%$ des valeurs $\le D_k$. Rang $r_k = kN/10$ (même règle d'arrondi au supérieur). Noter que $D_5 = Me$.
\item \textbf{Écart interquartile} $EI = Q_3 - Q_1$ : mesure la dispersion de la moitié centrale ($50\,\%$ des valeurs).
\end{itemize}
\end{definition}

\begin{methode}[Déterminer $Q_1$, $Me$, $Q_3$, $D_k$]
\begin{enumerate}
\item Ranger les valeurs, dresser les ECC.
\item Calculer le rang théorique $r = \frac{pN}{100}$ (pour $p=25, 50, 75, 10k$).
\item Si $r$ est entier $\rightarrow$ valeur de rang $r$. Sinon $\rightarrow$ valeur de rang $\lfloor r \rfloor + 1$ (entier immédiatement supérieur).
\item Lire la valeur correspondante dans le tableau des ECC.
\item Interpréter $EI = Q_3 - Q_1$.
\end{enumerate}
\end{methode}

\begin{exoresolu}[Durées d'appels au service client (CAMTEL)]
$60$ appels, durées en minutes :
\begin{center}
\begin{tabular}{|l|c|c|c|c|c|c|}
\hline
\rowcolor{popblueL} Durée & 2 & 4 & 6 & 8 & 10 & 12 \\
\hline
Effectif & 6 & 12 & 18 & 12 & 9 & 3 \\
\hline
\end{tabular}
\end{center}
Déterminer $Q_1$, $Me$, $Q_3$, $D_9$ et interpréter $EI$.
\tcblower
\textbf{Solution.} $N=60$. ECC : $6,\ 18,\ 36,\ 48,\ 57,\ 60$.
\begin{itemize}
\item $Q_1$ : $r = 60/4 = 15$ (entier) $\rightarrow$ valeur rang $15$. $6 < 15 \le 18 \Rightarrow Q_1 = 4$ min.
\item $Me$ : rangs $30$ et $31$. $18 < 30,31 \le 36 \Rightarrow$ valeurs $6$. $Me = 6$ min.
\item $Q_3$ : $r = 3\!\times\!60/4 = 45$ (entier) $\rightarrow$ valeur rang $45$. $36 < 45 \le 48 \Rightarrow Q_3 = 8$ min.
\item $D_9$ : $r = 9\!\times\!60/10 = 54$ (entier) $\rightarrow$ valeur rang $54$. $48 < 54 \le 57 \Rightarrow D_9 = 10$ min.
\end{itemize}
$EI = Q_3 - Q_1 = 8 - 4 = 4$ min.
\textbf{Interprétation} : la moitié centrale des appels dure entre $4$ et $8$ minutes. $90\,\%$ des appels durent $\le 10$ min.
\end{exoresolu}

\coursec{Séries à deux caractères : nuage de points et point moyen}

\begin{definition}[Série statistique à deux caractères]
On observe deux caractères quantitatifs $X$ et $Y$ sur une même population de taille $n$. On obtient $n$ couples $(x_i\,;\,y_i)$.
Le \cle{nuage de points} est l'ensemble des points $M_i(x_i\,;\,y_i)$ dans un repère orthonormé.
Le \cle{point moyen} $G$ a pour coordonnées les moyennes des deux caractères :
\[ x_G = \bar{x} = \frac{1}{n}\sum x_i \qquad y_G = \bar{y} = \frac{1}{n}\sum y_i. \]
Toute droite d'ajustement affine pertinente passe par $G$.
\end{definition}

\begin{methode}[Construire le nuage et trouver $G$]
\begin{enumerate}
\item Choisir les axes : $X$ en abscisses (souvent variable explicative ou temps), $Y$ en ordonnées. Échelles adaptées (origine non nulle possible).
\item Placer les points $M_i(x_i\,;\,y_i)$.
\item Calculer $\bar{x}$ et $\bar{y}$ pour obtenir $G(\bar{x}\,;\,\bar{y})$.
\item Observer la forme : alignement $\rightarrow$ ajustement affine pertinent ; forme courbe ou nuage diffus $\rightarrow$ ajustement affine non pertinent.
\end{enumerate}
\end{methode}

\begin{exoresolu}[Chiffre d'affaires d'un cybercafé à Douala]
\begin{center}
\begin{tabular}{|l|c|c|c|c|c|c|}
\hline
\rowcolor{popblueL} Mois $x_i$ & 1 & 2 & 3 & 4 & 5 & 6 \\
\hline
CA $y_i$ (100k FCFA) & 8 & 10 & 11 & 13 & 14 & 16 \\
\hline
\end{tabular}
\end{center}
Représenter le nuage et déterminer $G$.
\tcblower
\textbf{Solution.} $n=6$.
$\bar{x} = (1+2+3+4+5+6)/6 = 3{,}5$.
$\bar{y} = (8+10+11+13+14+16)/6 = 72/6 = 12$.
$G(3{,}5\,;\,12)$.
\begin{popfigure}
\begin{tikzpicture}[scale=0.8]
\begin{axis}[
    width=\linewidth, height=7cm,
    xlabel={Mois $x$}, ylabel={CA $y$ (100k FCFA)},
    xmin=0, xmax=7, ymin=5, ymax=18,
    grid=major, grid style={dashed, popmuted},
    scatter, only marks, mark=*, mark size=2.5, popblue,
]
\addplot coordinates {(1,8) (2,10) (3,11) (4,13) (5,14) (6,16)};
\addplot[mark=*, mark size=3, poporange, thick] coordinates {(3.5,12)};
\node[poporange, above right] at (axis cs:3.5,12) {$G$};
\end{axis}
\end{tikzpicture}
\legende{Nuage de points et point moyen $G(3{,}5;12)$. Les points suggèrent une tendance linéaire croissante.}
\end{popfigure}
\end{exoresolu}

\coursec{Ajustement affine par la méthode de Mayer et prévisions}

\begin{methode}[Méthode de Mayer (droite des points moyens)]
\begin{enumerate}
\item \textbf{Trier} les $n$ points par $x$ croissant.
\item \textbf{Partager} en deux sous-nuages d'effectifs $n_1$ et $n_2$ égaux (ou $|n_1-n_2|=1$ si $n$ impair). Le partage se fait \textbf{uniquement selon $x$}.
\item \textbf{Calculer les points moyens} de chaque sous-nuage :
\[ G_1\Bigl(\frac{\sum_{1}^{n_1} x_i}{n_1}\,;\,\frac{\sum_{1}^{n_1} y_i}{n_1}\Bigr), \quad
   G_2\Bigl(\frac{\sum_{n_1+1}^{n} x_i}{n_2}\,;\,\frac{\sum_{n_1+1}^{n} y_i}{n_2}\Bigr). \]
\item \textbf{Équation de la droite $(G_1G_2)$} :
\[ a = \frac{y_{G_2}-y_{G_1}}{x_{G_2}-x_{G_1}}, \qquad b = y_{G_1} - a\,x_{G_1} \quad (\text{ou avec } G_2). \]
La droite d'équation $y = ax + b$ passe par $G_1$, $G_2$ et $G$.
\item \textbf{Prévision} : pour $x_0$ donné, estimer $y_0 = a x_0 + b$ (en précisant que c'est une extrapolation si $x_0$ hors intervalle observé).
\end{enumerate}
\end{methode}

\begin{exoresolu}[Prévision ventes forfaits internet à Bafoussam]
\begin{center}
\begin{tabular}{|l|c|c|c|c|c|c|}
\hline
\rowcolor{popblueL} Semaine $x_i$ & 1 & 2 & 3 & 4 & 5 & 6 \\
\hline
Ventes $y_i$ (dizaines) & 12 & 14 & 15 & 18 & 19 & 22 \\
\hline
\end{tabular}
\end{center}
1. Droite de Mayer. 2. Prévoir semaine 8.
\tcblower
\textbf{Solution.}
\textbf{1. Droite de Mayer.} $n=6 \rightarrow$ deux sous-nuages de $3$ points.
\begin{itemize}
\item Sous-nuage 1 : $(1,12),\ (2,14),\ (3,15)$.
$G_1\bigl(\frac{6}{3}=2\,;\,\frac{41}{3}\approx 13{,}67\bigr)$.
\item Sous-nuage 2 : $(4,18),\ (5,19),\ (6,22)$.
$G_2\bigl(\frac{15}{3}=5\,;\,\frac{59}{3}\approx 19{,}67\bigr)$.
\end{itemize}
$a = \frac{\frac{59}{3}-\frac{41}{3}}{5-2} = \frac{18/3}{3} = 2$.
$b = \frac{41}{3} - 2\times 2 = \frac{41-12}{3} = \frac{29}{3} \approx 9{,}67$.
Droite : $y = 2x + \frac{29}{3}$.
Vérification : $G$ global $(\bar{x}=3{,}5\,;\,\bar{y}=99/6=16{,}5)$. $2\times 3{,}5 + 29/3 = 7 + 9{,}67 = 16{,}67 \approx 16{,}5$ (arrondi). OK.

\textbf{2. Prévision semaine 8 ($x_0=8$).}
$y_0 = 2\times 8 + \frac{29}{3} = \frac{77}{3} \approx 25{,}67$ dizaines $\approx 257$ forfaits.
\begin{popfigure}
\begin{tikzpicture}[scale=0.8]
\begin{axis}[
    width=\linewidth, height=7cm,
    xlabel={Semaine $x$}, ylabel={Ventes $y$ (dizaines)},
    xmin=0, xmax=9, ymin=10, ymax=28,
    grid=major, grid style={dashed, popmuted},
]
% Points
\addplot[only marks, mark=*, mark size=2.5, popblue] coordinates {(1,12) (2,14) (3,15) (4,18) (5,19) (6,22)};
% Droite de Mayer
\addplot[domain=0:8, samples=2, poporange, thick] {2*x + 29/3};
\node[poporange, above left] at (axis cs:7, 2*7+29/3) {$y=2x+\frac{29}{3}$};
% Points moyens G1, G2
\addplot[only marks, mark=*, mark size=3, popgreen] coordinates {(2,41/3) (5,59/3)};
\node[popgreen, above] at (axis cs:2,41/3) {$G_1$};
\node[popgreen, above] at (axis cs:5,59/3) {$G_2$};
% Point moyen global G
\addplot[only marks, mark=*, mark size=3, poppink] coordinates {(3.5,16.5)};
\node[poppink, below right] at (axis cs:3.5,16.5) {$G$};
\end{axis}
\end{tikzpicture}
\legende{Nuage, points moyens $G_1, G_2, G$ et droite de Mayer.}
\end{popfigure}
\textbf{Conclusion} : on prévoit environ $260$ forfaits pour la semaine 8, sous réserve du maintien de la tendance.
\end{exoresolu}

\begin{attention}[Les 5 erreurs qui coûtent des points au Bac]
\begin{enumerate}
\item \textbf{Oublier de vérifier $N = \sum n_i$ et $\sum f_i = 1$.} Une erreur de recopie d'effectif propage l'erreur sur tous les paramètres (moyenne, médiane, quartiles, droite de Mayer).
\item \textbf{Confondre médiane et moyenne.} La médiane se détermine \textbf{uniquement via les rangs et les ECC}, jamais par calcul direct sur les valeurs. Pour $N$ pair : demi-somme des deux valeurs centrales (même valeur $\rightarrow$ médiane = cette valeur).
\item \textbf{Utiliser les bornes au lieu des centres de classes.} Pour la moyenne d'une série groupée, la formule $\bar{x} = \frac{\sum n_i c_i}{N}$ avec $c_i = \frac{a_i+b_i}{2}$ est \textbf{obligatoire}. L'amplitude ne sert que pour l'histogramme (densité).
\item \textbf{Mal partager le nuage dans la méthode de Mayer.} Le partage se fait \textbf{selon $x$ croissant} en deux groupes d'effectifs égaux (ou $\pm 1$). Ne jamais partager selon $y$, ni tracer « à l'œil ». La droite doit passer par $G_1$, $G_2$ et $G$.
\item \textbf{Négliger la rédaction finale.} Tout résultat numérique doit être suivi d'une phrase de conclusion avec l'unité (FCFA, min, notes, forfaits...). Pour une prévision : « Sous réserve que la tendance observée se poursuive, on prévoit... ».
\end{enumerate}
\end{attention}

\begin{savaistu}[Statistiques au Cameroun]
L'\textbf{INS} (Institut National de la Statistique) réalise les recensements (RGPH) et les enquêtes (ECAM) qui produisent les séries statistiques officielles : population par région, taux de scolarisation, pauvreté, inflation. Les quartiles et déciles servent à définir les seuils de pauvreté (ex: 1er quintile = 20\% les plus pauvres). La méthode de Mayer, bien que simple, est l'ancêtre des régressions linéaires utilisées aujourd'hui par la BEAC, la BICEC ou les opérateurs MTN/Orange pour prévoir la consommation de data ou le trafic voix.
\end{savaistu}

\newpage

\coursec{Exercices}

\coursec{Vocabulaire et organisation des données}

\begin{definition}[Population, individu, caractère]
On appelle \textbf{population} l'ensemble des êtres ou objets étudiés. Chaque élément de la population est un \textbf{individu} (ou unité statistique). Un \textbf{caractère} est une propriété observée sur chaque individu. Il peut être :
\begin{itemize}
    \item \textbf{qualitatif} (ex: opérateur téléphonique MTN/Orange, couleur, sexe) ;
    \item \textbf{quantitatif discret} (valeurs numérables, ex: nombre de sacs de maïs, nombre de buts) ;
    \item \textbf{quantitatif continu} (valeurs dans un intervalle, ex: taille, distance, temps).
\end{itemize}
\end{definition}

\begin{definition}[Effectif, fréquence, effectifs cumulés]
Pour une modalité $x_i$ (ou une classe), on note :
\begin{itemize}
    \item $n_i$ : \textbf{effectif} (nombre d'individus prenant cette valeur) ;
    \item $f_i = \dfrac{n_i}{N}$ : \textbf{fréquence} ($N = \sum n_i$ effectif total) ; souvent exprimée en \%;
    \item $N_i$ : \textbf{effectif cumulé croissant (ECC)} : somme des effectifs des valeurs $\le x_i$ ;
    \item $F_i = \dfrac{N_i}{N}$ : \textbf{fréquence cumulée croissante (FCC)}.
\end{itemize}
\end{definition}

\begin{propriete}[Série groupée en classes]
Pour un caractère continu, on regroupe les valeurs en \textbf{classes} d'amplitude $a = b-a$ (ex: $[a\,;\,b[$). Le \textbf{centre de classe} (ou milieu) est $c = \dfrac{a+b}{2}$. Il représente la classe pour les calculs (moyenne, etc.).
\end{propriete}

\begin{methode}[Construire un tableau statistique complet]
\begin{enumerate}
    \item Trier les données (ordre croissant) si série brute.
    \item Lister les modalités (ou classes) sans omission.
    \item Compter les effectifs $n_i$.
    \item Calculer $N = \sum n_i$, puis les fréquences $f_i = n_i/N$ (\%).
    \item Calculer les ECC $N_i$ et FCC $F_i$.
    \item Pour classes : ajouter la colonne des centres $c_i$.
\end{enumerate}
\end{methode}

\begin{exemplebox}[Exemple : Ventes de maïs au marché Mokolo]
Série brute (10 jours) : $12\ ;\ 15\ ;\ 9\ ;\ 12\ ;\ 18\ ;\ 12\ ;\ 10\ ;\ 15\ ;\ 12\ ;\ 20$.
\begin{center}
\begin{tabularx}{\linewidth}{|c|X|X|X|X|X|}
\hline
\rowcolor{popblueL}
$x_i$ & $n_i$ & $f_i\ (\%)$ & $N_i$ & $F_i\ (\%)$ \\
\hline
9 & 1 & 10 & 1 & 10 \\
10 & 1 & 10 & 2 & 20 \\
12 & 4 & 40 & 6 & 60 \\
15 & 2 & 20 & 8 & 80 \\
18 & 1 & 10 & 9 & 90 \\
20 & 1 & 10 & 10 & 100 \\
\hline
\end{tabularx}
\end{center}
\end{exemplebox}

\coursec{Représentations graphiques}

\begin{propriete}[Diagrammes usuels]
\begin{itemize}
    \item \textbf{Diagramme circulaire (camembert)} : caractère qualitatif ou discret à peu de modalités. Angle du secteur $= 360^\circ \times f_i$.
    \item \textbf{Diagramme à bâtons} : caractère quantitatif discret. Bâton de hauteur $n_i$ (ou $f_i$) centré sur $x_i$. Largeur constante, espaces entre bâtons.
    \item \textbf{Histogramme} : caractère quantitatif continu groupé en classes. Rectangle sur chaque classe $[a\,;\,b[$ de hauteur $h_i = \dfrac{n_i}{a}$ (densité) pour que \textbf{l'aire} $= n_i$. Si amplitudes égales, hauteur $= n_i$.
    \item \textbf{Polygone des fréquences cumulées croissantes} : série groupée. Points $M_i(b_i\,;\,F_i)$ (borne supérieure, FCC en \%). Reliés par segments. On ajoute $(a_1\,;\,0)$.
\end{itemize}
\end{propriete}

\begin{attention}[Pièges graphiques]
\begin{itemize}
    \item Histogramme $\neq$ diagramme à bâtons : pas d'espaces entre classes, aire proportionnelle à l'effectif (densité si amplitudes inégales).
    \item Diagramme circulaire : somme des angles $= 360^\circ$, somme des fréquences $= 100\%$.
    \item Polygone des FCC : abscisses = bornes \textbf{supérieures} des classes (pas les centres).
\end{itemize}
\end{attention}

\begin{savaistu}[Le diagramme en boîte (box-plot)]
Bien que non explicitement au programme, c'est la représentation standard pour visualiser médiane, quartiles, extrêmes et valeurs aberrantes. Utile en sciences économiques et biologie.
\end{savaistu}

\coursec{Paramètres de position : mode, moyenne, médiane}

\begin{definition}[Mode]
Valeur (ou classe) d'effectif maximal. Une série peut être \textbf{unimodale}, \textbf{bimodale}, etc. Pour une classe modale, on donne l'intervalle.
\end{definition}

\begin{definition}[Moyenne]
$\bar{x} = \dfrac{\sum_{i=1}^k n_i x_i}{\sum_{i=1}^k n_i} = \dfrac{\text{Somme des valeurs}}{N}$.
Pour classe : on utilise le centre $c_i$.
\emph{Propriété} : $\sum n_i(x_i - \bar{x}) = 0$ (barycentre).
\end{definition}

\begin{definition}[Médiane $M_e$]
Valeur qui partage la série en deux effectifs égaux : au moins $50\%$ des valeurs $\le M_e$ et au moins $50\% \ge M_e$.
\begin{itemize}
    \item \textbf{Série discrète} (effectif $N$) : 
    \begin{itemize}
        \item $N$ impair : $M_e =$ valeur au rang $\frac{N+1}{2}$.
        \item $N$ pair : $M_e =$ moyenne des valeurs aux rangs $\frac{N}{2}$ et $\frac{N}{2}+1$.
    \end{itemize}
    \item \textbf{Série groupée} : lecture graphique sur le polygone des FCC (abscisse pour $F=50\%$). L'interpolation linéaire (formule) est \textbf{hors programme}.
\end{itemize}
\end{definition}

\begin{aretenir}[Choix du paramètre]
\begin{itemize}
    \item \textbf{Moyenne} : utilise toutes les valeurs, sensible aux valeurs extrêmes (outliers). Indiquée pour distributions symétriques.
    \item \textbf{Médiane} : robuste (insensible aux extrêmes). Indiquée pour distributions asymétriques (salaires, tailles).
    \item \textbf{Mode} : seule mesure pour caractère qualitatif. Pic de l'histogramme.
\end{itemize}
\end{aretenir}

\coursec{Quartiles et déciles}

\begin{definition}[Quartiles et déciles]
Les \textbf{quartiles} partagent la série en 4 parties d'effectifs égaux :
$Q_1$ (25\%), $Q_2 = M_e$ (50\%), $Q_3$ (75\%).
Les \textbf{déciles} $D_1, \dots, D_9$ partagent en 10 parties (10\%, 20\%, ..., 90\%).
$D_5 = M_e$.
\end{definition}

\begin{methode}[Détermination (sans interpolation linéaire)]
\begin{enumerate}
    \item \textbf{Série discrète} : Calculer les FCC. $Q_1$ = plus petite valeur $x_i$ telle que $F_i \ge 25\%$. $Q_3$ = plus petite $x_i$ telle que $F_i \ge 75\%$. Idem pour $D_k$ (seuil $10k\%$).
    \item \textbf{Série groupée} : \textbf{Lecture graphique} sur le polygone des FCC. Abscisse du point d'ordonnée $25\%, 75\%, 10\%, 90\%$, etc.
\end{enumerate}
\end{methode}

\begin{definition}[Écart interquartile (EIQ)]
$EIQ = Q_3 - Q_1$. Mesure la dispersion du \textbf{centre} de la série (50\% des valeurs centrales). Robuste aux outliers.
\end{definition}

\begin{exemplebox}[Quartiles sur série discrète (Notes sur 20, $N=40$)]
FCC : 5\%, 12,5\%, 27,5\%, 50\%, 70\%, 85\%, 95\%, 100\%.
$Q_1$ : premier FCC $\ge 25\%$ est 27,5\% $\to$ note 8.
$Q_3$ : premier FCC $\ge 75\%$ est 85\% $\to$ note 14.
$EIQ = 6$.
\end{exemplebox}

\coursec{Séries à deux caractères : nuage de points et point moyen}

\begin{definition}[Série à deux caractères quantitatifs]
On observe deux caractères $X$ et $Y$ sur la même population de taille $N$. Données : $(x_1,y_1), \dots, (x_N,y_N)$.
\end{definition}

\begin{definition}[Nuage de points]
Représentation dans un repère orthogonal des points $M_i(x_i\,;\,y_i)$. $X$ en abscisse (variable explicative), $Y$ en ordonnée (variable à expliquer).
\end{definition}

\begin{definition}[Point moyen $G$]
$G(\bar{x}\,;\,\bar{y})$ où $\bar{x} = \frac{\sum x_i}{N},\ \bar{y} = \frac{\sum y_i}{N}$. C'est le barycentre du nuage (même poids $1/N$).
\end{definition}

\begin{propriete}[Liaison affine]
Si les points $M_i$ sont \textbf{proches d'une droite}, on dit qu'il y a une liaison affine entre $X$ et $Y$.
\begin{itemize}
    \item Droite montante ($a>0$) : corrélation positive (ex: révision $\uparrow \Rightarrow$ note $\uparrow$).
    \item Droite descendante ($a<0$) : corrélation négative.
    \item Nuage "circulaire" : pas de liaison affine.
\end{itemize}
\end{propriete}

\coursec{Ajustement affine par la méthode de Mayer et prévisions}

\begin{methode}[Méthode de Mayer (démarche officielle au Bac)]
\begin{enumerate}
    \item Ordonner les points par $x$ croissant.
    \item Partager en \textbf{deux sous-nuages} de même effectif (si $N$ impair, exclure le point central $G$).
    \item Calculer les points moyens $G_1(\bar{x}_1\,;\,\bar{y}_1)$ et $G_2(\bar{x}_2\,;\,\bar{y}_2)$ de chaque sous-nuage.
    \item La \textbf{droite de Mayer} est la droite $(G_1G_2)$.
    \item Équation $y = ax + b$ avec $a = \dfrac{\bar{y}_2 - \bar{y}_1}{\bar{x}_2 - \bar{x}_1}$ et $b = \bar{y}_1 - a\bar{x}_1$ (ou passer par $G$).
\end{enumerate}
\end{methode}

\begin{propriete}[Propriété fondamentale]
La droite de Mayer \textbf{passe toujours par le point moyen $G$} du nuage total.
\emph{Preuve} : $G$ est le barycentre de $G_1$ et $G_2$ avec coefficients $N/2$ et $N/2$, donc milieu de $[G_1G_2]$.
\end{propriete}

\begin{attention}[Prévisions : interpolation vs extrapolation]
\begin{itemize}
    \item \textbf{Interpolation} : prévoir $y$ pour un $x$ \textbf{à l'intérieur} de l'intervalle des $x_i$ observés. Fiable si liaison affine forte.
    \item \textbf{Extrapolation} : prévoir pour $x$ \textbf{en dehors} de l'intervalle observé. \textbf{Très risqué} : la relation linéaire n'est pas garantie hors champ d'observation (ex: pub 100 M FCFA vs max 10 M observés).
\end{itemize}
\end{attention}

\begin{savaistu}[Autres méthodes d'ajustement]
La \textbf{méthode des moindres carrés} (régression linéaire) minimise la somme des carrés des résidus. Elle donne la "meilleure" droite au sens mathématique, mais nécessite calculatrices/logiciels. Mayer est une approximation géométrique simple, sans machine, au programme du Bac camerounais.
\end{savaistu}

\courssub{Je vérifie mes connaissances}

\begin{exercice}[QCM]
Pour chaque question, une seule réponse est exacte. Recopie la lettre correspondante sur ta copie.

\textbf{1.} Dans une série statistique, le \emph{mode} est :
\begin{itemize}
    \item[a)] la valeur qui partage la série en deux effectifs égaux ;
    \item[b)] la valeur (ou la classe) d'effectif maximal ;
    \item[c)] la moyenne des deux valeurs centrales.
\end{itemize}

\textbf{2.} Le point moyen $G$ d'un nuage de points $(x_i\,;\,y_i)$ a pour coordonnées :
\begin{itemize}
    \item[a)] $\left(\dfrac{\sum x_i}{n}\,;\,\dfrac{\sum y_i}{n}\right)$ ;
    \item[b)] $\left(\dfrac{\sum x_i y_i}{n}\,;\,\dfrac{\sum x_i}{n}\right)$ ;
    \item[c)] $\left(\dfrac{x_1+x_n}{2}\,;\,\dfrac{y_1+y_n}{2}\right)$.
\end{itemize}

\textbf{3.} La méthode de Mayer consiste à :
\begin{itemize}
    \item[a)] tracer la droite passant par le premier et le dernier point du nuage ;
    \item[b)] partager le nuage en deux sous-nuages et tracer la droite passant par leurs points moyens ;
    \item[c)] tracer la droite passant par l'origine et le point moyen.
\end{itemize}

\textbf{4.} Dans un diagramme circulaire, une modalité de fréquence $20\,\%$ est représentée par un secteur d'angle :
\begin{itemize}
    \item[a)] $20^{\circ}$ ; \qquad b) $36^{\circ}$ ; \qquad c) $72^{\circ}$.
\end{itemize}
\end{exercice}

\begin{exercice}[Vrai ou faux]
Réponds par \textbf{vrai} ou \textbf{faux} en justifiant brièvement chaque réponse.
\begin{enumerate}
    \item La médiane d'une série statistique est toujours une valeur de la série.
    \item La moyenne d'une série est toujours comprise entre le premier quartile $Q_1$ et le troisième quartile $Q_3$.
    \item Au moins $50\,\%$ des valeurs d'une série sont inférieures ou égales à la médiane.
    \item La droite d'ajustement obtenue par la méthode de Mayer passe toujours par le point moyen du nuage.
    \item Un histogramme est adapté à une série groupée en classes ; un diagramme à bâtons est adapté à une série discrète.
\end{enumerate}
\end{exercice}

\begin{exercice}[Définitions et vocabulaire]
\begin{enumerate}
    \item Définis : population, caractère, effectif, fréquence, effectif cumulé croissant.
    \item Qu'appelle-t-on \emph{centre de classe} d'une classe $[a\,;\,b[$ ? À quoi sert-il dans les calculs ?
    \item Cite les quatre diagrammes usuels étudiés dans cette leçon et précise, pour chacun, le type de série auquel il est le mieux adapté.
    \item Qu'est-ce que le neuvième décile $D_9$ d'une série statistique ?
\end{enumerate}
\end{exercice}

\begin{exercice}[Calculs directs]
On a relevé le nombre de sacs de maïs vendus chaque jour par un commerçant du marché de Mokolo (Yaoundé) pendant $10$ jours :
\[ 12\ ;\ 15\ ;\ 9\ ;\ 12\ ;\ 18\ ;\ 12\ ;\ 10\ ;\ 15\ ;\ 12\ ;\ 20. \]
\begin{enumerate}
    \item Range la série par ordre croissant.
    \item Détermine le mode, la médiane et la moyenne de cette série.
    \item Détermine le premier quartile $Q_1$ et le troisième quartile $Q_3$ (rappelle la méthode utilisée).
\end{enumerate}
\end{exercice}

\courssub{Je m'entraîne}

\begin{exercice}[Diagramme circulaire]
Une enquête menée auprès de $360$ élèves d'un lycée de Douala sur leur opérateur téléphonique principal a donné les résultats suivants :
\begin{center}
\begin{tabularx}{\linewidth}{|l|X|X|X|X|}
\hline
\rowcolor{popblueL}
Opérateur & MTN & Orange & Nexttel & CAMTEL \\
\hline
Effectif & 180 & 108 & 54 & 18 \\
\hline
\end{tabularx}
\end{center}
\begin{enumerate}
    \item Calcule la fréquence (en pourcentage) de chaque modalité.
    \item Calcule l'angle du secteur circulaire correspondant à chaque modalité.
    \item Construis le diagramme circulaire de cette série.
\end{enumerate}
\end{exercice}

\begin{popfigure}
\begin{tikzpicture}
    \pie[sum=auto, radius=3, text=legend, color={popblue, poporange, popgreen, poppurple}]{
        50/MTN,
        30/Orange,
        15/Nexttel,
        5/CAMTEL
    }
\end{tikzpicture}
\legende{Diagramme circulaire des opérateurs téléphoniques (360 élèves).}
\end{popfigure}

\begin{exercice}[Série discrète : notes d'un devoir]
Les notes sur $20$ obtenues par les $40$ élèves d'une classe de Tle A à un devoir de mathématiques sont résumées dans le tableau suivant :
\begin{center}
\begin{tabularx}{\linewidth}{|l|X|X|X|X|X|X|X|X|}
\hline
\rowcolor{popblueL}
Note $x_i$ & 4 & 6 & 8 & 10 & 12 & 14 & 16 & 18 \\
\hline
Effectif $n_i$ & 2 & 3 & 6 & 9 & 8 & 6 & 4 & 2 \\
\hline
\end{tabularx}
\end{center}
\begin{enumerate}
    \item Construis le diagramme à bâtons des effectifs.
    \item Détermine le mode de cette série.
    \item Calcule la moyenne de la classe à ce devoir (arrondis au dixième).
    \item Détermine la médiane, puis les quartiles $Q_1$ et $Q_3$ (méthode des FCC).
    \item Calcule l'écart interquartile et interprète le résultat.
\end{enumerate}
\end{exercice}

\begin{popfigure}
\begin{tikzpicture}
\begin{axis}[
    ybar, bar width=6pt,
    width=\linewidth, height=8cm,
    xlabel={Note $x_i$}, ylabel={Effectif $n_i$},
    xtick={4,6,8,10,12,14,16,18},
    ymin=0, ymax=10,
    enlarge x limits=0.1,
    tick label style={font=\small},
    title={Diagramme à bâtons : Notes de devoir (40 élèves)}
]
\addplot[fill=popblue, draw=popdark] coordinates {
    (4,2) (6,3) (8,6) (10,9) (12,8) (14,6) (16,4) (18,2)
};
\end{axis}
\end{tikzpicture}
\legende{Diagramme à bâtons des effectifs.}
\end{popfigure}

\begin{exercice}[Série groupée en classes]
On a mesuré la distance (en km) parcourue chaque jour par $50$ motos-taxis de la ville de Bafoussam :
\begin{center}
\begin{tabularx}{\linewidth}{|l|X|X|X|X|X|}
\hline
\rowcolor{popblueL}
Distance (km) & $[40\,;\,60[$ & $[60\,;\,80[$ & $[80\,;\,100[$ & $[100\,;\,120[$ & $[120\,;\,140[$ \\
\hline
Effectif & 6 & 12 & 16 & 11 & 5 \\
\hline
\end{tabularx}
\end{center}
\begin{enumerate}
    \item Dresse le tableau complet : centres de classes, fréquences, effectifs cumulés croissants (ECC).
    \item Construis l'histogramme des effectifs (amplitudes égales $\to$ hauteur = effectif).
    \item Calcule la distance moyenne parcourue par jour (arrondis à l'unité).
    \item Détermine la classe modale et la classe médiane.
\end{enumerate}
\end{exercice}

\begin{popfigure}
\begin{tikzpicture}
\begin{axis}[
    ybar interval, 
    width=\linewidth, height=8cm,
    xlabel={Distance (km)}, ylabel={Effectif},
    xtick={40,60,80,100,120,140},
    ymin=0, ymax=18,
    enlarge x limits=0,
    tick label style={font=\small},
    title={Histogramme des distances quotidiennes (50 motos-taxis)}
]
\addplot[fill=poporange!70, draw=popdark] coordinates {
    (40,6) (60,12) (80,16) (100,11) (120,5) (140,0)
};
\end{axis}
\end{tikzpicture}
\legende{Histogramme des effectifs (amplitude constante = 20 km).}
\end{popfigure}

\begin{exercice}[Nuage de points et point moyen]
Le tableau suivant donne le nombre $x$ d'heures de révision hebdomadaires et la note $y$ sur $20$ obtenue par $8$ élèves à un devoir :
\begin{center}
\begin{tabularx}{\linewidth}{|l|X|X|X|X|X|X|X|X|}
\hline
\rowcolor{popblueL}
$x_i$ (heures) & 2 & 3 & 4 & 5 & 6 & 7 & 8 & 9 \\
\hline
$y_i$ (note) & 6 & 7 & 9 & 10 & 11 & 13 & 14 & 16 \\
\hline
\end{tabularx}
\end{center}
\begin{enumerate}
    \item Représente le nuage de points dans un repère orthogonal (unités : \SI{1}{cm} pour $1$ h en abscisse, \SI{1}{cm} pour $2$ points en ordonnée).
    \item Calcule les coordonnées du point moyen $G$ et place-le sur le graphique.
    \item Le nuage justifie-t-il un ajustement affine ? Justifie.
\end{enumerate}
\end{exercice}

\begin{popfigure}
\begin{tikzpicture}
\begin{axis}[
    width=\linewidth, height=10cm,
    xlabel={Heures de révision $x$}, ylabel={Note $y$},
    xmin=1, xmax=10, ymin=4, ymax=18,
    xtick={2,3,4,5,6,7,8,9}, ytick={6,8,10,12,14,16,18},
    grid=major, grid style={popline},
    title={Nuage de points : Révision vs Note},
    scatter/classes={a={mark=*, popblue}, b={mark=*, poporange}},
    legend pos=north west
]
\addplot[only marks, scatter, mark size=2.5] coordinates {
    (2,6) (3,7) (4,9) (5,10) (6,11) (7,13) (8,14) (9,16)
};
\addplot[only marks, mark=*, mark size=3, popred] coordinates {(5.5,10.75)};
\addlegendentry{$M_i(x_i,y_i)$}
\addlegendentry{$G(5,55; 10,75)$}
\end{axis}
\end{tikzpicture}
\legende{Nuage de points et point moyen $G$. Liaison affine évidente.}
\end{popfigure}

\courssub{Je me prépare au Bac}

\begin{exercice}[Type Bac — Tailles des élèves d'un lycée de Bafoussam]
On a mesuré les tailles (en cm) de $100$ élèves de Terminale d'un lycée de Bafoussam. Les résultats sont regroupés dans le tableau suivant :
\begin{center}
\begin{tabularx}{\linewidth}{|l|X|X|X|X|X|X|}
\hline
\rowcolor{popblueL}
Taille (cm) & $[150\,;\,155[$ & $[155\,;\,160[$ & $[160\,;\,165[$ & $[165\,;\,170[$ & $[170\,;\,175[$ & $[175\,;\,180[$ \\
\hline
Effectif & 8 & 18 & 30 & 24 & 14 & 6 \\
\hline
\end{tabularx}
\end{center}
\begin{enumerate}
    \item Recopie et complète le tableau avec les effectifs cumulés croissants (ECC) et les fréquences cumulées croissantes (FCC) exprimées en pourcentage.
    \item Construis l'histogramme des effectifs (unités : \SI{2}{cm} pour $5$ cm de taille en abscisse, \SI{1}{cm} pour $4$ élèves en ordonnée).
    \item Construis, dans un autre repère, le polygone des fréquences cumulées croissantes.
    \item Détermine graphiquement, à l'aide du polygone des FCC, la médiane $M_e$ et les quartiles $Q_1$ et $Q_3$ de la série. Fais apparaître les traits de construction.
    \item Calcule la taille moyenne des élèves (on utilisera les centres de classes et on arrondira au dixième).
    \item Détermine la classe modale de cette série.
    \item Le proviseur affirme : « au moins les trois quarts des élèves mesurent moins de $175$ cm ». A-t-il raison ? Justifie à l'aide du tableau des FCC.
\end{enumerate}
\end{exercice}

\begin{popfigure}
\begin{tikzpicture}
\begin{axis}[
    ybar interval,
    width=\linewidth, height=8cm,
    xlabel={Taille (cm)}, ylabel={Effectif},
    xtick={150,155,160,165,170,175,180},
    ymin=0, ymax=35,
    enlarge x limits=0,
    tick label style={font=\small},
    title={Histogramme des tailles (100 élèves)}
]
\addplot[fill=popgreen!70, draw=popdark] coordinates {
    (150,8) (155,18) (160,30) (165,24) (170,14) (175,6) (180,0)
};
\end{axis}
\end{tikzpicture}
\legende{Histogramme des tailles.}
\end{popfigure}

\begin{popfigure}
\begin{tikzpicture}
\begin{axis}[
    width=\linewidth, height=8cm,
    xlabel={Taille (cm)}, ylabel={FCC (\%)},
    xmin=145, xmax=185, ymin=0, ymax=105,
    xtick={150,155,160,165,170,175,180},
    ytick={0,25,50,75,100},
    grid=major, grid style={popline},
    title={Polygone des fréquences cumulées croissantes},
    clip=false
]
% Points: (borne sup, FCC)
\addplot[popcurve, mark=*, mark size=2] coordinates {
    (150,0) (155,8) (160,26) (165,56) (170,80) (175,94) (180,100)
};
% Construction lines for Median (50%), Q1 (25%), Q3 (75%)
\draw[dashed, popred, thick] (145,50) -- (163.5,50) node[pos=0, left] {$50\%$};
\draw[dashed, popred, thick] (163.5,50) -- (163.5,0) node[below] {$M_e \approx 163,5$};
\draw[dashed, popblue, thick] (145,25) -- (159,25) node[pos=0, left] {$25\%$};
\draw[dashed, popblue, thick] (159,25) -- (159,0) node[below] {$Q_1 \approx 159$};
\draw[dashed, poporange, thick] (145,75) -- (169,75) node[pos=0, left] {$75\%$};
\draw[dashed, poporange, thick] (169,75) -- (169,0) node[below] {$Q_3 \approx 169$};
\end{axis}
\end{tikzpicture}
\legende{Polygone des FCC avec lecture graphique de $M_e, Q_1, Q_3$.}
\end{popfigure}

\begin{exercice}[Type Bac — Buts du championnat MTN Elite One]
Lors des $22$ journées d'une saison du championnat MTN Elite One, on a relevé le nombre de buts marqués par journée :
\begin{center}
\begin{tabularx}{\linewidth}{|l|X|X|X|X|X|X|X|}
\hline
\rowcolor{popblueL}
Nombre de buts $x_i$ & 0 & 1 & 2 & 3 & 4 & 5 & 6 \\
\hline
Nombre de journées $n_i$ & 1 & 3 & 5 & 6 & 4 & 2 & 1 \\
\hline
\end{tabularx}
\end{center}
\begin{enumerate}
    \item Construis le diagramme à bâtons des effectifs.
    \item Détermine le mode de cette série et interprète le résultat.
    \item Calcule le nombre moyen de buts marqués par journée (arrondis au centième).
    \item Détermine la médiane de la série, puis les déciles $D_1$ et $D_9$ (méthode des FCC).
    \item Interprète concrètement la valeur de $D_9$ dans le contexte du championnat.
    \item Un journaliste sportif écrit : « la moitié des journées a vu au moins $4$ buts ». Cette affirmation est-elle exacte ? Justifie à l'aide des effectifs cumulés.
\end{enumerate}
\end{exercice}

\begin{popfigure}
\begin{tikzpicture}
\begin{axis}[
    ybar, bar width=8pt,
    width=\linewidth, height=8cm,
    xlabel={Nombre de buts $x_i$}, ylabel={Nombre de journées $n_i$},
    xtick={0,1,2,3,4,5,6},
    ymin=0, ymax=7,
    enlarge x limits=0.1,
    tick label style={font=\small},
    title={Buts par journée (MTN Elite One, 22 journées)}
]
\addplot[fill=poppurple!70, draw=popdark] coordinates {
    (0,1) (1,3) (2,5) (3,6) (4,4) (5,2) (6,1)
};
\end{axis}
\end{tikzpicture}
\legende{Diagramme à bâtons des buts par journée.}
\end{popfigure}

\begin{exercice}[Type Bac — Publicité et chiffre d'affaires]
Une entreprise de distribution de boissons basée à Douala a relevé, pendant $6$ trimestres, ses dépenses de publicité $x_i$ (en millions de FCFA) et son chiffre d'affaires $y_i$ (en centaines de millions de FCFA) :
\begin{center}
\begin{tabularx}{\linewidth}{|l|X|X|X|X|X|X|}
\hline
\rowcolor{popblueL}
Trimestre & 1 & 2 & 3 & 4 & 5 & 6 \\
\hline
Dépenses $x_i$ & 2 & 3 & 5 & 6 & 8 & 10 \\
\hline
Chiffre d'affaires $y_i$ & 12 & 15 & 19 & 22 & 27 & 31 \\
\hline
\end{tabularx}
\end{center}
\begin{enumerate}
    \item Représente le nuage de points $M_i(x_i\,;\,y_i)$ dans un repère orthogonal (unités : \SI{1}{cm} pour $1$ million de FCFA en abscisse, \SI{1}{cm} pour $2$ centaines de millions de FCFA en ordonnée).
    \item Calcule les coordonnées du point moyen $G$ du nuage et place-le sur le graphique.
    \item On partage le nuage en deux sous-nuages : les trois premiers points et les trois derniers.
    \begin{enumerate}
        \item Calcule les coordonnées des points moyens $G_1$ et $G_2$ de ces deux sous-nuages.
        \item Détermine une équation de la droite $(G_1G_2)$, droite d'ajustement de Mayer du nuage, sous la forme $y = ax + b$ (arrondis $a$ et $b$ au centième).
        \item Trace cette droite sur le graphique. Vérifie qu'elle passe par le point moyen $G$.
    \end{enumerate}
    \item L'entreprise prévoit d'investir $12$ millions de FCFA en publicité le trimestre prochain. Estime, à l'aide de la droite de Mayer, le chiffre d'affaires attendu.
    \item Le directeur utilise la même droite pour prévoir le chiffre d'affaires correspondant à des dépenses de $100$ millions de FCFA et trouve un résultat qu'il juge « miraculeux ». Explique pourquoi une telle prévision n'est pas crédible.
\end{enumerate}
\end{exercice}

\begin{popfigure}
\begin{tikzpicture}
\begin{axis}[
    width=\linewidth, height=10cm,
    xlabel={Dépenses pub. $x$ (M FCFA)}, ylabel={Chiffre d'affaires $y$ (100 M FCFA)},
    xmin=0, xmax=13, ymin=10, ymax=35,
    xtick={2,3,4,5,6,7,8,9,10,11,12}, ytick={12,15,18,21,24,27,30,33},
    grid=major, grid style={popline},
    title={Nuage de points, point moyen $G$ et droite de Mayer},
    legend pos=north west
]
% Nuage
\addplot[only marks, mark=*, mark size=2.5, popblue] coordinates {
    (2,12) (3,15) (5,19) (6,22) (8,27) (10,31)
};
% Point moyen G
\addplot[only marks, mark=*, mark size=3, popred] coordinates {(5.667,21)};
% Mayer line: y = 2.43x + 7.24 (approx)
\addplot[popcurve, domain=0:12, samples=2, thick, poporange] {2.43*x + 7.24};
\addlegendentry{$M_i(x_i,y_i)$}
\addlegendentry{$G(5,67; 21)$}
\addlegendentry{Droite de Mayer}
% G1 G2
\addplot[only marks, mark=square*, mark size=3, popgreen] coordinates {(3.333,15.333)};
\addplot[only marks, mark=square*, mark size=3, poppurple] coordinates {(8,26.667)};
\addlegendentry{$G_1, G_2$}
\end{axis}
\end{tikzpicture}
\legende{Ajustement affine par la méthode de Mayer.}
\end{popfigure}

\newpage

\coursec{Corrigés des exercices}

\begin{corrige}[Exercice 1 — QCM]
\textbf{1. Réponse b).} Le mode est la valeur (ou la classe) d'\cle{effectif maximal}. La réponse a) définit la médiane ; la réponse c) décrit le calcul de la médiane dans le cas d'un effectif pair.

\textbf{2. Réponse a).} Le point moyen est le barycentre du nuage affecté de poids égaux : $G\left(\bar{x}\,;\,\bar{y}\right)$ avec $\bar{x}=\dfrac{\sum x_i}{n}$ et $\bar{y}=\dfrac{\sum y_i}{n}$.

\textbf{3. Réponse b).} La méthode de Mayer consiste à partager le nuage en deux sous-nuages de même effectif (ou d'effectifs les plus proches possibles), à calculer leurs points moyens $G_1$ et $G_2$, puis à tracer la droite $(G_1G_2)$. La réponse a) ne tient pas compte de la répartition des points ; la réponse c) impose un passage par l'origine sans justification.

\textbf{4. Réponse c).} L'angle du secteur est $\alpha = 360^{\circ}\times f = 360^{\circ}\times 0{,}20 = 72^{\circ}$.
\end{corrige}

\begin{corrige}[Exercice 2 — Vrai ou faux]
\begin{enumerate}
    \item \textbf{Faux.} Lorsque l'effectif total $N$ est pair, la médiane est la moyenne des valeurs aux rangs $\frac{N}{2}$ et $\frac{N}{2}+1$, qui peut ne pas être une valeur de la série. Contre-exemple : la série $1\ ;\ 2\ ;\ 3\ ;\ 10$ a pour médiane $\frac{2+3}{2}=2{,}5$, qui n'est pas une valeur de la série.
    \item \textbf{Faux.} La moyenne est sensible aux valeurs extrêmes et peut sortir de l'intervalle interquartile. Contre-exemple : la série $1\ ;\ 2\ ;\ 3\ ;\ 4\ ;\ 100$ a pour moyenne $\bar{x}=22$, alors que $Q_3=4$.
    \item \textbf{Vrai.} C'est la définition même de la médiane : au moins $50\,\%$ des valeurs sont inférieures ou égales à $M_e$ (et au moins $50\,\%$ lui sont supérieures ou égales).
    \item \textbf{Vrai.} $G$ est le barycentre de $G_1$ et $G_2$ affectés des coefficients $\frac{N}{2}$ et $\frac{N}{2}$ ; c'est donc le milieu du segment $[G_1G_2]$, qui appartient à la droite $(G_1G_2)$.
    \item \textbf{Vrai.} L'histogramme (rectangles juxtaposés, aire proportionnelle à l'effectif) est adapté à un caractère continu groupé en classes ; le diagramme à bâtons (bâtons séparés) est adapté à un caractère discret.
\end{enumerate}
\end{corrige}

\begin{corrige}[Exercice 3 — Définitions et vocabulaire]
\begin{enumerate}
    \item \textbf{Population} : ensemble des êtres ou objets sur lesquels porte l'étude statistique. \textbf{Caractère} : propriété observée sur chaque individu de la population (qualitatif, quantitatif discret ou continu). \textbf{Effectif} $n_i$ : nombre d'individus prenant la modalité $x_i$. \textbf{Fréquence} $f_i=\dfrac{n_i}{N}$ : part de la modalité dans l'effectif total, souvent exprimée en pourcentage. \textbf{Effectif cumulé croissant} $N_i$ : somme des effectifs des modalités inférieures ou égales à $x_i$.
    \item Le \textbf{centre de classe} de $[a\,;\,b[$ est $c=\dfrac{a+b}{2}$. Dans une série groupée, les valeurs exactes sont inconnues : on représente chaque classe par son centre pour calculer la moyenne ($\bar{x}=\frac{\sum n_i c_i}{N}$) et construire le diagramme à bâtons des centres.
    \item \textbf{Diagramme circulaire} : caractère qualitatif ou discret à peu de modalités. \textbf{Diagramme à bâtons} : caractère quantitatif discret. \textbf{Histogramme} : caractère quantitatif continu groupé en classes. \textbf{Polygone des fréquences cumulées croissantes} : série groupée en classes, pour la lecture graphique de la médiane et des quartiles.
    \item Le neuvième décile $D_9$ est la plus petite valeur de la série telle qu'au moins $90\,\%$ des valeurs lui soient inférieures ou égales (et donc au moins $10\,\%$ lui sont supérieures ou égales).
\end{enumerate}
\end{corrige}

\begin{corrige}[Exercice 4 — Calculs directs]
\begin{enumerate}
    \item Série rangée par ordre croissant :
    \[ 9\ ;\ 10\ ;\ 12\ ;\ 12\ ;\ 12\ ;\ 12\ ;\ 15\ ;\ 15\ ;\ 18\ ;\ 20. \]
    \item \textbf{Mode} : la valeur $12$ apparaît $4$ fois, c'est l'effectif maximal. Le mode est $\boxed{12\text{ sacs}}$.
    
    \textbf{Médiane} : $N=10$ est pair. La médiane est la moyenne des valeurs aux rangs $\frac{10}{2}=5$ et $6$ : $M_e=\dfrac{12+12}{2}=\boxed{12\text{ sacs}}$.
    
    \textbf{Moyenne} :
    \[ \bar{x}=\frac{9+10+12+12+12+12+15+15+18+20}{10}=\frac{135}{10}=\boxed{13{,}5\text{ sacs}}. \]
    \item \textbf{Méthode des FCC} : les fréquences cumulées croissantes sont : $10\,\%$ (9), $20\,\%$ (10), $60\,\%$ (12), $80\,\%$ (15), $90\,\%$ (18), $100\,\%$ (20).
    
    $Q_1$ est la plus petite valeur dont la FCC est $\ge 25\,\%$ : la première FCC atteignant ou dépassant $25\,\%$ est $60\,\%$, donc $\boxed{Q_1=12}$.
    
    $Q_3$ est la plus petite valeur dont la FCC est $\ge 75\,\%$ : la première FCC $\ge 75\,\%$ est $80\,\%$, donc $\boxed{Q_3=15}$.
\end{enumerate}
\end{corrige}

\begin{corrige}[Exercice 5 — Diagramme circulaire]
\begin{enumerate}
    \item \textbf{Fréquences} : $N=360$.
    \[ f_{\text{MTN}}=\frac{180}{360}=50\,\%\ ;\quad f_{\text{Orange}}=\frac{108}{360}=30\,\%\ ;\quad f_{\text{Nexttel}}=\frac{54}{360}=15\,\%\ ;\quad f_{\text{CAMTEL}}=\frac{18}{360}=5\,\%. \]
    Vérification : $50+30+15+5=100\,\%$.
    \item \textbf{Angles} : $\alpha = 360^{\circ}\times f$.
    \[ \alpha_{\text{MTN}}=180^{\circ}\ ;\quad \alpha_{\text{Orange}}=108^{\circ}\ ;\quad \alpha_{\text{Nexttel}}=54^{\circ}\ ;\quad \alpha_{\text{CAMTEL}}=18^{\circ}. \]
    Vérification : $180+108+54+18=360^{\circ}$.
    \item \textbf{Construction} : tracer un cercle, puis reporter au rapporteur les secteurs de $180^{\circ}$ (un demi-disque pour MTN), $108^{\circ}$, $54^{\circ}$ et $18^{\circ}$, en les coloriant différemment et en ajoutant une légende. Le diagramme obtenu est celui de la figure de l'énoncé.
\end{enumerate}
\end{corrige}

\begin{corrige}[Exercice 6 — Série discrète : notes d'un devoir]
\begin{enumerate}
    \item Le diagramme à bâtons est construit en plaçant, pour chaque note $x_i$, un bâton vertical de hauteur $n_i$ (voir figure de l'énoncé) : hauteurs $2, 3, 6, 9, 8, 6, 4, 2$ aux abscisses $4, 6, 8, 10, 12, 14, 16, 18$.
    \item \textbf{Mode} : l'effectif maximal est $9$, atteint pour la note $10$. Le mode est $\boxed{10}$.
    \item \textbf{Moyenne} :
    \begin{align*}
    \sum n_i x_i &= 4\times 2+6\times 3+8\times 6+10\times 9+12\times 8+14\times 6+16\times 4+18\times 2\\
    &= 8+18+48+90+96+84+64+36 = 444.
    \end{align*}
    \[ \bar{x}=\frac{444}{40}=11{,}1. \]
    La moyenne de la classe est $\boxed{11{,}1}$ sur $20$.
    \item \textbf{ECC et FCC} :
    \begin{center}
    \begin{tabularx}{\linewidth}{|l|X|X|X|X|X|X|X|X|}
    \hline
    \rowcolor{popblueL}
    $x_i$ & 4 & 6 & 8 & 10 & 12 & 14 & 16 & 18 \\
    \hline
    ECC & 2 & 5 & 11 & 20 & 28 & 34 & 38 & 40 \\
    \hline
    FCC (\%) & 5 & 12,5 & 27,5 & 50 & 70 & 85 & 95 & 100 \\
    \hline
    \end{tabularx}
    \end{center}
    \textbf{Médiane} : $N=40$ est pair ; on prend la moyenne des valeurs aux rangs $20$ et $21$. La $20^{\text{e}}$ valeur est $10$ (ECC $=20$) et la $21^{\text{e}}$ est $12$. Donc $M_e=\dfrac{10+12}{2}=\boxed{11}$.
    
    \textbf{Quartiles} : $Q_1$ est la plus petite note de FCC $\ge 25\,\%$ : c'est $\boxed{Q_1=8}$ (FCC $=27{,}5\,\%$). $Q_3$ est la plus petite note de FCC $\ge 75\,\%$ : c'est $\boxed{Q_3=14}$ (FCC $=85\,\%$).
    \item $EIQ = Q_3-Q_1 = 14-8 = \boxed{6}$. Interprétation : les $50\,\%$ centraux des élèves ont des notes comprises entre $8$ et $14$, sur une plage de $6$ points ; la dispersion centrale est modérée.
\end{enumerate}
\end{corrige}

\begin{corrige}[Exercice 7 — Série groupée en classes]
\begin{enumerate}
    \item \textbf{Tableau complet} :
    \begin{center}
    \begin{tabularx}{\linewidth}{|l|X|X|X|X|X|}
    \hline
    \rowcolor{popblueL}
    Classe & $[40;60[$ & $[60;80[$ & $[80;100[$ & $[100;120[$ & $[120;140[$ \\
    \hline
    Centre $c_i$ & 50 & 70 & 90 & 110 & 130 \\
    \hline
    Effectif $n_i$ & 6 & 12 & 16 & 11 & 5 \\
    \hline
    Fréquence (\%) & 12 & 24 & 32 & 22 & 10 \\
    \hline
    ECC & 6 & 18 & 34 & 45 & 50 \\
    \hline
    \end{tabularx}
    \end{center}
    \item Les classes ont toutes la même amplitude ($20$ km), donc la hauteur de chaque rectangle est égale à l'effectif : rectangles de hauteurs $6, 12, 16, 11, 5$ juxtaposés sur les intervalles successifs (voir figure de l'énoncé).
    \item \textbf{Distance moyenne} (avec les centres de classes) :
    \begin{align*}
    \sum n_i c_i &= 6\times 50+12\times 70+16\times 90+11\times 110+5\times 130\\
    &= 300+840+1440+1210+650 = 4440.
    \end{align*}
    \[ \bar{x}=\frac{4440}{50}=88{,}8 \approx \boxed{89\text{ km}}. \]
    \item \textbf{Classe modale} : effectif maximal $16$, donc la classe modale est $\boxed{[80\,;\,100[}$.
    
    \textbf{Classe médiane} : $\frac{N}{2}=25$ ; le premier ECC $\ge 25$ est $34$, correspondant à la classe $\boxed{[80\,;\,100[}$.
\end{enumerate}
\end{corrige}

\begin{corrige}[Exercice 8 — Nuage de points et point moyen]
\begin{enumerate}
    \item On place les huit points $M_1(2\,;\,6)$, $M_2(3\,;\,7)$, $M_3(4\,;\,9)$, $M_4(5\,;\,10)$, $M_5(6\,;\,11)$, $M_6(7\,;\,13)$, $M_7(8\,;\,14)$, $M_8(9\,;\,16)$ dans le repère (voir figure de l'énoncé).
    \item \textbf{Point moyen} :
    \[ \bar{x}=\frac{2+3+4+5+6+7+8+9}{8}=\frac{44}{8}=5{,}5 \ ;\quad \bar{y}=\frac{6+7+9+10+11+13+14+16}{8}=\frac{86}{8}=10{,}75. \]
    Donc $\boxed{G(5{,}5\,;\,10{,}75)}$. On le place sur le graphique avec un symbole distinct.
    \item \textbf{Oui}, un ajustement affine est justifié : les points du nuage sont sensiblement alignés autour d'une droite montante. Il existe donc une liaison affine (corrélation positive) entre le nombre d'heures de révision et la note obtenue.
\end{enumerate}
\end{corrige}

\begin{corrige}[Exercice 9 — Type Bac : Tailles des élèves d'un lycée de Bafoussam]
\begin{enumerate}
    \item \textbf{Tableau complété} :
    \begin{center}
    \begin{tabularx}{\linewidth}{|l|X|X|X|X|X|X|}
    \hline
    \rowcolor{popblueL}
    Taille (cm) & $[150;155[$ & $[155;160[$ & $[160;165[$ & $[165;170[$ & $[170;175[$ & $[175;180[$ \\
    \hline
    Effectif & 8 & 18 & 30 & 24 & 14 & 6 \\
    \hline
    ECC & 8 & 26 & 56 & 80 & 94 & 100 \\
    \hline
    FCC (\%) & 8 & 26 & 56 & 80 & 94 & 100 \\
    \hline
    \end{tabularx}
    \end{center}
    \item \textbf{Histogramme} : amplitudes égales ($5$ cm), donc hauteurs = effectifs : $8, 18, 30, 24, 14, 6$. Rectangles juxtaposés, sans espace (voir figure de l'énoncé).
    \item \textbf{Polygone des FCC} : on place les points $(150\,;\,0)$, $(155\,;\,8)$, $(160\,;\,26)$, $(165\,;\,56)$, $(170\,;\,80)$, $(175\,;\,94)$, $(180\,;\,100)$ — abscisses = bornes \emph{supérieures} des classes — puis on les relie par des segments (voir figure de l'énoncé).
    \item \textbf{Lecture graphique} : on trace les horizontales d'ordonnées $25\,\%$, $50\,\%$, $75\,\%$ jusqu'au polygone, puis on abaisse les verticales sur l'axe des abscisses :
    \[ Q_1 \approx 159\text{ cm}\ ;\quad M_e \approx 163{,}5\text{ cm}\ ;\quad Q_3 \approx 169\text{ cm}. \]
    (Toute valeur lue dans un intervalle de tolérance raisonnable, par exemple $\pm 1$ cm, est acceptée.)
    \item \textbf{Taille moyenne} (centres de classes : $152{,}5\,;\,157{,}5\,;\,162{,}5\,;\,167{,}5\,;\,172{,}5\,;\,177{,}5$) :
    \begin{align*}
    \sum n_i c_i &= 8\times 152{,}5+18\times 157{,}5+30\times 162{,}5+24\times 167{,}5+14\times 172{,}5+6\times 177{,}5\\
    &= 1220+2835+4875+4020+2415+1065 = 16430.
    \end{align*}
    \[ \bar{x}=\frac{16430}{100}=\boxed{164{,}3\text{ cm}}. \]
    \item \textbf{Classe modale} : effectif maximal $30$, donc la classe modale est $\boxed{[160\,;\,165[}$.
    \item La FCC à la borne $175$ vaut $94\,\%$ : $94\,\%$ des élèves mesurent moins de $175$ cm. Or $94\,\% > 75\,\%$ (les trois quarts). \textbf{Le proviseur a donc raison.}
\end{enumerate}

\textbf{Barème indicatif (sur 10)} : tableau complété 1,5 pt ; histogramme 1,5 pt ; polygone des FCC 2 pts ; lectures graphiques avec traits de construction 2 pts ; moyenne 1,5 pt ; classe modale 0,5 pt ; justification finale 1 pt.

\textbf{Points de vigilance} :
\begin{itemize}
    \item Pour le polygone des FCC, les abscisses sont les \textbf{bornes supérieures} des classes, et on n'oublie pas le point $(150\,;\,0)$.
    \item Les traits de construction (pointillés horizontaux et verticaux) doivent \textbf{apparaître sur la copie} : une valeur donnée sans trait de lecture n'est pas justifiée.
    \item Pour la moyenne d'une série groupée, citer explicitement l'utilisation des \textbf{centres de classes}.
    \item L'interpolation linéaire étant hors programme, une valeur \emph{lue graphiquement} (avec tolérance) est attendue, pas une formule.
\end{itemize}
\end{corrige}

\begin{corrige}[Exercice 10 — Type Bac : Buts du championnat MTN Elite One]
\begin{enumerate}
    \item \textbf{Diagramme à bâtons} : bâtons de hauteurs $1, 3, 5, 6, 4, 2, 1$ aux abscisses $0, 1, 2, 3, 4, 5, 6$ (voir figure de l'énoncé).
    \item \textbf{Mode} : effectif maximal $6$ pour $x_i=3$. Le mode est $\boxed{3\text{ buts}}$. Interprétation : le nombre de buts le plus fréquemment observé par journée est $3$.
    \item \textbf{Moyenne} :
    \[ \sum n_i x_i = 0\times 1+1\times 3+2\times 5+3\times 6+4\times 4+5\times 2+6\times 1 = 0+3+10+18+16+10+6 = 63. \]
    \[ \bar{x}=\frac{63}{22}\approx \boxed{2{,}86\text{ buts par journée}}. \]
    \item \textbf{ECC et FCC} :
    \begin{center}
    \begin{tabularx}{\linewidth}{|l|X|X|X|X|X|X|X|}
    \hline
    \rowcolor{popblueL}
    $x_i$ & 0 & 1 & 2 & 3 & 4 & 5 & 6 \\
    \hline
    ECC & 1 & 4 & 9 & 15 & 19 & 21 & 22 \\
    \hline
    FCC (\%) & 4,5 & 18,2 & 40,9 & 68,2 & 86,4 & 95,5 & 100 \\
    \hline
    \end{tabularx}
    \end{center}
    \textbf{Médiane} : $N=22$ est pair ; moyenne des valeurs aux rangs $11$ et $12$. Ces deux valeurs sont égales à $3$ (ECC : $9 < 11 \le 15$), donc $\boxed{M_e=3}$.
    
    \textbf{Décile $D_1$} : plus petite valeur de FCC $\ge 10\,\%$ : la FCC $18{,}2\,\%$ est la première à dépasser $10\,\%$, donc $\boxed{D_1=1}$.
    
    \textbf{Décile $D_9$} : plus petite valeur de FCC $\ge 90\,\%$ : la FCC $95{,}5\,\%$ est la première à dépasser $90\,\%$, donc $\boxed{D_9=5}$.
    \item \textbf{Interprétation de $D_9=5$} : au moins $90\,\%$ des journées ont vu au plus $5$ buts ; autrement dit, seulement environ $10\,\%$ des journées (ici $1$ journée sur $22$) ont vu $6$ buts ou plus.
    \item Les journées avec \textbf{au moins $4$ buts} sont celles avec $4$, $5$ ou $6$ buts : $4+2+1=7$ journées, soit $\frac{7}{22}\approx 31{,}8\,\%$. (De manière équivalente : l'ECC à $3$ buts est $15$, donc $22-15=7$ journées dépassent $3$ buts.) Comme $7 \neq 11$ (la moitié de $22$), \textbf{l'affirmation du journaliste est fausse}.
\end{enumerate}

\textbf{Barème indicatif (sur 10)} : diagramme à bâtons 1,5 pt ; mode + interprétation 1,5 pt ; moyenne 1,5 pt ; médiane 1,5 pt ; déciles 2 pts ; interprétation de $D_9$ 1 pt ; vérification de l'affirmation 1 pt.

\textbf{Points de vigilance} :
\begin{itemize}
    \item Bien distinguer « au moins $4$ buts » ($\ge 4$) de « plus de $4$ buts » ($>4$) : c'est l'erreur classique dans ce type de question.
    \item Pour les déciles d'une série discrète, la méthode exigée est celle des \textbf{FCC} : citer le seuil ($10\,\%$, $90\,\%$) et la première FCC qui l'atteint ou le dépasse.
    \item Une interprétation \emph{en contexte} (journées, buts) est attendue : une réponse purement numérique est incomplète.
\end{itemize}
\end{corrige}

\begin{corrige}[Exercice 11 — Type Bac : Publicité et chiffre d'affaires]
\begin{enumerate}
    \item On place les points $M_1(2\,;\,12)$, $M_2(3\,;\,15)$, $M_3(5\,;\,19)$, $M_4(6\,;\,22)$, $M_5(8\,;\,27)$, $M_6(10\,;\,31)$ dans le repère (voir figure de l'énoncé).
    \item \textbf{Point moyen} :
    \[ \bar{x}=\frac{2+3+5+6+8+10}{6}=\frac{34}{6}=\frac{17}{3}\approx 5{,}67\ ;\quad \bar{y}=\frac{12+15+19+22+27+31}{6}=\frac{126}{6}=21. \]
    Donc $\boxed{G\left(\frac{17}{3}\,;\,21\right)\approx G(5{,}67\,;\,21)}$.
    \item \begin{enumerate}
        \item \textbf{Premier sous-nuage} $(M_1, M_2, M_3)$ :
        \[ \bar{x}_1=\frac{2+3+5}{3}=\frac{10}{3}\approx 3{,}33\ ;\quad \bar{y}_1=\frac{12+15+19}{3}=\frac{46}{3}\approx 15{,}33. \]
        $\boxed{G_1\left(\frac{10}{3}\,;\,\frac{46}{3}\right)}$.
        
        \textbf{Second sous-nuage} $(M_4, M_5, M_6)$ :
        \[ \bar{x}_2=\frac{6+8+10}{3}=8\ ;\quad \bar{y}_2=\frac{22+27+31}{3}=\frac{80}{3}\approx 26{,}67. \]
        $\boxed{G_2\left(8\,;\,\frac{80}{3}\right)}$.
        \item \textbf{Équation de la droite $(G_1G_2)$} :
        \[ a=\frac{\bar{y}_2-\bar{y}_1}{\bar{x}_2-\bar{x}_1}=\frac{\frac{80}{3}-\frac{46}{3}}{8-\frac{10}{3}}=\frac{\frac{34}{3}}{\frac{14}{3}}=\frac{34}{14}=\frac{17}{7}\approx 2{,}43. \]
        \[ b=\bar{y}_1-a\bar{x}_1=\frac{46}{3}-\frac{17}{7}\times\frac{10}{3}=\frac{46}{3}-\frac{170}{21}=\frac{322-170}{21}=\frac{152}{21}\approx 7{,}24. \]
        La droite de Mayer a pour équation :
        \[ \boxed{y = 2{,}43\,x + 7{,}24}. \]
        \item \textbf{Vérification du passage par $G$} : $a\bar{x}+b = \dfrac{17}{7}\times\dfrac{17}{3}+\dfrac{152}{21}=\dfrac{289}{21}+\dfrac{152}{21}=\dfrac{441}{21}=21=\bar{y}$. La droite passe bien par $G$.
    \end{enumerate}
    \item \textbf{Prévision pour $x=12$} :
    \[ y = 2{,}43\times 12+7{,}24 = 29{,}16+7{,}24 = 36{,}40. \]
    Le chiffre d'affaires attendu est d'environ $36{,}4$ centaines de millions de FCFA, soit $\boxed{\text{environ }3{,}64\text{ milliards de FCFA}}$. Il s'agit d'une \emph{extrapolation proche} ($12$ est légèrement au-delà du maximum observé $10$) : elle reste utilisable avec prudence.
    \item Pour $x=100$, on obtiendrait $y\approx 2{,}43\times 100+7{,}24\approx 250$, soit $25$ milliards de FCFA. Cette prévision n'est \textbf{pas crédible} car il s'agit d'une \textbf{extrapolation très éloignée} du champ d'observation ($x\in[2\,;\,10]$) : rien ne garantit que la relation affine reste valable pour des dépenses dix fois supérieures (saturation du marché, concurrence, capacité de production limitée, etc.). La droite de Mayer n'est un modèle fiable qu'à l'intérieur (ou au voisinage immédiat) de l'intervalle des valeurs observées.
\end{enumerate}

\textbf{Barème indicatif (sur 10)} : nuage de points 1,5 pt ; point moyen $G$ 1 pt ; points $G_1$ et $G_2$ 2 pts ; équation de la droite 2 pts ; tracé + vérification du passage par $G$ 1 pt ; prévision 1,5 pt ; critique de l'extrapolation 1 pt.

\textbf{Points de vigilance} :
\begin{itemize}
    \item Le partage du nuage se fait sur des points \textbf{ordonnés par $x$ croissant} (ici déjà le cas) ; avec $N=6$ pair, les deux sous-nuages ont chacun $3$ points.
    \item La vérification $G\in(G_1G_2)$ doit être \textbf{rédigée} (calcul de $a\bar{x}+b$ comparé à $\bar{y}$), pas seulement constatée graphiquement.
    \item Pour la prévision, préciser les \textbf{unités} : $y$ est exprimé en centaines de millions de FCFA ; la conversion finale est attendue.
    \item La question 5 exige le vocabulaire du cours : \emph{extrapolation}, \emph{champ d'observation}, \emph{modèle non garanti hors intervalle}.
\end{itemize}
\end{corrige}

\newpage

\coursec{Fiche bilan}

\coursec{Statistiques : Organisation et gestion des données}

\begin{popfigure}
\begin{center}
\begin{tikzpicture}[mindmap, grow cyclic, every node/.style={concept}, root concept/.append style={concept color=popblue, font=\large\bfseries, text=white}, level 1/.append style={level distance=3.4cm, sibling angle=60}, level 1 concept/.append style={font=\small\bfseries, text=white}]
\node[root concept]{Statistiques}
  child[concept color=poporange]{ node[align=center] {1. Vocabulaire\\ et organisation} }
  child[concept color=popgreen]{ node[align=center] {2. Représentations\\ graphiques} }
  child[concept color=poppurple]{ node[align=center] {3. Position :\\ Mode, Moyenne, Médiane} }
  child[concept color=popgold]{ node[align=center] {4. Dispersion :\\ Quartiles, Déciles} }
  child[concept color=poppink]{ node[align=center] {5. Nuage de points\\ Point moyen} }
  child[concept color=popink]{ node[align=center] {6. Ajustement affine\\ Méthode de Mayer} };
\end{tikzpicture}
\end{center}
\legende{Carte mentale de la leçon : les six grandes idées à maîtriser pour le Baccalauréat.}
\end{popfigure}

\noindent Bienvenue dans cette leçon essentielle de Terminale. Les statistiques sont l'outil indispensable pour comprendre le monde qui nous entoure : évolution du prix du kilogramme de riz sur les marchés de Douala et Yaoundé, consommation d'électricité des ménages abonnés à ENEO, nombre d'abonnés MTN versus Orange, rendements agricoles du cacao dans le Sud-Ouest... Savoir organiser, représenter, synthétiser et prévoir à partir de données chiffrées est une compétence majeure, évaluée au Baccalauréat et utile pour toute carrière. Nous allons progresser pas à pas, du vocabulaire de base jusqu'à la prévision par la méthode de Mayer.

\coursec{Vocabulaire et organisation des données}

\courssub{Population, individu, caractère}

\begin{definition}[Population statistique]
On appelle \cle{population statistique} l'ensemble des êtres ou des objets sur lesquels porte l'étude. Chaque élément de cette population est un \cle{individu} (ou unité statistique).
\end{definition}

\begin{definition}[Caractère statistique]
Le \cle{caractère} étudié est la propriété observée sur chaque individu. On distingue :
\begin{itemize}
    \item Le \cle{caractère qualitatif} : il exprime une qualité, une catégorie (ex: la marque de téléphone \og MTN \fg, \og Orange \fg, \og Camtel \fg ; le type de carburant \og Super \fg, \og Gasoil \fg).
    \item Le \cle{caractère quantitatif} : il exprime une quantité, un nombre.
    \begin{itemize}
        \item \textbf{Discret} : il ne prend que des valeurs isolées, souvent entières (ex: nombre d'enfants par ménage, nombre de pannes sur le réseau CAMTEL par mois).
        \item \textbf{Continu} : il peut prendre toute valeur dans un intervalle (ex: taille en cm, poids en kg, facture d'électricité en FCFA, temps de trajet Douala--Yaoundé en minutes).
    \end{itemize}
\end{itemize}
\end{definition}

\begin{attention}[Piège fréquent]
Ne confondez pas \og caractère quantitatif discret \fg (valeurs comptables, sauts) et \og caractère quantitatif continu \fg (valeurs mesurables, continuum). La distinction guide le choix du diagramme : bâtons/histogramme pour le discret, histogramme/polygone pour le continu.
\end{attention}

\courssub{Effectifs, fréquences, tableau de données}

Soit une série statistique à une variable $X$ prenant les valeurs $x_1, x_2, \dots, x_p$ (pour un caractère discret ou un caractère continu groupé en classes).

\begin{definition}[Effectif, Effectif total, Fréquence]
\begin{itemize}
    \item L'\cle{effectif} $n_i$ de la valeur $x_i$ est le nombre d'individus pour lesquels $X = x_i$.
    \item L'\cle{effectif total} est $N = \sum_{i=1}^{p} n_i$.
    \item La \cle{fréquence} de la valeur $x_i$ est $f_i = \dfrac{n_i}{N}$. On a $\sum f_i = 1$ (ou $100\%$).
\end{itemize}
\end{definition}

\begin{definition}[Effectifs cumulés]
Pour une série ordonnée ($x_1 < x_2 < \dots < x_p$) :
\begin{itemize}
    \item L'\cle{effectif cumulé croissant (ECC)} de $x_i$ est $N_i = \sum_{k=1}^{i} n_k$. C'est le nombre d'individus tels que $X \le x_i$.
    \item L'\cle{effectif cumulé décroissant (ECD)} de $x_i$ est $N'_i = \sum_{k=i}^{p} n_k$. C'est le nombre d'individus tels que $X \ge x_i$.
\end{itemize}
\end{definition}

\begin{methode}[Construction d'un tableau d'effectifs cumulés]
\begin{enumerate}
    \item Trier les valeurs du caractère par ordre croissant.
    \item Reporter les effectifs $n_i$.
    \item Calculer les fréquences $f_i = n_i / N$ (souvent en $\%$).
    \item Calculer les ECC : $N_1 = n_1$, puis $N_i = N_{i-1} + n_i$. Vérifier : $N_p = N$.
    \item (Optionnel) Calculer les ECD : $N'_p = n_p$, puis $N'_i = N'_{i+1} + n_i$. Vérifier : $N'_1 = N$.
\end{enumerate}
\end{methode}

\begin{exemplebox}[Exemple : Consommation d'eau (CAMWATER)]
Une enquête auprès de $20$ foyers à Yaoundé donne la consommation mensuelle (en $\text{m}^3$) :
\[ 12,\; 15,\; 14,\; 12,\; 18,\; 15,\; 12,\; 16,\; 14,\; 15,\; 17,\; 15,\; 13,\; 14,\; 15,\; 16,\; 14,\; 12,\; 15,\; 14 \]
\textbf{Tableau :}
\begin{center}
\begin{tabularx}{\linewidth}{|l|*{7}{>{\centering\arraybackslash}X|}}\hline
\rowcolor{popblueL} \textbf{Valeur $x_i$} & 12 & 13 & 14 & 15 & 16 & 17 & 18 \\ \hline
\textbf{Effectif $n_i$} & 4 & 1 & 5 & 6 & 2 & 1 & 1 \\ \hline
\textbf{Fréquence $f_i$} & 0,20 & 0,05 & 0,25 & 0,30 & 0,10 & 0,05 & 0,05 \\ \hline
\textbf{ECC $N_i$} & 4 & 5 & 10 & 16 & 18 & 19 & 20 \\ \hline
\end{tabularx}
\end{center}
$N = 20$. La valeur la plus fréquente (mode) est $15~\text{m}^3$.
\end{exemplebox}

\courssub{Groupement en classes (caractère continu)}

Pour un caractère continu, on regroupe les valeurs en \cle{classes} d'amplitude $a$ (souvent constante).
\begin{itemize}
    \item Une classe s'écrit $[a ; b[$ (borne inférieure incluse, supérieure exclue), sauf la dernière $[a ; b]$.
    \item Le \cle{centre de la classe} (ou milieu) est $x_i = \dfrac{a+b}{2}$. C'est la valeur qui représentera la classe pour les calculs (moyenne, etc.).
    \item L'\cle{effectif de la classe} est le nombre de valeurs dans l'intervalle.
\end{itemize}

\begin{attention}[Règle de groupement]
Les classes doivent être \textbf{exhaustives} (toutes les valeurs y sont) et \textbf{mutuellement exclusives} (pas de chevauchement). On utilise la notation $[a ; b[$ pour éviter l'ambiguïté aux frontières.
\end{attention}

\coursec{Représentations graphiques}

Le choix du diagramme dépend de la nature du caractère.

\courssub{Caractère qualitatif ou quantitatif discret : Diagramme à bâtons et Diagramme à bandes}

\begin{definition}[Diagramme à bâtons (ou en barres verticales)]
Chaque valeur $x_i$ (ou modalité) est représentée par un rectangle (bâton) de largeur constante, dont la \textbf{hauteur} est proportionnelle à l'effectif $n_i$ (ou à la fréquence $f_i$). Les bâtons sont \textbf{séparés} (espace entre eux) pour bien montrer le caractère discret.
\end{definition}

\begin{definition}[Diagramme à bandes (ou barres horizontales)]
Même principe, mais les rectangles sont horizontaux. Pratique quand les modalités ont de longs libellés (ex: \og Agriculture \fg, \og Industrie \fg, \og Services \fg).
\end{definition}

\courssub{Diagramme circulaire (Camembert)}

\begin{definition}[Diagramme circulaire]
Un cercle (le \og camembert \fg) est divisé en secteurs angulaires. L'angle au centre du secteur correspondant à $x_i$ est :
\[ \alpha_i = 360^\circ \times f_i = 360^\circ \times \frac{n_i}{N} \]
L'aire du secteur est proportionnelle à l'effectif.
\end{definition}

\begin{methode}[Construction d'un diagramme circulaire]
\begin{enumerate}
    \item Calculer les fréquences $f_i$ (ou pourcentages).
    \item Calculer les angles $\alpha_i = 360^\circ \times f_i$. Vérifier $\sum \alpha_i = 360^\circ$.
    \item Tracer le cercle, reporter les angles au rapporteur, colorier/étiqueter.
\end{enumerate}
\end{methode}

\begin{exemplebox}[Parts de marché Mobile Money au Cameroun (Données fictives)]
\begin{center}
\begin{tabular}{|c|c|c|c|}\hline
\rowcolor{popblueL} Opérateur & Effectif $n_i$ & Fréquence $f_i$ & Angle $\alpha_i$ \\ \hline
MTN MoMo & 1200 & 0,40 & $144^\circ$ \\
Orange Money & 1000 & 0,333... & $120^\circ$ \\
Express Union & 500 & 0,166... & $60^\circ$ \\
Autres & 300 & 0,10 & $36^\circ$ \\ \hline
Total & 3000 & 1 & $360^\circ$ \\ \hline
\end{tabular}
\end{center}
\end{exemplebox}

\begin{popfigure}
\begin{center}
\begin{tikzpicture}
\def\R{2.5}
\fill[popblue]   (0,0) -- (0:\R)   arc (0:144:\R)   -- cycle;
\fill[poporange] (0,0) -- (144:\R) arc (144:264:\R) -- cycle;
\fill[popgreen]  (0,0) -- (264:\R) arc (264:324:\R) -- cycle;
\fill[poppurple] (0,0) -- (324:\R) arc (324:360:\R) -- cycle;
\draw[popink, line width=1.2pt] (0,0) circle (\R);
\foreach \a in {0,144,264,324} \draw[popink, line width=1.2pt] (0,0) -- (\a:\R);
\node[white, font=\small\bfseries] at (72:1.5)  {40\,\%};
\node[white, font=\small\bfseries] at (204:1.5) {33,3\,\%};
\node[white, font=\small\bfseries] at (294:1.6) {16,7\,\%};
\node[white, font=\footnotesize\bfseries] at (342:1.8) {10\,\%};
\begin{scope}[xshift=3.6cm, yshift=1.2cm, font=\small]
\foreach \c/\t [count=\k] in {popblue/MTN MoMo, poporange/Orange Money, popgreen/Express Union, poppurple/Autres}
  {\fill[\c] (0,-0.6*\k) rectangle ++(0.35,0.35); \node[anchor=west] at (0.45,-0.6*\k+0.175) {\t};}
\end{scope}
\end{tikzpicture}
\end{center}
\legende{Diagramme circulaire des parts de marché Mobile Money (angles : $144^\circ$, $120^\circ$, $60^\circ$, $36^\circ$).}
\end{popfigure}

\courssub{Caractère quantitatif continu : Histogramme}

\begin{definition}[Histogramme]
Pour une série groupée en classes $[b_i ; b_{i+1}[$ d'amplitude $a_i = b_{i+1} - b_i$ :
\begin{itemize}
    \item L'abscisse porte les bornes des classes.
    \item On trace des rectangles \textbf{joints} (pas d'espace) sur chaque classe.
    \item L'\textbf{aire} du rectangle est proportionnelle à l'effectif de la classe.
    \item La \textbf{hauteur} (densité) est $h_i = \dfrac{n_i}{a_i}$ (effectif divisé par amplitude).
\end{itemize}
Si toutes les classes ont la \textbf{même amplitude}, la hauteur est simplement proportionnelle à l'effectif (on peut prendre $h_i = n_i$ ou $h_i = f_i$).
\end{definition}

\begin{attention}[Erreur classique au Bac]
Ne tracez \textbf{jamais} un histogramme avec des hauteurs égales aux effectifs si les amplitudes des classes sont différentes ! L'aire ne serait plus proportionnelle à l'effectif. Calculez toujours la densité $h_i = n_i / a_i$.
\end{attention}

\begin{exemplebox}[Histogramme : Factures ENEO (en FCFA) - Classes d'amplitudes inégales]
\begin{center}
\begin{tabular}{|c|c|c|c|}\hline
\rowcolor{popblueL} Classe (FCFA) & Amplitude $a_i$ & Effectif $n_i$ & Hauteur $h_i = n_i/a_i$ \\ \hline
$[0 ; 5000[$ & 5000 & 40 & $40/5000 = 0,008$ \\
$[5000 ; 10000[$ & 5000 & 60 & $60/5000 = 0,012$ \\
$[10000 ; 20000[$ & 10000 & 50 & $50/10000 = 0,005$ \\
$[20000 ; 50000[$ & 30000 & 30 & $30/30000 = 0,001$ \\ \hline
\end{tabular}
\end{center}
\end{exemplebox}

\begin{popfigure}
\begin{center}
\begin{tikzpicture}
\begin{axis}[
    width=\linewidth, height=7cm,
    xlabel={Facture (FCFA)}, ylabel={Densité $h_i$},
    xmin=0, xmax=52000, ymin=0, ymax=0.014,
    xtick={0,5000,10000,20000,50000},
    xticklabels={0, 5\,000, 10\,000, 20\,000, 50\,000},
    ytick={0.001,0.005,0.008,0.012},
    yticklabels={0{,}001, 0{,}005, 0{,}008, 0{,}012},
    scaled x ticks=false, scaled y ticks=false,
    axis lines*=left,
]
\addplot[fill=popblue, draw=popdark] coordinates {(0,0) (5000,0) (5000,0.008) (0,0.008)} -- cycle;
\addplot[fill=popblue!70, draw=popdark] coordinates {(5000,0) (10000,0) (10000,0.012) (5000,0.012)} -- cycle;
\addplot[fill=popblue!50, draw=popdark] coordinates {(10000,0) (20000,0) (20000,0.005) (10000,0.005)} -- cycle;
\addplot[fill=popblue!30, draw=popdark] coordinates {(20000,0) (50000,0) (50000,0.001) (20000,0.001)} -- cycle;
\end{axis}
\end{tikzpicture}
\end{center}
\legende{Histogramme des factures ENEO : les aires des rectangles sont proportionnelles aux effectifs.}
\end{popfigure}

\courssub{Polygone des fréquences cumulées croissantes}

\begin{definition}[Polygone des fréquences cumulées croissantes]
C'est la courbe reliant les points $M_i(x_i ; F_i)$ où $x_i$ est la \textbf{borne supérieure} de la classe (ou la valeur pour le discret) et $F_i = \dfrac{N_i}{N}$ est la fréquence cumulée croissante (en $\%$ ou en décimal). On ajoute souvent le point origine $(x_0 ; 0)$ où $x_0$ est la borne inférieure de la première classe.
\end{definition}

\begin{propriete}[Lecture graphique de la médiane et des quartiles]
Sur le polygone des fréquences cumulées :
\begin{itemize}
    \item La \cle{médiane} $Me$ correspond à l'abscisse du point d'ordonnée $0,5$ (ou $50\%$).
    \item Le \cle{premier quartile} $Q_1$ correspond à l'abscisse d'ordonnée $0,25$ ($25\%$).
    \item Le \cle{troisième quartile} $Q_3$ correspond à l'abscisse d'ordonnée $0,75$ ($75\%$).
\end{itemize}
C'est la méthode graphique officielle (pas d'interpolation linéaire calculée).
\end{propriete}

\begin{popfigure}
\begin{center}
\begin{tikzpicture}
\begin{axis}[
    width=\linewidth, height=8cm,
    xlabel={Consommation (m$^3$)}, ylabel={Fréquence cumulée $F$},
    xmin=10, xmax=20, ymin=0, ymax=1.05,
    axis lines=middle, grid=major,
    xtick={10,11,12,13,14,15,16,17,18,19,20},
    ytick={0,0.25,0.5,0.75,1},
    yticklabels={0\%, 25\%, 50\%, 75\%, 100\%},
]
% Example data points for cumulative frequency polygon
\addplot[popcurve, thick, mark=*, mark size=2, color=popblue] coordinates {
    (11.5, 0) (12.5, 0.2) (13.5, 0.25) (14.5, 0.5) (15.5, 0.8) (16.5, 0.9) (17.5, 0.95) (18.5, 1)
};
% Median lines
\draw[dashed, poporange, thick] (11.5, 0.5) -- (14.5, 0.5) -- (14.5, 0) node[below] {$Me \approx 14,5$};
\draw[dashed, popgreen, thick] (11.5, 0.25) -- (13.5, 0.25) -- (13.5, 0) node[below] {$Q_1 \approx 13,5$};
\draw[dashed, popgreen, thick] (11.5, 0.75) -- (15.5, 0.75) -- (15.5, 0) node[below] {$Q_3 \approx 15,5$};
\end{axis}
\end{tikzpicture}
\end{center}
\legende{Polygone des fréquences cumulées : lecture graphique de $Q_1$, $Me$, $Q_3$.}
\end{popfigure}

\coursec{Paramètres de position : Mode, Moyenne, Médiane}

Ces indicateurs résument la \og position \fg centrale de la série.

\courssub{Le Mode}

\begin{definition}[Mode]
Le \cle{mode} (noté $Mod$) est la valeur (ou la classe) qui possède l'\textbf{effectif maximal}.
\begin{itemize}
    \item Série discrète : valeur $x_i$ telle que $n_i = \max(n_k)$.
    \item Série groupée : \cle{classe modale} = classe d'effectif maximal.
\end{itemize}
Une série peut être \textbf{unimodale} (un mode), \textbf{bimodale} (deux modes), ou \textbf{amodale} (effectifs tous égaux).
\end{definition}

\courssub{La Moyenne}

\begin{definition}[Moyenne arithmétique]
La \cle{moyenne} (notée $\bar{x}$) est la somme des valeurs pondérées par leurs effectifs, divisée par l'effectif total.
\[ \bar{x} = \frac{1}{N} \sum_{i=1}^{p} n_i x_i = \sum_{i=1}^{p} f_i x_i \]
Pour une série groupée en classes, on utilise le \textbf{centre de la classe} $x_i$ comme valeur représentative.
\end{definition}

\begin{propriete}[Propriétés de la moyenne]
\begin{enumerate}
    \item $\sum n_i (x_i - \bar{x}) = 0$ (somme des écarts à la moyenne nulle).
    \item Si on ajoute une constante $c$ à toutes les valeurs, la moyenne augmente de $c$.
    \item Si on multiplie toutes les valeurs par un coefficient $k$, la moyenne est multipliée par $k$.
\end{enumerate}
\end{propriete}

\courssub{La Médiane}

\begin{definition}[Médiane]
La \cle{médiane} (notée $Me$) est la valeur qui partage la population en deux sous-populations de même effectif (ou aussi proches que possible) : $50\%$ des valeurs sont $\le Me$ et $50\%$ sont $\ge Me$.
\end{definition}

\begin{methode}[Calcul de la médiane (Sans interpolation linéaire --- Hors programme)]
\begin{enumerate}
    \item Calculer $N/2$.
    \item \textbf{Série discrète (valeurs brutes) :}
    \begin{itemize}
        \item Si $N$ est \textbf{impair} : $Me$ = valeur de rang $\frac{N+1}{2}$ dans la série ordonnée.
        \item Si $N$ est \textbf{pair} : $Me$ = demi-somme des valeurs de rangs $\frac{N}{2}$ et $\frac{N}{2}+1$.
    \end{itemize}
    \item \textbf{Série groupée en classes :} Repérer la \textbf{première classe} dont l'ECC est \textbf{supérieur ou égal} à $N/2$. C'est la \cle{classe médiane}. La médiane est \textbf{n'importe quelle valeur} dans cette classe (souvent on prend le centre de la classe, ou on dit \og la médiane appartient à la classe $[a;b[$ \fg). \textbf{Pas de formule d'interpolation.}
    \item \textbf{Lecture graphique :} Sur le polygone des fréquences cumulées, abscisse du point d'ordonnée $0,5$.
\end{enumerate}
\end{methode}

\begin{attention}[Rédaction attendue au Bac]
\begin{itemize}
    \item Toujours écrire : \og $N = \dots$, donc $N/2 = \dots$ \fg.
    \item Pour $N$ pair : \og La médiane est la demi-somme des valeurs de rangs $N/2$ et $N/2+1$, soit $Me = \frac{x_{N/2} + x_{N/2+1}}{2} = \dots$ \fg.
    \item Pour classe médiane : \og $N/2 = \dots$ ; la première classe dont l'ECC $\ge N/2$ est $[a;b[$. La médiane appartient à cette classe. \fg (On peut donner le centre).
\end{itemize}
\end{attention}

\begin{exoresolu}[Médiane : Série discrète et groupée]
\textbf{1. Série discrète :} Effectifs : $n_1=3 (x=10), n_2=5 (x=12), n_3=2 (x=14)$. $N=10$ (pair). $N/2=5$.
ECC : $3, 8, 10$. Rangs 5 et 6 sont dans la valeur $12$ (ECC=8). $Me = \frac{12+12}{2} = 12$.

\textbf{2. Série groupée :} Classes $[0;10[$ ($n=5$), $[10;20[$ ($n=12$), $[20;30[$ ($n=8$). $N=25$ (impair). $N/2=12,5$.
ECC : $5, 17, 25$. Première classe avec ECC $\ge 12,5$ est $[10;20[$. Classe médiane = $[10;20[$. On peut dire $Me \approx 15$ (centre).
\tcblower
\textbf{Solution détaillée :}
1. $N=10$, pair. Rangs 5 et 6. Valeurs ordonnées : 10, 10, 10, 12, \textbf{12, 12}, 12, 12, 14, 14. $Me = (12+12)/2 = 12$.
2. $N=25$, $N/2=12,5$. ECC : 5 (classe 1), 17 (classe 2). $17 \ge 12,5$. Classe médiane = $[10;20[$. Centre = 15.
\end{exoresolu}

\coursec{Quartiles et Déciles}

Ils généralisent la médiane pour partager la population en 4 ou 10 parts égales.

\begin{definition}[Quartiles]
Les \cle{quartiles} divisent la série ordonnée en 4 parties d'effectifs égaux (ou proches) :
\begin{itemize}
    \item $Q_1$ (premier quartile) : $25\%$ des valeurs $\le Q_1$. ECC $\ge N/4$.
    \item $Q_2$ = Médiane ($Me$) : $50\%$. ECC $\ge N/2$.
    \item $Q_3$ (troisième quartile) : $75\%$ des valeurs $\le Q_3$. ECC $\ge 3N/4$.
\end{itemize}
\end{definition}

\begin{definition}[Déciles]
Les \cle{déciles} $D_1, D_2, \dots, D_9$ divisent la série en 10 parties égales :
$D_k$ est la première valeur dont l'ECC $\ge \dfrac{kN}{10}$.
$D_5 = Me$.
\end{definition}

\begin{methode}[Détermination des quartiles et déciles (Règle officielle : sans interpolation)]
\begin{enumerate}
    \item Calculer les seuils : $N/4$, $N/2$, $3N/4$ (quartiles) ou $kN/10$ (déciles).
    \item Sur le tableau d'ECC (série discrète) ou le polygone (série groupée), repérer la \textbf{première valeur (ou classe)} dont l'ECC \textbf{atteint ou dépasse} le seuil.
    \item \textbf{Discrète :} La valeur trouvée est le quartile/décile.
    \item \textbf{Groupée :} La classe trouvée est la classe du quartile/décile. On donne souvent le centre de la classe.
\end{enumerate}
\end{methode}

\begin{exemplebox}[Quartiles sur l'exemple Consommation d'eau ($N=20$)]
$N/4 = 5$ ; $N/2 = 10$ ; $3N/4 = 15$.
ECC : $x=12 \to 4$ ; $x=13 \to 5$ ; $x=14 \to 10$ ; $x=15 \to 16$ ; $x=16 \to 18$...
\begin{itemize}
    \item $Q_1$ : Seuil 5. Première valeur ECC $\ge 5$ est $13$ (ECC=5). $\Rightarrow Q_1 = 13$.
    \item $Me=Q_2$ : Seuil 10. Première valeur ECC $\ge 10$ est $14$ (ECC=10). $\Rightarrow Me = 14$. (Note : ici $N$ pair, médiane = demi-somme rangs 10 et 11. Rang 10 = 14, Rang 11 = 15 $\Rightarrow Me=14,5$. \textbf{Attention :} La règle ECC $\ge N/2$ donne la classe/valeur qui \og contient \fg la médiane. Pour le discret pair, il faut appliquer la règle des rangs. La règle ECC est parfaite pour le groupé et le graphique).
    \item $Q_3$ : Seuil 15. Première valeur ECC $\ge 15$ est $15$ (ECC=16). $\Rightarrow Q_3 = 15$.
\end{itemize}
\end{exemplebox}

\begin{attention}[Distinction Discret Pair / Groupé]
Pour une série \textbf{discrète à effectif total pair}, la médiane (et $Q_1, Q_3$ si les seuils tombent sur des entiers) se calcule par la \textbf{demi-somme des deux valeurs centrales} (rangs $N/2$ et $N/2+1$). La règle \og première valeur ECC $\ge$ seuil \fg donne la valeur qui \textbf{commence} la seconde moitié. Au Bac, pour le discret pair, \textbf{appliquez la règle des rangs (demi-somme)}. Pour le groupé et la lecture graphique, utilisez la règle ECC / graphique.
\end{attention}

\coursec{Séries statistiques à deux caractères : Nuage de points et Point moyen}

On observe deux caractères quantitatifs $X$ et $Y$ sur la même population de taille $n$. Données : $(x_1, y_1), (x_2, y_2), \dots, (x_n, y_n)$.

\courssub{Nuage de points}

\begin{definition}[Nuage de points]
C'est la représentation graphique dans un repère orthogonal des points $M_i(x_i ; y_i)$. L'axe des abscisses porte $X$, l'axe des ordonnées porte $Y$.
\end{definition}

\begin{savaistu}[Contexte camerounais]
Étudier la relation entre le \textbf{prix du kg de plantain} (abscisse $X$) et la \textbf{quantité vendue par jour} au marché Mfoundi (ordonnée $Y$) sur 12 mois. Ou entre l'\textbf{année} (rang $x$) et le \textbf{nombre d'abonnés Internet} (Cameroon Telecom / MTN / Orange) en millions ($y$).
\end{savaistu}

\courssub{Point moyen (Centre de gravité)}

\begin{definition}[Point moyen]
Le \cle{point moyen} (ou centre de gravité) du nuage est le point $G(\bar{x} ; \bar{y})$ où :
\[ \bar{x} = \frac{1}{n}\sum_{i=1}^{n} x_i \qquad \bar{y} = \frac{1}{n}\sum_{i=1}^{n} y_i \]
\end{definition}

\begin{propriete}[Propriété fondamentale]
Le point moyen $G$ est le centre de gravité du nuage de points (masses égales aux sommets). \textbf{Toute droite d'ajustement affine pertinente doit passer par $G$.}
\end{propriete}

\begin{exemplebox}[Nuage de points : Année / Abonnés Internet (Millions) - Données fictives]
\begin{center}
\begin{tabular}{|c|cccccc|}\hline
\rowcolor{popblueL} Année & 2018 & 2019 & 2020 & 2021 & 2022 & 2023 \\ \hline
Rang $x$ & 1 & 2 & 3 & 4 & 5 & 6 \\
Abonnés $y$ & 4,2 & 5,1 & 6,0 & 7,2 & 8,5 & 9,8 \\ \hline
\end{tabular}
\end{center}
$n=6$. $\bar{x} = \frac{1+2+3+4+5+6}{6} = 3,5$.
$\bar{y} = \frac{4,2+5,1+6,0+7,2+8,5+9,8}{6} = \frac{40,8}{6} = 6,8$.
$G(3,5 ; 6,8)$.
\end{exemplebox}

\begin{popfigure}
\begin{center}
\begin{tikzpicture}
\begin{axis}[
    width=\linewidth, height=8cm,
    xlabel={Rang de l'année $x$}, ylabel={Abonnés (Millions) $y$},
    xmin=0, xmax=7, ymin=3, ymax=11,
    axis lines=middle, grid=major,
    xtick={1,2,3,4,5,6}, ytick={4,5,6,7,8,9,10},
]
% Points
\addplot[only marks, mark=*, mark size=3, color=popblue] coordinates {
    (1,4.2) (2,5.1) (3,6.0) (4,7.2) (5,8.5) (6,9.8)
};
% Point moyen G
\addplot[only marks, mark=*, mark size=4, color=poporange] coordinates {(3.5,6.8)};
\node[poporange, above right] at (axis cs:3.5,6.8) {$G(3,5 ; 6,8)$};
\end{axis}
\end{tikzpicture}
\end{center}
\legende{Nuage de points et point moyen $G$. On devine une tendance linéaire croissante.}
\end{popfigure}

\coursec{Ajustement affine par la méthode de Mayer et prévisions}

Quand le nuage de points suggère une relation linéaire ($Y \approx aX + b$), on cherche la \cle{droite d'ajustement} (ou droite de régression) qui \og colle \fg le mieux au nuage. La méthode de Mayer est la méthode officielle au programme (pas de moindres carrés).

\courssub{Principe de la méthode de Mayer}

\begin{methode}[Méthode de Mayer : Étapes officielles]
\begin{enumerate}
    \item \textbf{Trier} les points par abscisses $x_i$ croissantes.
    \item \textbf{Partager} le nuage en \textbf{deux sous-nuages} d'effectifs \textbf{égaux} (si $n$ pair) ou \textbf{proches} (si $n$ impair, on écarte le point central ou on fait des groupes de tailles $(n-1)/2$ et $(n+1)/2$). La séparation se fait \textbf{verticalement} (selon $x$).
    \item Calculer le \textbf{point moyen} $G_1(\bar{x}_1 ; \bar{y}_1)$ du premier sous-nuage (les $x$ les plus petits).
    \item Calculer le \textbf{point moyen} $G_2(\bar{x}_2 ; \bar{y}_2)$ du second sous-nuage (les $x$ les plus grands).
    \item La \cle{droite d'ajustement de Mayer} est la droite $(G_1G_2)$.
    \item Son équation réduite est $y = ax + b$ avec :
    \[ a = \frac{\bar{y}_2 - \bar{y}_1}{\bar{x}_2 - \bar{x}_1} \qquad b = \bar{y}_1 - a\bar{x}_1 \quad (\text{ou } b = \bar{y} - a\bar{x} \text{ car } G \in (G_1G_2)) \]
\end{enumerate}
\end{methode}

\begin{propriete}[Droite de Mayer et Point moyen global]
La droite de Mayer $(G_1G_2)$ passe \textbf{toujours} par le point moyen global $G(\bar{x} ; \bar{y})$ de l'ensemble du nuage.
\textbf{Preuve :} $G$ est le barycentre de $G_1$ (poids $n_1$) et $G_2$ (poids $n_2$). Comme $n_1 = n_2$ (ou presque), $G$ est le milieu de $[G_1G_2]$ (ou très proche). Donc $G \in (G_1G_2)$.
\end{propriete}

\begin{attention}[Pièges de la méthode de Mayer]
\begin{itemize}
    \item \textbf{Triage obligatoire :} Les points doivent être classés par $x$ croissant \textbf{avant} de partager en deux groupes. Ne pas séparer au hasard !
    \item \textbf{Effectifs égaux :} Les deux groupes doivent avoir le même nombre de points (ou différer de 1 au maximum). Si $n$ est impair, le point du milieu (médiane des $x$) est généralement \textbf{exclu} des deux groupes pour calculer $G_1$ et $G_2$.
    \item \textbf{Coefficient directeur :} $\bar{x}_2 - \bar{x}_1 > 0$ car le groupe 2 a des $x$ plus grands. Le signe de $a$ donne la tendance (croissante si $a>0$, décroissante si $a<0$).
\end{itemize}
\end{attention}

\courssub{Application numérique complète (Méthode de Mayer)}

Reprenons l'exemple des abonnés Internet ($n=6$, pair).

\begin{exoresolu}[Droite de Mayer et Prévision 2025]
\textbf{Données triées par $x$ (déjà fait) :}
Groupe 1 (premiers 3 points) : $(1;4,2), (2;5,1), (3;6,0)$
Groupe 2 (derniers 3 points) : $(4;7,2), (5;8,5), (6;9,8)$

\textbf{1. Points moyens partiels :}
$G_1$ : $\bar{x}_1 = \frac{1+2+3}{3} = 2$ ; $\bar{y}_1 = \frac{4,2+5,1+6,0}{3} = \frac{15,3}{3} = 5,1$. $\Rightarrow G_1(2 ; 5,1)$.
$G_2$ : $\bar{x}_2 = \frac{4+5+6}{3} = 5$ ; $\bar{y}_2 = \frac{7,2+8,5+9,8}{3} = \frac{25,5}{3} = 8,5$. $\Rightarrow G_2(5 ; 8,5)$.

\textbf{2. Coefficient directeur $a$ :}
$a = \frac{8,5 - 5,1}{5 - 2} = \frac{3,4}{3} \approx 1,1333$ (Millions par an).

\textbf{3. Ordonnée à l'origine $b$ (via $G_1$) :}
$b = \bar{y}_1 - a\bar{x}_1 = 5,1 - 1,1333 \times 2 = 5,1 - 2,2666 = 2,8334$.
\textbf{Vérification via $G(3,5 ; 6,8)$ :} $b = 6,8 - 1,1333 \times 3,5 = 6,8 - 3,9666 = 2,8334$. \textbf{OK.}

\textbf{4. Équation de la droite de Mayer :}
\[ \boxed{y = 1,1333\,x + 2,8334} \]
(On arrondit selon consignes, ex: $y = 1,13x + 2,83$).

\textbf{5. Prévision pour 2025 (Rang $x=8$) :}
$y = 1,1333 \times 8 + 2,8334 = 9,0664 + 2,8334 = 11,8998 \approx 11,9$ millions d'abonnés.
\tcblower
\textbf{Solution détaillée :}
Voir calculs ci-dessus. La prévision s'obtient en remplaçant $x$ par le rang de l'année visée (8 pour 2025 car 2018 $\to$ 1). Attention : si $x$ est l'année civile (2018, 2019...), on utilise directement l'année. Ici $x$ est le rang.
\end{exoresolu}

\begin{popfigure}
\begin{center}
\begin{tikzpicture}
\begin{axis}[
    width=\linewidth, height=8cm,
    xlabel={Rang $x$}, ylabel={Abonnés $y$ (Millions)},
    xmin=0, xmax=9, ymin=3, ymax=13,
    axis lines=middle, grid=major,
    xtick={1,2,3,4,5,6,7,8}, ytick={4,5,6,7,8,9,10,11,12},
]
% Points
\addplot[only marks, mark=*, mark size=3, color=popblue] coordinates {
    (1,4.2) (2,5.1) (3,6.0) (4,7.2) (5,8.5) (6,9.8)
};
% Point moyen G
\addplot[only marks, mark=*, mark size=4, color=poporange] coordinates {(3.5,6.8)};
\node[poporange, above right] at (axis cs:3.5,6.8) {$G$};
% Mayer Line
\addplot[domain=0:8.5, samples=2, thick, poporange, dashed] {1.1333*x + 2.8334};
\node[poporange, right] at (axis cs:8, 11.9) {$y=1,13x+2,83$};
% Prevision 2025 (x=8)
\addplot[only marks, mark=square*, mark size=3, color=popgreen] coordinates {(8,11.9)};
\node[popgreen, above left] at (axis cs:8,11.9) {Prévision 2025};
% G1 G2
\addplot[only marks, mark=triangle*, mark size=3, color=poppurple] coordinates {(2,5.1) (5,8.5)};
\node[poppurple, below left] at (axis cs:2,5.1) {$G_1$};
\node[poppurple, above right] at (axis cs:5,8.5) {$G_2$};
\end{axis}
\end{tikzpicture}
\end{center}
\legende{Nuage, points moyens $G_1, G_2, G$, droite de Mayer et prévision pour $x=8$ (2025).}
\end{popfigure}

\courssub{Interprétation et limites}

\begin{aretenir}[Interprétation des coefficients]
\begin{itemize}
    \item $a$ (pente) : Variation moyenne de $Y$ quand $X$ augmente de 1 unité. Ici $a \approx 1,13$ : le nombre d'abonnés augmente en moyenne de 1,13 million par an.
    \item $b$ (ordonnée à l'origine) : Valeur théorique de $Y$ pour $X=0$. Ici $b \approx 2,83$ : extrapolation pour l'année "rang 0" (2017). Sens physique parfois absent.
    \item \textbf{Prévision :} Valable seulement \textbf{à l'intérieur} de l'intervalle des $x$ observés (interpolation) ou \textbf{proche} (extrapolation courte). Méfiance pour $x$ très lointain.
\end{itemize}
\end{aretenir}

\coursec{Synthèse et Entraînement Bac}

\begin{aretenir}[Synthèse des formules et méthodes clés (À connaître par cœur)]
\begin{center}
\begin{tabularx}{\linewidth}{|l|X|}\hline
\rowcolor{popblueL} \textbf{Notion} & \textbf{Règle / Formule} \\ \hline
\textbf{Fréquence} & $f_i = \dfrac{n_i}{N}$ \quad ($\sum f_i = 1$) \\ \hline
\textbf{Moyenne} & $\bar{x} = \dfrac{1}{N}\sum n_i x_i$ (Groupé : $x_i$ = centre de classe) \\ \hline
\textbf{Mode} & Valeur/Classe d'effectif maximal \\ \hline
\textbf{Médiane (Discrète)} & $N$ impair : valeur rang $\frac{N+1}{2}$ \quad $N$ pair : $\frac{1}{2}(x_{N/2} + x_{N/2+1})$ \\ \hline
\textbf{Médiane (Groupée)} & Classe où ECC $\ge N/2$ (Pas d'interpolation) \\ \hline
\textbf{Quartiles $Q_1, Q_3$ / Déciles $D_k$} & Première valeur/classe avec ECC $\ge N/4,\; 3N/4,\; kN/10$ \\ \hline
\textbf{Diagramme circulaire} & Angle $\alpha_i = 360^\circ \times f_i$ \\ \hline
\textbf{Histogramme} & Aire $\propto$ Effectif \quad Hauteur $h_i = n_i / a_i$ (densité) \\ \hline
\textbf{Polygone Freq. Cum.} & Points $(x_i^{\text{sup}} ; F_i)$ \quad Lecture $Q_1, Me, Q_3$ aux ordonnées $0,25 ; 0,5 ; 0,75$ \\ \hline
\textbf{Point moyen $G$} & $G\left(\frac{\sum x_i}{n} ; \frac{\sum y_i}{n}\right)$ \\ \hline
\textbf{Méthode de Mayer} & 1. Trier par $x$. 2. Deux groupes $n/2$. 3. $G_1, G_2$. 4. Droite $(G_1G_2)$ : $a=\frac{\bar{y}_2-\bar{y}_1}{\bar{x}_2-\bar{x}_1},\; b=\bar{y}-a\bar{x}$ \\ \hline
\textbf{Prévision} & Remplacer $x$ par la valeur future dans $y=ax+b$ (Attention : $x$ = rang ou année ?) \\ \hline
\end{tabularx}
\end{center}
\end{aretenir}

\begin{aretenir}[Checklist finale avant le Baccalauréat]
\begin{itemize}
\item[$\square$] Je distingue caractère qualitatif / quantitatif discret / quantitatif continu.
\item[$\square$] Je construis un tableau complet : valeurs, effectifs, fréquences (\%), ECC, ECD.
\item[$\square$] Je choisis et construis le bon diagramme : bâtons (discret), bandes (qualitatif long), circulaire (parts), histogramme (continu, \textbf{aire = effectif}), polygone cumulé (lecture $Q_1, Me, Q_3$).
\item[$\square$] Je calcule le mode (valeur ou classe modale).
\item[$\square$] Je calcule la moyenne (série brute ou groupée avec centres de classes).
\item[$\square$] Je détermine la médiane : cas $N$ pair (demi-somme) / $N$ impair (rang central) / groupé (classe ECC $\ge N/2$).
\item[$\square$] Je détermine $Q_1, Q_3, D_k$ par la règle ECC $\ge$ seuil (sans interpolation).
\item[$\square$] Je place un nuage de points, je calcule et place $G(\bar{x};\bar{y})$.
\item[$\square$] J'applique la méthode de Mayer : tri par $x$, partage en 2 groupes égaux, calcul $G_1, G_2$, équation $(G_1G_2)$, vérification $G \in (G_1G_2)$.
\item[$\square$] Je fais une prévision en remplaçant $x$ dans l'équation (vigilance sur le codage de $x$ : rang vs année).
\item[$\square$] Je conclus par une phrase : \og La médiane est 14,5 m$^3$, cela signifie que 50\% des foyers consomment moins de 14,5 m$^3$. \fg
\end{itemize}
\end{aretenir}

\coursec{Exercices d'entraînement (Style Bac)}

\begin{exercice}[Exercice 1 : Série discrète - Prix du carburant]
On relève le prix du litre de \og Super \fg (en FCFA) dans 15 stations-services de Douala :
\[ 730,\; 735,\; 730,\; 740,\; 735,\; 730,\; 745,\; 735,\; 730,\; 740,\; 735,\; 730,\; 740,\; 735,\; 750 \]
\begin{enumerate}
    \item Construire le tableau d'effectifs, fréquences (\%), ECC.
    \item Déterminer le mode, la moyenne, la médiane, $Q_1$ et $Q_3$.
    \item Représenter la série par un diagramme à bâtons et un diagramme circulaire.
\end{enumerate}
\end{exercice}

\begin{corrige}[Exercice 1]
\begin{enumerate}
    \item \textbf{Tableau :}
    \begin{center}
    \begin{tabular}{|c|ccccc|}\hline
    \rowcolor{popblueL} Prix $x_i$ & 730 & 735 & 740 & 745 & 750 \\ \hline
    Effectif $n_i$ & 5 & 5 & 3 & 1 & 1 \\ \hline
    Fréquence $f_i$ & 33,3\% & 33,3\% & 20\% & 6,7\% & 6,7\% \\ \hline
    ECC $N_i$ & 5 & 10 & 13 & 14 & 15 \\ \hline
    \end{tabular}
    \end{center}
    \item \textbf{Mode :} Bimodale : 730 et 735 (effectif max = 5).
    \textbf{Moyenne :} $\bar{x} = \frac{1}{15}(5\times730 + 5\times735 + 3\times740 + 745 + 750) = \frac{11020}{15} \approx 734,67$ FCFA.
    \textbf{Médiane :} $N=15$ (impair). Rang $\frac{16}{2}=8$. Valeur rang 8 : ECC 10 correspond à 735. $Me = 735$ FCFA.
    \textbf{$Q_1$ :} Seuil $15/4 = 3,75$. Première valeur ECC $\ge 3,75$ : 730 (ECC=5). $Q_1 = 730$.
    \textbf{$Q_3$ :} Seuil $3\times15/4 = 11,25$. Première valeur ECC $\ge 11,25$ : 740 (ECC=13). $Q_3 = 740$.
    \item Diagrammes : Bâtons (hauteurs 5, 5, 3, 1, 1). Circulaire : Angles $120^\circ, 120^\circ, 72^\circ, 24^\circ, 24^\circ$.
\end{enumerate}
\end{corrige}

\begin{exercice}[Exercice 2 : Série groupée - Rendements maïs (Tonnes/Ha)]
Rendements de 40 exploitations dans l'Ouest :
\begin{center}
\begin{tabular}{|c|cccc|}\hline
\rowcolor{popblueL} Classe (T/Ha) & $[1;2[$ & $[2;3[$ & $[3;4[$ & $[4;6[$ \\ \hline
Effectif & 5 & 15 & 12 & 8 \\ \hline
\end{tabular}
\end{center}
\begin{enumerate}
    \item Compléter le tableau (centres, fréquences, ECC). Construire l'histogramme.
    \item Calculer la moyenne. Déterminer la classe médiane, $Q_1$ et $Q_3$ (par calcul et lecture graphique).
\end{enumerate}
\end{exercice}

\begin{corrige}[Exercice 2]
\begin{enumerate}
    \item \textbf{Tableau complété :}
    \begin{center}
    \begin{tabular}{|c|c|c|c|c|c|}\hline
    \rowcolor{popblueL} Classe & Centre $x_i$ & Amp. $a_i$ & $n_i$ & $f_i$ & ECC $N_i$ \\ \hline
    $[1;2[$ & 1,5 & 1 & 5 & 0,125 & 5 \\
    $[2;3[$ & 2,5 & 1 & 15 & 0,375 & 20 \\
    $[3;4[$ & 3,5 & 1 & 12 & 0,30 & 32 \\
    $[4;6[$ & 5 & 2 & 8 & 0,20 & 40 \\ \hline
    \end{tabular}
    \end{center}
    \textbf{Histogramme :} Classes 1-3 ont amplitude 1 $\to$ hauteur = effectif (5, 15, 12). Classe $[4;6[$ amplitude 2 $\to$ hauteur = $8/2 = 4$ (densité). Aire = $4 \times 2 = 8$. OK.
    \item \textbf{Moyenne :} $\bar{x} = \frac{1}{40}(5\times1,5 + 15\times2,5 + 12\times3,5 + 8\times5) = \frac{7,5 + 37,5 + 42 + 40}{40} = \frac{127}{40} = 3,175$ T/Ha.
    \textbf{Médiane :} $N/2 = 20$. ECC : 5, 20, 32, 40. Première classe ECC $\ge 20$ : $[2;3[$ (ECC=20). Classe médiane = $[2;3[$. $Me \approx 2,5$.
    \textbf{$Q_1$ :} $N/4 = 10$. Première classe ECC $\ge 10$ : $[2;3[$ (ECC=20). Classe $Q_1 = [2;3[$. $Q_1 \approx 2,5$.
    \textbf{$Q_3$ :} $3N/4 = 30$. Première classe ECC $\ge 30$ : $[3;4[$ (ECC=32). Classe $Q_3 = [3;4[$. $Q_3 \approx 3,5$.
\end{enumerate}
\end{corrige}

\begin{exercice}[Exercice 3 : Méthode de Mayer - Évolution du PIB (Billions FCFA)]
Données annuelles (Année, PIB) : $(2015, 18), (2016, 19), (2017, 20,5), (2018, 22), (2019, 23,5), (2020, 24), (2021, 26), (2022, 27,5)$.
\begin{enumerate}
    \item On code l'année par son rang $x$ (2015 $\to$ 1). Donner le nuage de points $(x_i, y_i)$.
    \item Calculer le point moyen $G$.
    \item Déterminer l'équation de la droite de Mayer.
    \item Prévoir le PIB pour 2024. Commenter la fiabilité.
\end{enumerate}
\end{exercice}

\begin{corrige}[Exercice 3]
\begin{enumerate}
    \item $n=8$. Rangs $x=1..8$. $y$ : 18, 19, 20,5, 22, 23,5, 24, 26, 27,5.
    \item $G$ : $\bar{x} = \frac{1+...+8}{8} = 4,5$. $\bar{y} = \frac{18+19+20,5+22+23,5+24+26+27,5}{8} = \frac{180,5}{8} = 22,5625$. $G(4,5 ; 22,5625)$.
    \item \textbf{Mayer :} Groupes de 4.
    Groupe 1 ($x=1..4$) : $\bar{x}_1 = 2,5$. $\bar{y}_1 = \frac{18+19+20,5+22}{4} = \frac{79,5}{4} = 19,875$. $G_1(2,5 ; 19,875)$.
    Groupe 2 ($x=5..8$) : $\bar{x}_2 = 6,5$. $\bar{y}_2 = \frac{23,5+24+26+27,5}{4} = \frac{101}{4} = 25,25$. $G_2(6,5 ; 25,25)$.
    $a = \frac{25,25 - 19,875}{6,5 - 2,5} = \frac{5,375}{4} = 1,34375$.
    $b = \bar{y} - a\bar{x} = 22,5625 - 1,34375 \times 4,5 = 22,5625 - 6,046875 = 16,515625$.
    Droite : $\boxed{y = 1,34375\,x + 16,515625}$ (Arrondi Bac : $y = 1,34x + 16,52$).
    \item \textbf{2024 :} Rang $x = 10$. $y = 1,34375 \times 10 + 16,515625 = 13,4375 + 16,515625 = 29,953125 \approx 30,0$ milliards FCFA.
    \textbf{Fiabilité :} Extrapolation de 2 ans hors intervalle $[1;8]$. Tendance linéaire forte ($r$ proche de 1 visuellement), prévision plausible à court terme mais incertaine (chocs économiques, COVID, etc. non modélisés).
\end{enumerate}
\end{corrige}

\begin{savaistu}[Le saviez-vous ? Statistiques et Développement au Cameroun]
L'\cle{INS} (Institut National de la Statistique) réalise l'\cle{ECAM} (Enquête Camerounaise Auprès des Ménages) tous les 5-7 ans. Les indicateurs calculés (médiane des dépenses, quartiles de revenus, taux de pauvreté) guident la \cle{Stratégie Nationale de Développement (SND30)}. La méthode de Mayer, bien que simple, est utilisée par les services de planification (MINEPAT) pour des prévisions rapides de tendances démographiques ou de production agricole quand les modèles économétriques complexes ne sont pas nécessaires. Maîtriser ces outils, c'est participer à la décision publique éclairée !
\end{savaistu}

\findecours

\end{document}
